\documentclass[11pt]{article}
\usepackage[a4paper,margin=1in]{geometry}
\usepackage[T1]{fontenc}
\usepackage[utf8]{inputenc}
\usepackage{lmodern}
\usepackage{microtype}
\usepackage{amsmath,amssymb,amsthm,mathtools}
\usepackage{enumitem}
\usepackage{booktabs,longtable,array}
\usepackage{xcolor}
\usepackage{hyperref}
\usepackage{mathrsfs}
\allowdisplaybreaks

\hypersetup{
  colorlinks=true,
  linkcolor=blue!50!black,
  citecolor=blue!50!black,
  urlcolor=blue!50!black,
  pdftitle={A Research Compendium on the Sunflower Problem},
  pdfauthor={Samuel Mausberg and GPT-5.4 Pro}
}

\newtheorem{theorem}{Theorem}[section]
\newtheorem{lemma}[theorem]{Lemma}
\newtheorem{proposition}[theorem]{Proposition}
\newtheorem{corollary}[theorem]{Corollary}
\theoremstyle{definition}
\newtheorem{definition}[theorem]{Definition}
\newtheorem{example}[theorem]{Example}
\newtheorem{counterexample}[theorem]{Counterexample}
\newtheorem{question}[theorem]{Question}
\newtheorem{conjecture}[theorem]{Conjecture}
\newtheorem{remark}[theorem]{Remark}

\newcommand{\F}{\mathcal{F}}
\newcommand{\G}{\mathcal{G}}
\newcommand{\Hh}{\mathcal{H}}
\newcommand{\Kk}{\mathcal{K}}
\newcommand{\PP}{\mathcal{P}}
\newcommand{\Qq}{\mathcal{Q}}
\newcommand{\Tr}{\operatorname{Tr}}
\newcommand{\lk}{\operatorname{lk}}
\newcommand{\supp}{\operatorname{supp}}
\newcommand{\er}{\operatorname{er}}
\newcommand{\rank}{\operatorname{rk}}
\newcommand{\TV}{\operatorname{TV}}
\newcommand{\diag}{\operatorname{diag}}
\newcommand{\eps}{\varepsilon}
\newcommand{\binn}[2]{\binom{#1}{#2}}
\newcommand{\N}{\mathbb{N}}
\newcommand{\Z}{\mathbb{Z}}
\newcommand{\R}{\mathbb{R}}
\newcommand{\E}{\mathbb{E}}
\newcommand{\1}{\mathbf{1}}
\newcommand{\status}[1]{\par\smallskip\noindent\textbf{Status.} #1\par\smallskip}
\newcommand{\leannote}[1]{\par\smallskip\noindent\textbf{Lean formalization.} #1\par\smallskip}

\title{A Research Compendium on the Sunflower Problem\\[4pt]
\large Exact reductions, proved subclass theorems, counterexamples to false routes, and a status-layered record of an AI-assisted search program}
\author{Samuel Mausberg \and GPT-5.4 Pro}
\date{April 20, 2026}

\begin{document}
\maketitle

\begin{abstract}
This document is a research compendium for the sunflower problem,
recording the full program developed in the accompanying chat history.
We separate the material into four layers:
(1) fully proved statements,
(2) exact conditional reductions,
(3) counterexamples that kill tempting but false routes,
and (4) exploratory constructions whose local pieces are rigorous but whose final global rank bounds still require repair.

The proved core includes the exact leaf-stripping / terminal-kernel reduction,
an exact one-round sunflower theorem,
exact trace recursion,
the exact heavy-set / prefix-product formalism,
rigorous bounded-degree and degree-two terminal-kernel theorems,
sharp results for exact block products and positive-density near-product classes,
sharp LP reductions on projected branches,
a padding-invariant weighted choice-block invariant,
exact off-diagonal symbol-lift recurrences in fixed-alphabet code models,
and exact finite-state / sliding-window profitable-prefix and entropy-gap theorems.
We also prove or record rigorous counterexamples showing that low-order LP inverse theorems, low-correlation rigidity, broad support-piece dichotomies, and bounded-local-coordinate extraction are all false in natural forms.

On the exploratory side, we record exact local tensors for several hard fixed-alphabet diagonal models, including $(4,4)$, $(5,5)$, $(6,6)$, $(8,8)$, and the power-of-two diagonal ansatz. Their local support and global diagonalization steps are rigorous, but several current slice-rank counts use an invalid lower-tail estimate and are therefore quarantined here rather than used elsewhere.

The paper is written so that every theorem, proposition, lemma, and corollary in Parts~I--III is intended to be rigorous as written. Exploratory claims are explicitly labeled as provisional. The final appendices contain a ledger of all major claims touched in the project, together with their current status.
\end{abstract}

\tableofcontents

\section{Introduction and current status}
For integers $n,k\ge 2$, let $f(n,k)$ denote the least integer such that every family $\F$ of $n$-uniform sets with $|\F|\ge f(n,k)$ contains a $k$-sunflower.
Write
\[
M(n,k):=f(n,k)-1.
\]
Thus $M(n,k)$ is the largest size of a $k$-sunflower-free $n$-uniform family.
The Erdos--Rado sunflower conjecture asks whether for every fixed $k$ there exists a constant $c_k>0$ such that
\[
M(n,k)<c_k^n
\qquad\text{for all }n.
\]

The present public benchmark remains the logarithmic-base regime inaugurated by Alweiss--Lovett--Wu--Zhang and sharpened by Bell--Chueluecha--Warnke; Bloom's problem page still lists the global exponential-base question as open, while a December 2025 preprint of Fukuyama claims a sub-logarithmic-base improvement. We do not attempt to verify the latter here. Instead, this paper records a systematic AI-assisted exploration of the problem, with a strict separation between proved results and exploratory claims. See \cite{ErdosRado1960,ALWZ2020,Rao2020,BCW2021,Bloom2026,Fukuyama2025,Rao2025}.

\subsection*{Status convention}
\begin{itemize}[leftmargin=2em]
\item Statements labeled \emph{Theorem}, \emph{Proposition}, \emph{Lemma}, or \emph{Corollary} are intended to be rigorous and self-contained in this document.
\item Statements labeled \emph{Conjecture}, \emph{Question}, or discussed inside a \textbf{Status} paragraph as provisional are not claimed proved.
\item If a previously explored argument had a real proof gap, we record both the valid local part and the point where the proof currently breaks.
\end{itemize}

\subsection*{Main rigorous outputs of the project}
The main proved outputs are:
\begin{enumerate}[label=(\alph*),leftmargin=2.2em]
\item the exact leaf-stripping / terminal-kernel reduction (Theorem~\ref{thm:exact-reduction}),
\item the exact one-round sunflower theorem for stripping fibers (Theorem~\ref{thm:one-round-sunflower}),
\item the exact trace recursion and its honest recursive closure on solved trace classes (Theorem~\ref{thm:trace-recursion} and Theorem~\ref{thm:unionclass-closure}),
\item the exact heavy-set / prefix-product formalism (Theorem~\ref{thm:heavy-set} and Theorem~\ref{thm:adaptive-kernel-prefix}),
\item sharp theorems for exact block products (Theorem~\ref{thm:block-products}) and positive-density support-box families (Theorem~\ref{thm:positive-density-box}),
\item a sharp padding-invariant weighted choice-block lower bound for pairwise correlation (Theorem~\ref{thm:weighted-choice-block}),
\item exact bounded-degree terminal-kernel theorems, including the degree-two simplex classification (Theorem~\ref{thm:bounded-degree} and Theorem~\ref{thm:degree-two}),
\item exact fixed-order and sharper randomized deletion lifts in fixed-alphabet transversal/code models (Theorem~\ref{thm:extra-symbol} and Theorem~\ref{thm:randomized-deletion}),
\item exact profitable-prefix and entropy-gap theorems in fixed-memory finite-state classes (Theorem~\ref{thm:finite-state-prefix}, Theorem~\ref{thm:finite-state-statecount}, and Theorem~\ref{thm:finite-state-fixed-order}),
\item rigorous counterexamples killing several tempting routes, including broad support-piece dichotomies, low-order LP inverse theorems, low-$\sigma$ rigidity, bounded-local-coordinate extraction, and the naive A/B bridge (Section~\ref{sec:counterexamples}).
\end{enumerate}

\subsection*{Most promising unresolved directions}
Two unresolved directions survive repeated attempts:
\begin{enumerate}[label=(\roman*),leftmargin=2.2em]
\item a genuine diagonal hard-model theorem in fixed-alphabet code classes, most naturally a sharp exponential-rate theorem for $F_{4,4}(m)$ (now known to exceed $3^m$ already at $m=3$) or an exact theorem such as $F_{5,5}(m)=4^m$, and
\item a universal routing theorem sending every large terminal kernel into a proved trace class or directly yielding a profitable prefix block.
\end{enumerate}

\section{Background and notation}\label{sec:notation}
A \emph{$k$-sunflower} is a family $A_1,\dots,A_k$ of distinct sets such that $A_i\cap A_j$ is the same set for all $i\neq j$.
That common intersection is the \emph{kernel} (or \emph{core}); outside the kernel the sets are pairwise disjoint.

If $\Hh$ is a family of sets and $S$ is a set, define the link
\[
\Hh_S:=\{A\setminus S: A\in\Hh,\ S\subseteq A\}.
\]
If $\Hh$ is $r$-uniform, then $\Hh_S$ is $(r-|S|)$-uniform.

If $G$ is an $m$-uniform family and $P\subseteq \supp(G)$ is such that $|A\cap P|=r$ for every $A\in G$, we call $P$ a \emph{support piece of rank $r$}. Its \emph{trace family} is
\[
\Tr_P(G):=\{A\cap P: A\in G\}\subseteq \binom{P}{r},
\]
and for each trace $\tau\in \Tr_P(G)$ the associated \emph{projected branch} is
\[
G_\tau:=\{A\setminus P: A\in G,\ A\cap P=\tau\}.
\]
Then one always has the exact trace identity
\[
|G|=\sum_{\tau\in \Tr_P(G)} |G_\tau|.
\]

For an $n$-uniform family $F$ and a set $B$, define the degree
\[
d_F(B):=|\{A\in F:B\subseteq A\}|.
\]
Define
\[
\Gamma_t(F):=\sum_{|B|=t} d_F(B)^2,
\qquad
\Lambda_t(F):=\sum_{|B|=t} d_F(B)=\binom{n}{t}|F|.
\]
For $0\le t<n$, define
\[
\beta_t(F):=\frac{(t+1)\Gamma_{t+1}(F)}{(n-t)\Gamma_t(F)}.
\]
The products of the $\beta_t$'s are the central recursive invariants in the project.

For the fixed-alphabet transversal code model, we write
\[
F_{q,k}(m):=\max\{|C|: C\subseteq [q]^m,\ \text{no }k\text{ distinct codewords are coordinatewise all-equal-or-all-distinct}\}.
\]
When $q<k$, every subset of $[q]^m$ is automatically $k$-sunflower-free in this model.

\section{Exact leaf stripping and terminal kernels}\label{sec:leaf-stripping}
The first durable theorem of the project is that leaf stripping gives an exact reduction from arbitrary sunflower-free families to leafless terminal kernels.

\begin{definition}[Leaf stripping]
Let $\mathcal A$ be a finite family of sets and let
\[
L(\mathcal A):=\{x: d_{\mathcal A}(x)=1\}
\]
be the set of degree-one vertices.
Define
\[
\partial \mathcal A:=\{A\setminus L(\mathcal A): A\in \mathcal A\},
\]
with duplicates collapsed. Iterating,
\[
\mathcal F_0:=\mathcal F,
\qquad
\mathcal F_{t+1}:=\partial \mathcal F_t.
\]
The process stabilizes after finitely many steps at a family $K(\mathcal F)$ having no degree-one vertices. We call $K(\mathcal F)$ the \emph{terminal kernel} of $\mathcal F$.
\end{definition}

\begin{lemma}[Stripping preserves all nontrivial intersections]
If $\mathcal A\subseteq \mathcal F$ and $|\mathcal A|\ge 2$, then
\[
\bigcap_{E\in \mathcal A} E = \bigcap_{E\in \mathcal A} \bigl(E\cap L(\mathcal F)^c\bigr).
\]
In particular, for any two edges $E,F\in\mathcal F$ one has
\[
E\cap F=(E\setminus L(\mathcal F))\cap(F\setminus L(\mathcal F)).
\]
\end{lemma}

\begin{proof}
If $x\in \bigcap_{E\in\mathcal A}E$, then $x$ lies in at least two members of $\mathcal F$, so $x\notin L(\mathcal F)$, hence $x$ survives stripping. The reverse inclusion is immediate.
\end{proof}

\leannote{This lemma is formalized and proved in \texttt{FormalConjectures/Problems/Erdos/E20/Kernels/LeafStripping.lean}, principally as \texttt{familyInter\_stripped\_eq} and \texttt{pairwise\_intersection\_stripped\_eq}.}

\begin{theorem}[Exact one-round sunflower theorem]\label{thm:one-round-sunflower}
Let $\pi:\mathcal F\to \partial\mathcal F$ be the reduction map $E\mapsto E\setminus L(\mathcal F)$. For a set $K\in \partial\mathcal F$, define
\[
m_K:=|\pi^{-1}(K)|,
\]
and let $s_K$ be the maximum size of a family $\mathcal S\subseteq \partial\mathcal F\setminus\{K\}$ such that every two distinct members of $\mathcal S$ intersect exactly in $K$. Then the maximum size of a sunflower in $\mathcal F$ with kernel $K$ is exactly
\[
m_K+s_K.
\]
In particular, if $\mathcal F$ is $k$-sunflower-free, then every fiber of $\pi$ has size at most $k-1$, and hence
\[
|\mathcal F|\le (k-1)|\partial \mathcal F|.
\]
\end{theorem}

\begin{proof}
If $E_1,\dots,E_r$ all satisfy $\pi(E_i)=K$, then by the previous lemma each pair intersects exactly in $K$; hence any $r$ such sets already form an $r$-sunflower with kernel $K$.
Now if $\mathcal U\subseteq \mathcal F$ is any sunflower with kernel $K$, then for distinct $E,F\in\mathcal U$,
\[
\pi(E)\cap \pi(F)=E\cap F=K.
\]
If $\pi(E)=\pi(F)$, this common value must equal $K$, so only the kernel image can repeat. The repeated kernel images contribute at most $m_K$, while the distinct non-kernel images contribute at most $s_K$. Conversely, taking all $m_K$ preimages of $K$ and one preimage of each member of an optimal $\mathcal S$ yields a sunflower of size $m_K+s_K$.
\end{proof}

\leannote{This theorem is formalized and proved in \texttt{FormalConjectures/Problems/Erdos/E20/Kernels/LeafStripping.lean} as \texttt{exact\_one\_round\_sunflower\_theorem}; the wrapper \texttt{pasted\_exact\_one\_round\_sunflower\_theorem} appears in \texttt{FormalConjectures/Problems/Erdos/E20/InformalMap.lean}.}

\begin{lemma}[Strict parent lemma]
Suppose $t\ge 1$ and $B\in \mathcal F_t$ shrinks in the next round, i.e.
\[
\partial B\subsetneq B.
\]
Then there exists $A\in \mathcal F_{t-1}$ with
\[
\partial A=B
\qquad\text{and}\qquad
A\supsetneq B.
\]
\end{lemma}

\begin{proof}
Choose $x\in B\setminus \partial B$. Then $x$ has degree one in $\mathcal F_t$ but survived the previous round, so it had degree at least two in $\mathcal F_{t-1}$. All parents of $x$ in $\mathcal F_{t-1}$ map to the unique child $B$ containing $x$ in $\mathcal F_t$, and at least two such parents exist. One of them is strictly larger than $B$.
\end{proof}

\begin{corollary}[Length of the stripping process]
If $\mathcal F$ is $n$-uniform, then the stripping process stabilizes after at most $n$ rounds.
\end{corollary}

\begin{proof}
Choose a set that shrinks in the last nontrivial round and pull it back using the strict parent lemma. One obtains a strictly descending chain of sets of length equal to the number of nontrivial stripping rounds. Since the initial set has size $n$, the chain has length at most $n$.
\end{proof}

\begin{definition}[Terminal-kernel extremal function]
Let $M_{\ker}(n,k)$ be the maximum size of a $k$-sunflower-free family $\Hh$ of rank at most $n$ with no degree-one vertices.
\end{definition}

\begin{theorem}[Exact reduction to terminal kernels]\label{thm:exact-reduction}
For every $n,k\ge 2$,
\[
M_{\ker}(n,k)\le M(n,k)\le (k-1)^n M_{\ker}(n,k).
\]
Equivalently, if $f_{\ker}(n,k):=M_{\ker}(n,k)+1$, then
\[
f_{\ker}(n,k)\le f(n,k)\le (k-1)^n f_{\ker}(n,k).
\]
\end{theorem}

\begin{proof}
The upper bound is obtained by applying the one-round inequality $|\mathcal F_t|\le (k-1)|\mathcal F_{t+1}|$ at each stripping stage and using that there are at most $n$ nontrivial rounds.
For the lower bound, take any terminal $k$-sunflower-free family $\Hh$ of rank at most $n$, and pad each edge $H\in \Hh$ with a fresh private set $P_H$ of size $n-|H|$. The padded family is $n$-uniform, sunflower-free, and strips in one round back to $\Hh$.
\end{proof}

\leannote{The conditional terminal-kernel reduction algebra used by this theorem is formalized and proved in \texttt{FormalConjectures/Problems/Erdos/E20/Kernels/TerminalKernelReduction.lean}, including \texttt{terminalKernelReduction\_bound}, \texttt{terminalKernelReduction\_bound\_of\_class}, and the block-product specializations. The Lean file packages the iterated reduction factor as a hypothesis, so it does not by itself prove the displayed two-sided exact reduction theorem verbatim.}

\begin{example}[Sharpness of the factor $(k-1)^n$]
Let the ground set be the set of all nonempty words over $[k-1]$ of length at most $n$. For each $w=w_1\cdots w_n\in [k-1]^n$, define
\[
E_w:=\{w_1,w_1w_2,\dots,w_1\cdots w_n\}.
\]
Then $|E_w|=n$, the family has $(k-1)^n$ members, is $k$-sunflower-free, and each stripping round divides the size by exactly $k-1$. Thus the factor $(k-1)^n$ is best possible for pure leaf stripping.
\end{example}

\begin{remark}[Where leaf stripping stops]
Leaf stripping perfectly removes private padding, but it can stop immediately on a large leafless core. The exact block-product family with $n$ blocks of width $k-1$ is already terminal and sunflower-free, and therefore survives leaf stripping unchanged. Thus leaf stripping isolates the real obstruction perfectly, but does not resolve it.
\end{remark}

\section{Support pieces, traces, and exact trace recursion}\label{sec:support-pieces}
The second robust layer of the project is the support-piece / trace formalism.

\begin{definition}[Support piece and trace family]
Let $G$ be an $m$-uniform family and let $P\subseteq \supp(G)$ be such that $|A\cap P|=r$ for every $A\in G$. Then $P$ is a support piece of rank $r$. Its trace family and projected branches are
\[
\Tr_P(G):=\{A\cap P:A\in G\},
\qquad
G_\tau:=\{A\setminus P:A\in G,\ A\cap P=\tau\}.
\]
\end{definition}

\begin{theorem}[Exact trace recursion]\label{thm:trace-recursion}
For every support piece $P$ of rank $r$ in an $m$-uniform family $G$,
\[
|G|=\sum_{\tau\in \Tr_P(G)} |G_\tau|.
\]
Each projected branch $G_\tau$ is $(m-r)$-uniform. If $G$ is $k$-sunflower-free, then each $G_\tau$ is also $k$-sunflower-free. Consequently, if $\Tr_P(G)$ is itself $k$-sunflower-free, then
\[
|G|\le |\Tr_P(G)|\,f(m-r,k)\le f(r,k)f(m-r,k).
\]
\end{theorem}

\begin{proof}
The exact decomposition follows by partitioning $G$ by its traces on $P$. Uniformity is immediate. If some $G_\tau$ contained a $k$-sunflower, adjoining $\tau$ to each petal would give a $k$-sunflower in $G$. The final inequality follows because $\Tr_P(G)$ has rank $r$ and is assumed $k$-sunflower-free.
\end{proof}

\leannote{The decomposition, projected-branch uniformity, and projected-branch sunflower-freeness components are formalized and proved in \texttt{FormalConjectures/Problems/Erdos/E20/Compendium/Recurrences/TraceRecursion.lean} as \texttt{exact\_trace\_identity}, \texttt{projected\_branch\_uniform}, and \texttt{projected\_branch\_sunflower\_free}. The combined extremal inequality \texttt{trace\_recursion\_bound} in that file is still marked \texttt{sorry}, so this note does not certify the final displayed bound as a completed Lean theorem.}

The support-piece formalism is powerful but dangerous. The following universality lemma shows why broad ``balanced sparse versus heavy trace'' support-piece dichotomies are not genuine reductions.

\begin{proposition}[One-block universality of empty-core support pieces]\label{prop:one-block-universality}
Let $H$ be any coreless $r$-uniform $k$-sunflower-free family, and let $n\ge r$. Choose an $n$-set $A$ and, for each $B\in H$, choose a fresh private set $P_B$ of size $n-r$ disjoint from $A$ and from all other $P_{B'}$. Define
\[
E_B:=B\cup P_B,
\qquad
F:=\{A\}\cup\{E_B:B\in H\}.
\]
Then:
\begin{enumerate}[label=(\roman*),leftmargin=2.2em]
\item $F$ is $n$-uniform and $k$-sunflower-free,
\item $\{A\}$ is a maximal matching in $F$,
\item relative to this matching, the unique nontrivial support piece is the family $\{E_B:B\in H\}$,
\item its exact support traces on $A$ are exactly the members of $H$, and
\item the support piece has empty core because $H$ is coreless.
\end{enumerate}
In particular, empty-core support pieces can encode arbitrary smaller coreless sunflower-free families.
\end{proposition}

\begin{proof}
Uniformity is clear. Any sunflower entirely among the sets $E_B$ would project to a sunflower in $H$ because the private paddings are pairwise disjoint and cannot create extra intersections. A sunflower involving $A$ would force all non-core parts of the other members to be empty, which is impossible because the traces are distinct. Every $E_B$ meets $A$ in $B\neq\varnothing$, so $\{A\}$ is maximal. The remaining statements are immediate from construction.
\end{proof}

\leannote{This proposition is formalized and proved in \texttt{FormalConjectures/Problems/Erdos/E20/Compendium/Recurrences/TraceRecursion.lean} as \texttt{support\_piece\_universality}. The formal theorem makes explicit the ambient-renaming hypothesis \(\mathrm{Fintype.card}(\alpha)\le n\), together with \(1\le r\) and \(3\le k\), and proves the constructed \(n\)-uniform sunflower-free family, maximal one-edge matching condition, exact trace recovery, and empty trace core.}

\status{Proposition~\ref{prop:one-block-universality} is one of the most important negative results in the project. It shows that the empty-core support-piece problem still contains the full original sunflower problem in disguise. Any truly useful support-piece theorem must therefore use more than the bare existence of an empty-core piece.}

\section{The exact $\Gamma_t$/$\beta_t$ formalism and the kernel-prefix criterion}\label{sec:beta-gamma}
This section records the exact heavy-set / prefix-product machinery that survived all later revisions.

\begin{lemma}[Moment propagation]
For an $n$-uniform family $F$,
\[
\Lambda_t(F)=\binom{n}{t}|F|,
\qquad
A_t(F):=\frac{\Gamma_t(F)}{\Lambda_t(F)},
\qquad
A_{t+1}(F)=\beta_t(F)A_t(F).
\]
Hence
\[
A_t(F)=|F|\prod_{i=0}^{t-1}\beta_i(F).
\]
\end{lemma}

\begin{proof}
The identity for $\Lambda_t$ is a direct double count. The recursion for $A_t$ follows from the definitions of $\beta_t$ and $\Lambda_t$:
\[
\Lambda_{t+1}(F)=\frac{n-t}{t+1}\Lambda_t(F),
\qquad
\Gamma_{t+1}(F)=\frac{n-t}{t+1}\beta_t(F)\Gamma_t(F).
\]
Taking the ratio yields the claim.
\end{proof}

\leannote{These moment identities are formalized and proved in \texttt{FormalConjectures/Problems/Erdos/E20/Foundations/BetaChain.lean}, including \texttt{lambdaLevel\_eq\_choose\_mul\_card}, \texttt{averageDegreeLevel\_succ}, and \texttt{averageDegreeLevel\_eq\_card\_mul\_betaPrefix}.}

\begin{theorem}[Exact heavy-set lemma]\label{thm:heavy-set}
For every $0\le t\le n$,
\[
M_t(F):=\max_{|S|=t} d_F(S)
\ge
|F|\prod_{i=0}^{t-1}\beta_i(F).
\]
Equivalently,
\[
M_t(F)\ge A_t(F)=\frac{\Gamma_t(F)}{\Lambda_t(F)}.
\]
\end{theorem}

\begin{proof}
Since $\Gamma_t(F)=\sum_{|B|=t} d(B)^2\le M_t(F)\sum_{|B|=t} d(B)=M_t(F)\Lambda_t(F)$, we get
\[
M_t(F)\ge \frac{\Gamma_t(F)}{\Lambda_t(F)}=A_t(F).
\]
Now substitute the product expression for $A_t(F)$.
\end{proof}

\leannote{This theorem is formalized and proved in \texttt{FormalConjectures/Problems/Erdos/E20/Foundations/BetaChain.lean} as \texttt{heavy\_set\_average\_bound} and \texttt{heavy\_set\_betaPrefix\_bound}; the corresponding wrapper is \texttt{pasted\_exact\_heavy\_set\_lemma} in \texttt{FormalConjectures/Problems/Erdos/E20/InformalMap.lean}.}

\begin{corollary}[Block-extension form]
Let $0\le s<t\le n$. For any $s$-set $B$ with $d_F(B)>0$, the link $F_B$ is $(n-s)$-uniform of size $d_F(B)$. Hence there exists $T\supseteq B$ with $|T|=t$ such that
\[
d_F(T)\ge d_F(B)\prod_{j=0}^{t-s-1} \beta_j(F_B).
\]
\end{corollary}

\begin{proof}
Apply Theorem~\ref{thm:heavy-set} to the link $F_B$.
\end{proof}

\leannote{The link form is formalized and proved in \texttt{FormalConjectures/Problems/Erdos/E20/Foundations/BetaChain.lean} as \texttt{heavy\_link\_betaPrefix\_bound} and \texttt{exists\_heavy\_set\_in\_link}; \texttt{InformalMap.lean} records the corresponding wrapper \texttt{pasted\_heavy\_link\_lemma}.}

\begin{corollary}[Exact recursive upper bound]\label{cor:exact-recursion}
Let $g(n,k):=f(n,k)-1=M(n,k)$. Define
\[
B_{n,k}(t):=\inf\Bigl\{\prod_{i=0}^{t-1}\beta_i(F): F\text{ is $n$-uniform and $k$-sunflower-free}\Bigr\}.
\]
Then for every $1\le t\le n$,
\[
g(n,k)\le \frac{g(n-t,k)}{B_{n,k}(t)}.
\]
Equivalently,
\[
f(n,k)\le \left\lceil \frac{f(n-t,k)}{B_{n,k}(t)}\right\rceil.
\]
\end{corollary}

\begin{proof}
If $F$ is $k$-sunflower-free, then every $t$-link $F_S$ has size at most $g(n-t,k)$, so $M_t(F)\le g(n-t,k)$. Combine this with Theorem~\ref{thm:heavy-set}.
\end{proof}

\leannote{The corresponding beta-prefix quotient recursion is formalized and proved in \texttt{FormalConjectures/Problems/Erdos/E20/Foundations/BetaChain.lean} as \texttt{sunflower\_size\_le\_div\_of\_betaPrefix\_lower}, with wrapper \texttt{pasted\_beta\_chain\_recursive\_bound} in \texttt{FormalConjectures/Problems/Erdos/E20/InformalMap.lean}.}

\begin{remark}[What the full product says]
At the top level $t=n$, one has $\Gamma_n(F)=|F|$ and therefore
\[
|F|\prod_{i=0}^{n-1}\beta_i(F)=1.
\]
Thus the full $\beta$-product is exactly $1/|F|$. The game is to force sufficiently large \emph{initial} products before reaching level $n$.
\end{remark}

\subsection*{The weak chain from a cover}
If $H$ is $r$-uniform and $k$-sunflower-free, then a maximal matching has size at most $k-1$, so its union is a vertex cover of size at most $r(k-1)$. Assign each edge to one cover vertex. This gives
\[
\sum_x d_H(x)^2\ge \frac{|H|^2}{r(k-1)}.
\]
Applying this inside every link $F_B$ yields
\[
(t+1)\Gamma_{t+1}(F)\ge \frac{1}{(k-1)(n-t)}\Gamma_t(F),
\]
or equivalently
\[
\Gamma_t(F)\le (k-1)(t+1)(n-t)\Gamma_{t+1}(F).
\]
This explains exactly why the raw chain gives factorial-type losses: it is short by a factor of order $r^2$ from the constant-scale behavior suggested by genuine product examples.

\begin{theorem}[Adaptive kernel-prefix criterion]\label{thm:adaptive-kernel-prefix}
Let $\mathfrak K_k$ be a class of reduced kernel links equipped with an essential-rank function $\er$ satisfying the natural properties:
\begin{enumerate}[label=(\roman*),leftmargin=2.2em]
\item $\Hh_T\in \mathfrak K_k$ whenever $\Hh\in\mathfrak K_k$ and $\Hh_T\neq\varnothing$,
\item links remain $k$-sunflower-free,
\item $\er(\Hh_T)\le \er(\Hh)-|T|$,
\item $\er(\Hh)=0$ implies $\Hh=\{\varnothing\}$.
\end{enumerate}
Suppose there exists $\lambda_k>1$ such that every nontrivial reduced kernel link $\Hh\in\mathfrak K_k$ has some $t\in\{1,\dots,\er(\Hh)\}$ with
\[
\prod_{i=0}^{t-1}\beta_i(\Hh)\ge \lambda_k^{-t}.
\]
Then every $k$-sunflower-free $n$-uniform family satisfies
\[
g(n,k)\le \lambda_k^n,
\qquad
f(n,k)\le \lambda_k^n+1< (\lambda_k+1)^n.
\]
\end{theorem}

\begin{proof}
Apply the good prefix block on the current reduced kernel link, use Theorem~\ref{thm:heavy-set} to select a heavy $t$-set, pass to the corresponding link, and iterate. The sum of the chosen block lengths is at most the initial essential rank, so the total multiplicative loss is at most $\lambda_k^n$.
\end{proof}

\leannote{The abstract recurrence step behind this criterion is formalized and proved in \texttt{FormalConjectures/Problems/Erdos/E20/Recurrences/RecurrenceReduction.lean} as \texttt{exponential\_bound\_of\_hasProfitableDropRecurrence}; \texttt{InformalMap.lean} records it as \texttt{pasted\_profitable\_drop\_implies\_exponential\_bound}. The full reduced-kernel class hypothesis of Theorem~\ref{thm:adaptive-kernel-prefix} is not encoded as a single Lean theorem.}

\status{Theorem~\ref{thm:adaptive-kernel-prefix} is exact and rigorous, but it is only a reduction theorem: the missing input is a proof that every nontrivial reduced kernel link has \emph{some} profitable prefix block. The whole unrestricted sunflower conjecture is reduced to that one problem.}

\section{Correlation, weighted coordinates, and why raw rank is the wrong invariant}\label{sec:correlation}
A recurring theme in the project was that raw uniformity or even leaf-stripped uniformity is not the right quantity controlling pairwise correlation.

\begin{definition}[Normalized correlation]
For an $r$-uniform family $H$, define
\[
\sigma(H):=\frac{1}{r|H|^2}\sum_x d_H(x)^2.
\]
Equivalently,
\[
\sigma(H)=\frac{1}{r|H|^2}\sum_{A,B\in H}|A\cap B|.
\]
\end{definition}

The correct padding-invariant replacement for raw rank is a weighted packing of genuine one-out-of-$q$ coordinates.

\begin{definition}[Choice blocks and the weighted invariant]
Let $H^{\circ}$ be the edge-labeled $2$-core obtained from $H$ by repeatedly deleting degree-one vertices without merging coincident edges. A subset $B\subseteq V(H^{\circ})$ is a \emph{choice block} if every edge of $H^{\circ}$ meets $B$ in exactly one point. Define
\[
\sigma_{\mathrm{cb}}(H):=\max\Bigl\{\sum_B \frac{\lambda_B}{|B|}:
\lambda_B\ge 0,\ B\text{ a choice block},\ \sum_{B\ni x}\lambda_B\le 1\ \forall x\Bigr\},
\]
and normalize
\[
\rho_k(H):=(k-1)\sigma_{\mathrm{cb}}(H).
\]
\end{definition}

\begin{theorem}[Weighted choice-block lower bound]\label{thm:weighted-choice-block}
For every finite family $H$,
\[
\sum_x d_H(x)^2\ge \sigma_{\mathrm{cb}}(H)|H|^2
=\frac{1}{k-1}\rho_k(H)|H|^2.
\]
In particular, a one-out-of-$q$ coordinate contributes exactly $(k-1)/q$ units to $\rho_k$.
\end{theorem}

\begin{proof}
Fix any feasible weighted packing $(\lambda_B)$. Since every edge meets a choice block $B$ in exactly one point,
\[
\sum_{x\in B} d_H(x)=|H|.
\]
By Cauchy,
\[
\sum_{x\in B} d_H(x)^2\ge \frac{|H|^2}{|B|}.
\]
Multiply by $\lambda_B$ and sum over $B$. The vertex constraints $\sum_{B\ni x}\lambda_B\le 1$ then give the claim.
\end{proof}

\begin{counterexample}[Why raw core size is false]
Fix $s>r\ge 1$ and $t\ge 0$. Let there be $s$ safe blocks of size $k-1$ and $t$ noisy blocks of size $(k-1)^r$. Index edges by $u\in [k-1]^s$, and define
\[
E_u:=\{(i,u_i):1\le i\le s\}\cup \{(s+j,(u_1,\dots,u_r)):1\le j\le t\}.
\]
Call the resulting family $H_{s,t,r}$. Then no leaf stripping occurs, the family is $k$-sunflower-free, and
\[
\sigma(H_{s,t,r})=\frac{s}{k-1}+\frac{t}{(k-1)^r}.
\]
Thus noisy coordinates can contribute almost nothing to pairwise correlation while contributing one full unit to raw leaf-stripped uniformity. Raw uniformity is therefore the wrong invariant.
\end{counterexample}

\begin{counterexample}[Constant-scale correlation fails even on the $2$-core]\label{ce:simplex}
Let $H_r=\{E_1,\dots,E_{r+1}\}$ be the $r$-uniform simplex family on vertices $v_{ij}$ ($1\le i<j\le r+1$), where
\[
E_i:=\{v_{ij}:j\neq i\}.
\]
Every vertex has degree $2$, every two edges intersect in exactly one vertex, and $H_r$ is $k$-sunflower-free for every $k\ge 3$.
Moreover,
\[
\frac{\sum_x d(x)^2}{r|H_r|^2}=\frac{2}{r+1}\to 0.
\]
So no bound of the form
\[
\sum_x d(x)^2\ge c_k r|H|^2
\]
can hold even on leafless kernels.
\end{counterexample}

\leannote{This simplex obstruction is formalized and proved in \texttt{FormalConjectures/Problems/Erdos/E20/Families/Counterexample.lean}; \texttt{InformalMap.lean} restates the relevant facts as \texttt{simplexFamily\_uniform}, \texttt{simplexFamily\_sunflowerFree}, \texttt{simplexFamily\_degreeSquareSum}, \texttt{simplexFamily\_normalized\_ratio}, and \texttt{pasted\_no\_positive\_two\_core\_branching\_constant}.}

\begin{proposition}[Sharp universal replacement for the simplex obstruction]
If $H$ is an $r$-uniform $k$-sunflower-free family with $m:=|H|$, then
\[
\sum_x d_H(x)^2\ge \frac{m^2}{k-1}+(r-1)m.
\]
This is sharp, with equality attained by the disjoint union of $k-1$ simplex families.
\end{proposition}

\begin{proof}
Let $I(H)$ be the intersection graph on the edge set of $H$. Since $H$ has no $k$ pairwise disjoint edges, $\alpha(I(H))\le k-1$. By Turan,
\[
e(I(H))\ge \frac{m(m-k+1)}{2(k-1)}\ge \frac{m^2}{2(k-1)}-\frac{m}{2}.
\]
Every intersecting pair contributes at least one to $\sum_x \binom{d(x)}{2}$, so
\[
\sum_x d(x)^2 = 2\sum_x \binom{d(x)}{2}+rm\ge 2e(I(H))+rm\ge \frac{m^2}{k-1}+(r-1)m.
\]
\end{proof}

\subsection*{Bounded-degree kernels are polynomially harmless}
The preceding counterexample shows that constant-scale correlation is false even on the $2$-core. Nevertheless, bounded-degree terminal kernels are harmless for a different reason.

\begin{theorem}[Bounded-degree kernels]\label{thm:bounded-degree}
Let $J$ be an $s$-uniform hypergraph with matching number $\nu(J)$. If every edge of $J$ meets at most $D$ other edges, then
\[
|J|\le (D+1)\nu(J).
\]
In particular, if $J$ is a kernel of a $k$-sunflower-free family, then $\nu(J)\le k-1$, and so
\[
|J|\le (k-1)(D+1).
\]
If $\Delta(J)\le \Delta$, then each edge meets at most $s(\Delta-1)$ other edges, hence
\[
|J|\le (k-1)\bigl((\Delta-1)s+1\bigr).
\]
\end{theorem}

\begin{proof}
Consider the intersection graph $L(J)$ on the edge set of $J$. Its independence number equals $\nu(J)$. Turan's theorem implies $|J|\le (D+1)\nu(J)$ when the graph has maximum degree at most $D$.
\end{proof}

\leannote{The finite maximal-matching form of this bounded-overlap and bounded-degree argument is formalized and proved in \texttt{FormalConjectures/Problems/Erdos/E20/Kernels/KernelBounds.lean}, notably as \texttt{card\_le\_overlap\_bound\_of\_maximalMatching}, \texttt{edgeOverlapSet\_card\_le\_uniform\_degree}, \texttt{card\_le\_vertexDegree\_bound\_of\_maximalMatching}, and \texttt{card\_le\_vertexDegree\_bound\_of\_small\_maximalMatching}. The wrappers \texttt{pasted\_bounded\_overlap\_kernel\_bound}, \texttt{pasted\_bounded\_degree\_kernel\_bound}, and \texttt{pasted\_small\_matching\_bounded\_degree\_kernel\_bound} are in \texttt{InformalMap.lean}.}

\begin{theorem}[Degree-two classification]\label{thm:degree-two}
Let $J$ be an $s$-uniform hypergraph with $\Delta(J)\le 2$. Then
\[
|J|\le (s+1)\nu(J).
\]
If $s\ge 2$, equality holds if and only if $J$ is a disjoint union of $\nu(J)$ copies of the $s$-simplex family.
In particular, if $J$ is a kernel of a $k$-sunflower-free family, then
\[
|J|\le (k-1)(s+1),
\]
with equality exactly for a disjoint union of $k-1$ simplices.
\end{theorem}

\begin{proof}
When $\Delta(J)\le 2$, each edge meets at most $s$ other edges, so the general bounded-degree theorem gives the bound.
Equality forces the intersection graph to be a disjoint union of copies of $K_{s+1}$, and degree-two forces each component to realize the simplex incidence pattern.
\end{proof}

\section{LP bounded-block covers and projected-branch obstructions}\label{sec:lp}
The LP route cleanly reduces heavy seeds to a projected-branch obstruction problem.

\begin{definition}[Heavy seeds]
For an $n$-uniform family $F$ and a block $S$, define
\[
F(S):=\{A\in F:S\subseteq A\},
\qquad
\mu_F(S):=\frac{|F(S)|}{|F|}.
\]
Fix $\eps\in(0,1)$. We say that $S$ is \emph{$\eps$-heavy} if
\[
\mu_F(S)\ge \eps^{|S|}.
\]
\end{definition}

\begin{proposition}[Heavy seeds imply an exponential bound]
If there exists $\eps_k>0$ such that every nonempty $k$-sunflower-free $n$-uniform family contains a nonempty $\eps_k$-heavy block, then
\[
f(n,k)\le \eps_k^{-n}.
\]
\end{proposition}

\begin{proof}
If $S\neq \varnothing$ is $\eps_k$-heavy, then the projected family $F_S$ is $(n-|S|)$-uniform and $k$-sunflower-free, so
\[
|F|\le \eps_k^{-|S|} f(n-|S|,k).
\]
Iterate.
\end{proof}

\leannote{The corresponding natural-number recurrence form is formalized and proved in \texttt{FormalConjectures/Problems/Erdos/E20/Compendium/Recurrences/HeavySet.lean} as \texttt{heavy\_seeds\_imply\_exponential\_bound}. The nearby stronger theorem \texttt{recursive\_upper\_bound} in that same file remains marked \texttt{sorry} and is not being cited here.}

\begin{definition}[Fractional bounded-block cover]
For a family $H$ and parameters $B\in\N$, $\eps\in(0,1)$, define
\[
\tau_{B,\eps}(H):=\min \sum_{\substack{T\subseteq V(H)\\ |T|\le B}} x_T\,\eps^{|T|}
\]
subject to
\[
x_T\ge 0,
\qquad
\sum_{\substack{T\subseteq E\\ |T|\le B}} x_T\ge 1
\quad\text{for every }E\in H.
\]
\end{definition}

\begin{proposition}[Cheap covers force heavy blocks]
If $\tau_{B,\eps}(H)\le 1$, then $H$ contains some block $T$ with $|T|\le B$ such that
\[
\mu_H(T)\ge \eps^{|T|}.
\]
\end{proposition}

\begin{proof}
Average the covering inequality over edges of $H$.
\end{proof}

\leannote{The averaging/double-counting inequality used in this proposition is formalized and proved in \texttt{FormalConjectures/Problems/Erdos/E20/Compendium/Recurrences/LPCovers.lean} as \texttt{cheap\_cover\_total\_weight}. The existential heavy-block conclusion itself is not recorded there as a completed Lean theorem.}

\begin{proposition}[Dual obstruction]
The dual linear program is
\[
\tau_{B,\eps}(H)=\max \sum_{E\in H} y_E
\]
subject to
\[
y_E\ge 0,
\qquad
\sum_{E\supseteq T} y_E\le \eps^{|T|}
\quad\text{for all }|T|\le B.
\]
Thus if there is no heavy block, one obtains a dual probability measure on $H$ whose small-block marginals are all exponentially tiny.
\end{proposition}

\begin{remark}[Why projected branches matter]
Raw branches $H_u$ are trivial from the primal viewpoint because the single block $\{u\}$ covers everything. The informative object is the projected branch
\[
L_u:=\{E\setminus \{u\}: E\in H_u\}.
\]
If $\tau_{B,\eps}(L_u)\le 1$, then one obtains a heavy block in $L_u$ and hence a linear seed in the ambient family by reattaching $u$. The remaining difficulty is to eliminate the exact dual obstructions on $L_u$.
\end{remark}

\subsection*{Exact dual obstructions for the normalized block-box LP}
We now record the exact obstruction class for the normalized projected-branch LP.

\begin{definition}
Let $\Omega=\prod_{i=1}^m X_i$ with ambient product measure $\pi=\pi_1\otimes\cdots\otimes\pi_m$. For $t\ge 0$, let $\mathcal C_t$ be the family of subboxes depending on at most $t$ blocks. For $B\subseteq \Omega$, define the normalized cover value
\[
\tau_t(B):=\min\Bigl\{\sum_{C\in \mathcal C_t} \pi(C)\lambda_C: \lambda_C\ge 0,\ \sum_{C\ni x}\lambda_C\ge 1\ \forall x\in B\Bigr\}.
\]
\end{definition}

\begin{proposition}[Exact dual classification for the normalized LP]\label{prop:oa-dual}
One has $\tau_t(B)=1$ if and only if $B$ supports a probability measure $\mu$ such that for every subset $I\subseteq [m]$ with $|I|\le t$,
\[
\mu_I=\pi_I.
\]
Equivalently, the exact normalized dual obstructions are precisely relative orthogonal-array / $t$-wise independent measures with respect to the ambient product.
\end{proposition}

\begin{proof}
If $\tau_t(B)=1$, take an optimal dual solution and normalize it to a probability measure $\mu$. For every cylinder event depending on at most $t$ coordinates, dual feasibility and its complement force equality with the ambient product measure. Conversely, if such a measure exists, it is dual-feasible with objective value $1$.
\end{proof}

\begin{counterexample}[Parity/code-slice obstruction to product-like inverse theorems]\label{ce:parity-projective}
Let $q:=k-1$, let $G=\Z/q\Z$, and consider the transversal system
\[
L:=\{E_x:x\in G^m,\ x_1+\cdots+x_m=0\}
\]
on block system $[m]\times G$, with the uniform measure $p$ on $L$.
Then:
\begin{enumerate}[label=(\roman*),leftmargin=2.2em]
\item $L$ is $k$-sunflower-free,
\item for every admissible pattern $T$ with $|T|<m$,
\[
\Pr_{E\sim p}[T\subseteq E]=q^{-|T|},
\]
exactly as for the uniform block-product measure,
\item hence every bounded-order normalized-correlation statistic built from such inclusion events is exactly the product one,
\item nevertheless $p$ stays bounded away in total variation from every block-product measure.
\end{enumerate}
In particular, the projected-branch dichotomy ``close to product or else positive bounded-order correlation'' is false.
\end{counterexample}

\begin{proof}
The first point follows because in a block of size $q=k-1$ a $k$-sunflower cannot realize the all-distinct option. The exact low-order marginal calculation follows from the linear sum constraint. The total-variation lower bound follows by a character calculation on the subgroup $\{x\in G^m: x_1+\cdots+x_m=0\}$.
\end{proof}

\leannote{The Lean development proves the sunflower-free and uniform-prefix components of this parity-slice obstruction in \texttt{FormalConjectures/Problems/Erdos/E20/Families/TransversalCounterexample.lean} and \texttt{FormalConjectures/Problems/Erdos/E20/InformalMap.lean}, including \texttt{pasted\_parity\_slice\_prefix\_uniform} and \texttt{pasted\_transversal\_sunflower\_free\_below\_alphabet\_threshold}. The total-variation separation in item (iv) is not recorded as a completed Lean theorem.}

\begin{counterexample}[Parity slice with all proper marginals uniform]\label{ce:parity-proper-marginals}
Let $q\ge 2$ and
\[
P_n:=\{x\in (\Z/q\Z)^n: x_1+\cdots+x_n=0\}.
\]
Then $|P_n|=q^{n-1}$. For every proper subset $I\subsetneq [n]$, the projection of $P_n$ to the coordinates in $I$ is all of $(\Z/q\Z)^I$, and each fiber has size $q^{n-1-|I|}$. In particular every marginal on fewer than $n$ coordinates is exactly uniform. When $q=k-1$, the associated transversal family is still $k$-sunflower-free.
\end{counterexample}

\begin{proof}
Fix $I\subsetneq [n]$ and prescribe arbitrary values on the coordinates in $I$. Choose values freely on all but one of the coordinates in $[n]\setminus I$; there are $q^{n-|I|-1}$ such choices. The final remaining coordinate is then uniquely determined by the equation $x_1+\cdots+x_n=0$ in $\Z/q\Z$. Hence every point of $(\Z/q\Z)^I$ has exactly $q^{n-1-|I|}$ preimages, proving the uniform-marginal claim. The final statement is the usual width-$(k-1)$ transversal observation: in each coordinate there are only $k-1$ symbols, so a $k$-sunflower cannot realize the all-distinct option there.
\end{proof}

\leannote{The prefix-uniform specialization and the sunflower-free transversal conclusion are formalized and proved in \texttt{FormalConjectures/Problems/Erdos/E20/Families/TransversalCounterexample.lean} and exposed in \texttt{InformalMap.lean} as \texttt{pasted\_parity\_slice\_prefix\_uniform} and \texttt{pasted\_transversal\_sunflower\_free\_below\_alphabet\_threshold}. The arbitrary-subset marginal statement is not separately recorded as a completed Lean theorem.}

\begin{remark}[Decorated OA lifts]
The exact obstruction class is even larger than affine slices. If $L:[q]^m\to [q]^R$ is a quotient map whose fibers are exact $t$-wise independent objects, then any dense pullback $L^{-1}(A)$ with arbitrary $A\subseteq [q]^R$ still matches the ambient product measure on all block sets of size at most $t$. Thus growing-order block tests do not by themselves force bounded-codimension structure. Any positive inverse theorem in this direction must first quotient out the maximal OA factor.
\end{remark}

\section{Exact block products and near-products}\label{sec:block-products}
The genuine block-product regime turned out to be completely explicit.

\begin{definition}[Exact block product]
Let the ground set split into disjoint blocks $X_1\sqcup\cdots\sqcup X_r$. For each $i$, let $\mathcal A_i=\{A_{i,1},\dots,A_{i,w_i}\}$ be a family of pairwise disjoint nonempty subsets of $X_i$. The associated block product is
\[
\PP=\mathcal A_1\otimes\cdots\otimes \mathcal A_r:=\Bigl\{\bigcup_{i=1}^r A_{i,t_i}: t_i\in [w_i]\Bigr\}.
\]
The integer $w_i$ is the width of coordinate $i$.
\end{definition}

\begin{theorem}[Exact block-product dichotomy and stability]\label{thm:block-products}
Let $k\ge 3$ and let $\PP$ be an exact block product with widths $w_1,\dots,w_r$.
Then:
\begin{enumerate}[label=(\roman*),leftmargin=2.2em]
\item $\PP$ is $k$-sunflower-free if and only if $w_i\le k-1$ for every $i$.
\item If some coordinate has width at least $k$, then $\PP$ already contains a $k$-sunflower.
\item The number of unordered $k$-sunflowers in $\PP$ is
\[
N_k(\PP)=\frac{1}{k!}\left(\prod_{i=1}^r (w_i+(w_i)_k)-\prod_{i=1}^r w_i\right),
\]
where $(w)_k=w(w-1)\cdots(w-k+1)$.
\item Among all sunflower-free subproducts of a given ambient product with widths $w_i$, the largest has size exactly
\[
\prod_{i=1}^r \min\{w_i,k-1\}.
\]
\item If $\PP$ is sunflower-free and
\[
|\PP|\ge (k-1)^r e^{-\eps r},
\]
then all but $O_k(\eps r)$ coordinates have width exactly $k-1$.
\end{enumerate}
\end{theorem}

\begin{proof}
The sunflower condition factorizes coordinatewise because the local pieces in each coordinate are pairwise disjoint. In a sunflower, each coordinate is either in the core (all chosen local pieces equal) or in the petals (all chosen local pieces distinct). Thus width at least $k$ in one coordinate already creates a $k$-sunflower, while width at most $k-1$ in every coordinate forbids the all-distinct option everywhere and hence forces equality of all petals. The counting and stability formulas follow by the same factorization.
\end{proof}

\leannote{The block-product size, sunflower-free width-bound, maximal-subproduct, and leaf-stripping components are formalized and proved in \texttt{FormalConjectures/Problems/Erdos/E20/Families/ExplicitKernels.lean} and \texttt{FormalConjectures/Problems/Erdos/E20/Compendium/Models/BlockProducts.lean}, including \texttt{card\_blockProductFamily}, \texttt{blockProductFamily\_sunflowerFree}, \texttt{strippedSupport\_blockProductFamily\_eq}, \texttt{sunflower\_free\_block\_product\_bound}, and \texttt{max\_sunflower\_free\_subproduct}. The sunflower-counting and stability clauses are not cited here as completed Lean theorems.}

\status{Theorem~\ref{thm:block-products} shows that genuine block products are not the hard obstruction. The only true product obstruction is width $k-1$ in every coordinate. Wide coordinates create sunflowers immediately, and near-extremizers are close to the width-$(k-1)$ transversal template.}

\begin{theorem}[Positive-density support-box theorem]\label{thm:positive-density-box}
Fix $k\ge 2$ and $0<\delta\le 1$. Let $H$ be an $r$-uniform partite family using exactly one vertex from each of $r$ coordinate parts, and suppose
\[
|H|\ge \delta \prod_{i=1}^r |\pi_i(H)|.
\]
Set
\[
Q:=\left\lfloor \frac{k-1}{\delta}\right\rfloor.
\]
If $H$ is $k$-sunflower-free, then:
\begin{enumerate}[label=(\roman*),leftmargin=2.2em]
\item every active coordinate width satisfies $|\pi_i(H)|\le Q$,
\item every nonempty iterated link $J$ of $H$ satisfies
\[
\sigma(J)\ge \frac{1}{Q}\ge \frac{\delta}{k-1},
\]
\item consequently,
\[
|H|\le Q^r\le \left(\frac{k-1}{\delta}\right)^r.
\]
For $\delta=1$, this yields the sharp exact-product bound $|H|\le (k-1)^r$.
\end{enumerate}
\end{theorem}

\begin{proof}
If some fiber over a fixed choice of the other coordinates had size at least $k$, those $k$ edges would form a sunflower. Thus every fiber has size at most $k-1$. Compare this with the density lower bound to obtain $|\pi_i(H)|\le Q$. Then use Cauchy on the degree sequence in each coordinate to get $\sigma(H)\ge \frac{1}{r}\sum_i \frac1{|\pi_i(H)|}\ge 1/Q$. The recursion $|H|\le Q|H_x|$ on a heavy vertex then yields $|H|\le Q^r$.
\end{proof}

\leannote{The sharp \(\delta=1\) counting specialization is formalized and proved in \texttt{FormalConjectures/Problems/Erdos/E20/Compendium/Models/BlockProducts.lean} as \texttt{positive\_density\_support\_box\_delta\_one}. The full positive-density theorem for arbitrary \(0<\delta\le 1\), including the Cauchy lower bound on \(\sigma(J)\), is not recorded as a completed Lean theorem.}

\section{Fixed-memory finite-state branches}\label{sec:finite-state}
The finite-state lane is one of the most useful proved model classes. The right invariant is the transfer root, not the one-step local alphabet.

\begin{definition}[Finite-state branch]
Let $\mathcal A=(Q,E,s)$ be a right-resolving finite automaton emitting one symbol per edge. For a state $q\in Q$ and $m\ge 0$, let $\mathcal B_q(m)$ be the set of admissible words of length $m$ starting from $q$. The \emph{prefix complexity} is
\[
N_*(t):=\max_{q\in Q} |\mathcal B_q(t)|.
\]
If $T$ is the transfer matrix, its Perron root is denoted $\rho(T)$.
\end{definition}

\begin{theorem}[Exact profitable-prefix theorem]\label{thm:finite-state-prefix}
Let $F$ be a reduced kernel link of residual length $m$ lying inside a finite-state branch class, and let $1\le t\le m$. Then
\[
B^{\ker}_{m,k}(t)\ge N_*(t)^{-1}.
\]
Equivalently, every such reduced link has a contiguous prefix block of length $t$ whose best exact block profit is at least $N_*(t)^{-1}$.
If the transition graph is irreducible with Perron right eigenvector $r$ and condition number
\[
\kappa(r):=\frac{\max r}{\min r},
\]
then
\[
N_*(t)\le \kappa(r)\rho(T)^t,
\qquad
B^{\ker}_{m,k}(t)\ge \kappa(r)^{-1}\rho(T)^{-t}.
\]
Consequently,
\[
h_{\mathcal A}(m,k)\le \kappa(r)\rho(T)^m,
\]
where $h_{\mathcal A}(m,k)$ is the extremal size among $k$-sunflower-free reduced links in the automaton class.
\end{theorem}

\begin{proof}
A prefix cell of length $t$ is determined by an admissible word of length $t$, so there are at most $N_*(t)$ such cells. Hence the heaviest cell has density at least $1/N_*(t)$. The Perron-root estimate follows from $N_t=T^t\1$ and the standard super-eigenfunction comparison.
\end{proof}

\leannote{The finite-state profitable-chain and automaton heavy-prefix/recurrence forms are formalized and proved in \texttt{FormalConjectures/Problems/Erdos/E20/Recurrences/FiniteStatePrefixes.lean} and \texttt{FormalConjectures/Problems/Erdos/E20/Recurrences/AutomatonBranches.lean}, including \texttt{exists\_phi\_chain}, \texttt{exists\_heavy\_prefix}, \texttt{exists\_heavy\_prefix\_ratio}, and \texttt{hA\_recurrence}; \texttt{InformalMap.lean} exposes these as \texttt{pasted\_finite\_state\_profitable\_chain} and \texttt{pasted\_automaton\_branch\_recurrence}. The Perron-root comparison is represented through these formal recurrence/superharmonic statements rather than as this theorem verbatim.}

\begin{theorem}[State-count entropy gap]\label{thm:finite-state-statecount}
Let $X$ be an irreducible right-resolving $q$-ary finite-state branch with $m$ states. Then exactly one of the following holds:
\begin{enumerate}[label=(\roman*),leftmargin=2.2em]
\item $X$ is the full $q$-shift, in which case $|L_n(X)|=q^n$ for all $n$;
\item or else, writing $n=tm+s$ with $0\le s<m$,
\[
|L_n(X)|\le q^s (q^m-1)^t.
\]
In particular,
\[
\limsup_{n\to\infty} |L_n(X)|^{1/n}\le (q^m-1)^{1/m}<q.
\]
\end{enumerate}
\end{theorem}

\begin{proof}
If every state has all $q$ labels leaving it, then the language is the full shift. Otherwise some state misses a label. Strong connectivity implies that from any state, some forbidden length-$m$ word can be reached, so every state has at most $q^m-1$ admissible words of length $m$. Chop a length-$n$ word into blocks of size $m$.
\end{proof}

\leannote{The core block-count inequality behind this theorem is formalized and proved in \texttt{FormalConjectures/Problems/Erdos/E20/Compendium/Models/FiniteState.lean} as \texttt{entropy\_gap\_block\_bound}. The full irreducible right-resolving dichotomy is not recorded as a single completed Lean theorem.}

\begin{theorem}[Sharp fixed-memory entropy gap]\label{thm:finite-state-fixed-order}
Let $X$ be an irreducible $r$-step $q$-ary sliding-window branch. Then exactly one of the following holds:
\begin{enumerate}[label=(\roman*),leftmargin=2.2em]
\item $X$ is the full $q$-shift;
\item or else
\[
|L_n(X)|\le \frac{\lambda_{q,r}^{n+1}}{q-1}
\qquad (n\ge 0),
\]
where $\lambda_{q,r}\in(q-1,q)$ is the unique root of
\[
\lambda^{r+1}-q\lambda^r+q-1=0,
\]
equivalently
\[
1=(q-1)\sum_{j=1}^r \lambda^{-j}.
\]
This bound is sharp in exponential rate, with equality attained by the branch forbidding only the word $0^r$.
\end{enumerate}
In particular, there is no $k$-only entropy gap across the entire finite-state class when $q=k-1$.
\end{theorem}

\begin{proof}
Any proper $r$-step branch forbids some length-$r$ block. The standard prefix automaton for avoiding that block yields the comparison vector used in the project, and hence the upper bound by the Perron root of the associated companion matrix. The single-forbidden-block branch shows sharpness. Letting $r\to\infty$ shows that no $k$-only gap uniform in memory can hold.
\end{proof}

\leannote{The exact forbidden-block recursion and its closed-form block bound are formalized and proved in \texttt{FormalConjectures/Problems/Erdos/E20/Recurrences/FixedMemoryBound.lean} as \texttt{card\_language\_add\_le} and \texttt{card\_language\_bound}; the polynomial-equivalence and root-existence components are in \texttt{FormalConjectures/Problems/Erdos/E20/Compendium/Models/FiniteState.lean} as \texttt{sliding\_window\_polynomial\_equiv} and \texttt{sliding\_window\_root\_exists}. The sharper Perron-bound theorem as stated here is only partially represented by these Lean components.}

\status{Finite-state branches are therefore a genuine theorem-producing class, but only after some separate routing theorem forces bounded state count or bounded memory. By themselves they do not solve the unrestricted sunflower problem.}



\subsection*{Hereditary support-pattern selectors and the greedy-branch obstruction}
The finite-state theorems of this section apply to hereditary recursive states, not merely to the surface language of emitted symbols.

\begin{definition}[Hereditarily finite-state support-pattern selector]\label{def:hereditary-pattern-selector}
Fix a uniform bound $B$ on the size of the chosen transversals.
A \emph{residual support-pattern system} consists of residual states $v$, a chosen transversal $H(v)$ with $|H(v)|\le B$, a realized pattern set
\[
\Pi(v)\subseteq 2^{H(v)}\setminus \{\varnothing\},
\]
and, for each $\sigma\in \Pi(v)$, a child state $v\!\restriction \sigma$ obtained by conditioning on the exact support pattern $\sigma$ and then cone-reducing.

A selector $s$ on support patterns is called \emph{hereditarily finite-state} if there exist a finite set $Q$ and a map
\[
\tau:\{\text{residual states}\}\to Q
\]
such that, after relabeling $H(v)\cong [m]$ for the relevant value of $m$,
\begin{enumerate}[label=(\roman*),leftmargin=2.2em]
\item the realized pattern set $\Pi(v)\subseteq 2^{[m]}\setminus\{\varnothing\}$ depends only on $\tau(v)$;
\item for every $\sigma\in \Pi(v)$, the child state $\tau(v\!\restriction \sigma)$ depends only on $\tau(v)$ and $\sigma$;
\item the chosen pattern $s(v)$ depends only on $\tau(v)$.
\end{enumerate}
Equivalently, the relevant Myhill--Nerode congruence is the hereditary congruence on support-pattern contexts, not merely the word-language congruence of emitted symbols.
\end{definition}

\begin{proposition}[Finite-profile criterion for support-pattern selectors]\label{prop:selector-finite-profile}
Fix $B<\infty$.
A selector on exact support patterns gives a uniformly solved finite-state branch class if and only if it factors through a finite hereditary quotient as in Definition~\ref{def:hereditary-pattern-selector}.

More concretely, if $q(v)\in Q$ is such a quotient, then after relabeling $H(v)\cong [m]$ the associated child-profile
\[
\theta(v):=\bigl(m,\Pi(v),\chi_v\bigr),
\qquad
\chi_v:\Pi(v)\to Q,
\qquad
\chi_v(\sigma):=q(v\!\restriction \sigma),
\]
takes values in the finite set
\[
\Theta_B(Q):=
\Bigl\{
(m,\Pi,\chi):
1\le m\le B,\ 
\Pi\subseteq 2^{[m]}\setminus\{\varnothing\},\ 
\chi:\Pi\to Q
\Bigr\}.
\]
Conversely, any finite recursive presentation of the selector yields such a quotient by taking the recursive state itself.
\end{proposition}

\begin{proof}
Assume first that a finite hereditary quotient $q(v)\in Q$ exists.
For each state $v$, the profile $\theta(v)$ is determined by $q(v)$.
Because $1\le m\le B$, there are only finitely many possible subsets $\Pi\subseteq 2^{[m]}\setminus\{\varnothing\}$, and once $\Pi$ is fixed there are only finitely many maps $\chi:\Pi\to Q$.
Hence $\Theta_B(Q)$ is finite.
The transition rule from $v$ to $v\!\restriction \sigma$ depends only on the profile and on $\sigma$, and the selected branch symbol depends only on the same finite profile.
Therefore the selector is implemented by a finite recursive state space.

Conversely, suppose the selector already has a finite recursive presentation.
Take $Q$ to be the finite set of recursive states, and let $q(v)$ be the state assigned to the residual state $v$.
By definition of a correct recursive presentation, the realized pattern set, the child transitions, and the selected pattern all depend only on the current recursive state.
Thus $q$ is a finite hereditary quotient in the sense of Definition~\ref{def:hereditary-pattern-selector}.
\end{proof}

\begin{counterexample}[The greedy heaviest support-pattern selector need not be finite-state]\label{ce:greedy-support-branch}
Fix $k\ge 3$.
There exists a rigid class of $k$-sunflower-free uniform families with canonical bounded transversals and singleton support patterns for which the greedy heaviest support-pattern selector is \emph{not} hereditarily finite-state.
\end{counterexample}

\begin{proof}
We encode rooted trees into uniform families.

\smallskip
\noindent\emph{Step 1: the path-family encoding.}
Let $T$ be a finite rooted tree of maximum outdegree at most $k-1$, and let $N$ be its maximum root-to-leaf depth.
For each edge $e\in E(T)$ introduce a ground element $x_e$.
For each leaf $\ell$, let $P(\ell)\subseteq E(T)$ be the root-to-leaf edge set, and choose a filler set $U_\ell$ of size $N-|P(\ell)|$, with all filler sets pairwise disjoint and disjoint from the edge atoms $\{x_e:e\in E(T)\}$.
Define
\[
A_\ell:=\{x_e:e\in P(\ell)\}\cup U_\ell,
\qquad
\F_T:=\{A_\ell:\ell\in L(T)\}.
\]
Then every $A_\ell$ has size $N$, so $\F_T$ is $N$-uniform.

\smallskip
\noindent\emph{Step 2: $\F_T$ is $k$-sunflower-free.}
Take distinct leaves $\ell_1,\dots,\ell_k$.
Let $u$ be their deepest common vertex.
Since $u$ has at most $k-1$ children, two of the leaves, say $\ell_1$ and $\ell_2$, pass through the same child $c$ of $u$.
Because $u$ is the deepest common vertex, not all $k$ leaves can pass through that same child, so some $\ell_j$ passes through a different child.
The intersection $A_{\ell_1}\cap A_{\ell_2}$ contains the edge atom corresponding to the edge $(u,c)$, whereas $A_{\ell_1}\cap A_{\ell_j}$ does not.
Hence the pairwise intersections are not all equal, so $\{\ell_1,\dots,\ell_k\}$ does not form a sunflower.
Therefore $\F_T$ is $k$-sunflower-free.

\smallskip
\noindent\emph{Step 3: exact support-pattern recursion is canonical.}
For a vertex $u$, let $\F_T(u)$ be the residual family of leaves below $u$, after deleting the common root-to-$u$ path atoms.
If $u_1,\dots,u_r$ are the children of $u$, define
\[
H(u):=\{x_{(u,u_1)},\dots,x_{(u,u_r)}\}.
\]
Every residual set in $\F_T(u)$ contains exactly one element of $H(u)$.
Thus $H(u)$ is a canonical transversal, $|H(u)|=r\le k-1$, and the only realized support patterns are the singleton sets
\[
\bigl\{\{x_{(u,u_i)}\}:1\le i\le r\bigr\}.
\]
Conditioning on the pattern $\{x_{(u,u_i)}\}$ yields exactly the child subtree $T_{u_i}$.

Under the uniform measure on leaves, the mass of the $i$th child is
\[
\mu_u(i)=\frac{L(u_i)}{L(u)},
\]
where $L(u_i)$ is the number of descendant leaves below $u_i$.
Therefore the greedy heaviest support-pattern selector is exactly the heavy-path rule in the rooted tree.

\smallskip
\noindent\emph{Step 4: the hereditary congruence has infinite index.}
Work now inside the binary subcategory, which is allowed because every binary tree has outdegree at most $k-1$ when $k\ge 3$.
For each $m\ge 1$, let $T_m$ be a finite rooted binary tree with exactly $m$ leaves.
For $t\ge 1$, let $C_t[\square]$ denote the unary context obtained by attaching the input tree at the left child of a new root and a fixed tree $T_t$ at the right child.

In the family corresponding to $C_t[T_m]$, the first support-pattern choice compares the number of leaves in the two child subtrees.
Hence the greedy first move is:
\begin{enumerate}[label=(\alph*),leftmargin=2.2em]
\item left if $m>t$,
\item right if $m<t$,
\item and determined by the fixed tie-breaking rule if $m=t$.
\end{enumerate}

Suppose, for contradiction, that the greedy selector were hereditarily finite-state.
Then the support-pattern-context congruence would have finite index.
Applying the pigeonhole principle to the infinite sequence $T_2,T_4,T_6,\dots$, choose distinct even integers $a>b$ such that $T_a$ and $T_b$ lie in the same hereditary congruence class.
Set $t:=b+1$.
Then $b<t<a$, so in the context $C_t[\square]$ the greedy first move on $T_a$ is left, while the greedy first move on $T_b$ is right.
Thus $T_a$ and $T_b$ are distinguished by a unary context, contradicting that they belong to the same hereditary congruence class.

Therefore the hereditary support-pattern-context congruence has infinite index, and the greedy heaviest support-pattern selector is not hereditarily finite-state.
\end{proof}

\begin{remark}[What the counterexample actually says]\label{rem:greedy-support-selector}
In the tree model above, the emitted word language of greedy left/right choices is still regular: every word in $\{L,R\}^\ast$ occurs along some heavy path.
The obstruction is not the surface word language.
It is the hereditary continuation state, which remembers the descendant mass parameter tested by the contexts $C_t[\square]$.
Thus bounded transversals alone do \emph{not} put the greedy support-pattern branch inside the solved finite-state class.
What survives is the weaker finite-profile viewpoint of Proposition~\ref{prop:selector-finite-profile}, namely a finite hereditary solved cover of the \emph{full} exact support-pattern decomposition tree.
\end{remark}

\section{Fixed-alphabet code models: rigorous recurrences and provisional local tensors}\label{sec:fixed-alphabet}
In the fixed-alphabet transversal code model the most durable rigorous statements are the exact off-diagonal recurrences. Several diagonal local-tensor bounds remain exploratory because the current monomial counts use the wrong tail direction.

\subsection{Exact off-diagonal recurrences}
\begin{theorem}[One extra symbol]\label{thm:extra-symbol}
For all integers $q,k,m\ge 1$,
\[
F_{q+1,k}(m)\le \sum_{r=0}^m \binom{m}{r}F_{q,k}(r).
\]
More generally, for every integer $t\ge 0$,
\[
F_{q+t,k}(m)\le \sum_{r=0}^m \binom{m}{r} t^{m-r}F_{q,k}(r).
\]
\end{theorem}

\begin{proof}
Partition a code $C\subseteq [q+1]^m$ by the support of the extra symbol $q+1$. Deleting the extra-symbol coordinates from a fixed support class gives an injective map into $[q]^{m-r}$, and sunflower-freeness is preserved because deleted coordinates were constant. Summing over support sets yields the recurrence. The iterated formula follows by induction on $t$.
\end{proof}

\leannote{Lean status: the corresponding declarations \texttt{one\_extra\_symbol\_recurrence} and \texttt{iterated\_extra\_symbol\_recurrence} are present in \texttt{FormalConjectures/Problems/Erdos/E20/Compendium/Models/CodeModels.lean}, but both are currently marked \texttt{sorry}. No completed Lean proof is claimed for this theorem.}

\begin{theorem}[Sharper randomized deletion lift]\label{thm:randomized-deletion}
For all integers $q,t,k,m\ge 1$,
\[
F_{q+t,k}(m)\le \left(\frac{q+t}{q}\right)^m F_{q,k}(m).
\]
In particular,
\[
F_{6,5}(m)\le \left(\frac65\right)^m F_{5,5}(m).
\]
\end{theorem}

\begin{proof}
Choose independently in each coordinate a uniformly random $q$-subset of $[q+t]$. Any codeword survives with probability $(q/(q+t))^m$. The surviving code is identified with a subset of $[q]^m$ and remains $k$-sunflower-free. Taking expectations yields the inequality.
\end{proof}

\leannote{Lean status: the corresponding declaration \texttt{randomized\_deletion\_lift} is present in \texttt{FormalConjectures/Problems/Erdos/E20/Compendium/Models/CodeModels.lean}, but it is currently marked \texttt{sorry}. No completed Lean proof is claimed for this theorem.}

\begin{remark}[Consequences of a diagonal input]
The sharper lift shows that to force
\[
F_{6,5}(m)<5^m
\]
it is enough to prove
\[
F_{5,5}(m)<\left(\frac{25}{6}\right)^m\approx 4.1667^m.
\]
This is much weaker than proving a base strictly below $4$, which is impossible anyway because $[4]^m\subset[5]^m$ is already 5-sunflower-free.
\end{remark}


\subsection{Rigorous diagonal counterexamples and exact reductions}\label{subsec:exact-diagonal-reductions}
The diagonal fixed-alphabet models also admit several exact results that survive completely rigorous checking. They sharpen the status of the $(4,4)$ and $(5,5)$ cases even though they do not yet finish the corresponding global upper bounds.

\begin{counterexample}[The naive power-of-two diagonal bound already fails at $q=4$]\label{ce:444-power-of-two-false}
In the fixed-alphabet code model,
\[
F_{4,4}(3)\ge 28>27=3^3.
\]
Consequently the proposed formula
\[
F_{2^d,2^d}(m)\le (2^d-1)^m
\qquad(d\ge 2)
\]
is false.
\end{counterexample}

\begin{proof}
Set
\[
\begin{aligned}
A_0&:=\{(0,0),(0,1),(0,2),(2,2),(2,3),(3,0),(3,1),(3,2)\},\\
A_1&:=\{(0,0),(0,1),(0,3),(1,0),(2,1),(2,3),(3,0),(3,1),(3,3)\},\\
A_2&:=\{(0,0),(0,1),(0,2),(2,1),(2,2),(2,3),(3,0),(3,1),(3,2)\},\\
A_3&:=\{(2,1),(2,2)\},
\end{aligned}
\]
and let
\[
S:=(\{0\}\times A_0)\cup(\{1\}\times A_1)\cup(\{2\}\times A_2)\cup(\{3\}\times A_3)\subset [4]^3.
\]
Then $|S|=8+9+9+2=28$.

We claim that $S$ is $4$-sunflower-free. First note that a $4$-sunflower in $[4]^2$ is exactly one of the following: a full row, a full column, or a perfect matching in the $4\times 4$ grid. Indeed, with two coordinates, each column of the $4\times 2$ coordinate matrix must be either constant or injective.

Now each fiber $A_i\subset [4]^2$ is itself $4$-sunflower-free. The sets $A_0$ and $A_2$ have empty row $1$, so they contain no perfect matching and no full row or full column. The set $A_1$ has empty column $2$, so it also contains no perfect matching and no full row or full column. Finally, $A_3$ has size $2$.

Therefore a $4$-sunflower in $S$ cannot have constant first coordinate. So after reordering it would have the form
\[
(0,u_0),(1,u_1),(2,u_2),(3,u_3)
\qquad\text{with }u_i\in A_i.
\]
Write $u_i=(s_i,t_i)$. Since $A_3=\{(2,1),(2,2)\}$, we always have $s_3=2$. Because the second coordinate of a sunflower must be either constant or injective, there are two cases.

\emph{Case 1: $s_0,s_1,s_2,s_3$ are constant.}
Then all four equal $2$, so
\[
t_0\in\{2,3\},\qquad
 t_1\in\{1,3\},\qquad
 t_2\in\{1,2,3\},\qquad
 t_3\in\{1,2\}.
\]
These four values cannot all be equal, because the intersection of the four displayed sets is empty. They also cannot be all distinct, because only the values $1,2,3$ can occur. Thus the third coordinate is neither constant nor injective.

\emph{Case 2: $s_0,s_1,s_2,s_3$ are injective.}
Since only $A_1$ contains a point with first entry $1$, we must have $u_1=(1,0)$. Also $s_3=2$, so $s_0$ and $s_2$ are $0$ and $3$ in some order. Hence
\[
t_1=0,
\qquad t_3\in\{1,2\},
\qquad t_0,t_2\in\{0,1,2\}.
\]
Again the four values cannot all be equal, because $t_1=0$ while $t_3\in\{1,2\}$, and they cannot be all distinct because $3$ never appears.

Both cases are impossible. So $S$ is $4$-sunflower-free, and hence $F_{4,4}(3)\ge |S|=28$.
\end{proof}

\begin{lemma}[Product lemma for diagonal code models]\label{lem:code-product}
If $X\subseteq [q]^m$ and $Y\subseteq [q]^n$ are $q$-sunflower-free, then their concatenation product
\[
X\times Y:=\{(x,y):x\in X,\ y\in Y\}\subseteq [q]^{m+n}
\]
is also $q$-sunflower-free.
\end{lemma}

\begin{proof}
Suppose that $(x^{(1)},y^{(1)}),\dots,(x^{(q)},y^{(q)})\in X\times Y$ are distinct and form a $q$-sunflower. In the $X$-block, every coordinate is constant or injective, so the $q$ words $x^{(1)},\dots,x^{(q)}$ are either all equal or all distinct. The same is true for $y^{(1)},\dots,y^{(q)}$. Since the pairs are distinct, the two blocks cannot both be constant. Hence at least one block is all distinct, and that block gives a genuine $q$-sunflower in $X$ or in $Y$, a contradiction.
\end{proof}

\begin{corollary}[Asymptotic lower bound in the $(4,4)$ model]\label{cor:444-asymp-lower}
For every $t\ge 0$ and $r\in\{0,1,2\}$,
\[
F_{4,4}(3t+r)\ge 28^t3^r.
\]
In particular,
\[
\limsup_{m\to\infty} F_{4,4}(m)^{1/m}\ge 28^{1/3}\approx 3.03659>3.
\]
Under the standard transversal encoding,
\[
M(3t+r,4)\ge 28^t3^r.
\]
So any exponential bound of the form $M(n,4)\le c^n$ must satisfy $c\ge 28^{1/3}$.
\end{corollary}

\begin{proof}
The cube $[3]^r\subset [4]^r$ is trivially $4$-sunflower-free, and Counterexample~\ref{ce:444-power-of-two-false} provides a $28$-point $4$-sunflower-free code in length $3$. Applying Lemma~\ref{lem:code-product} repeatedly gives $F_{4,4}(3t+r)\ge 28^t3^r$. The limsup claim follows immediately. The standard transversal encoding from $[4]^m$ to $m$-uniform set families preserves the sunflower condition coordinatewise, so the same lower bound transfers to $M(m,4)$.
\end{proof}

\begin{theorem}[Exact two-moment characterization of admissible $5$-columns]\label{thm:5column-two-moment}
Let
\[
\mathcal S_5:=\{(a_1,\dots,a_5)\in \mathbb F_5^5: \text{the column is constant or injective}\}.
\]
Then for every $(a_1,\dots,a_5)\in \mathbb F_5^5$,
\[
(a_1,\dots,a_5)\in \mathcal S_5
\iff
\sum_{i=1}^5 a_i=0
\quad\text{and}\quad
\sum_{i=1}^5 a_i^2=0
\]
in $\mathbb F_5$.
\end{theorem}

\begin{proof}
Set
\[
s_1:=\sum_{i=1}^5 a_i,
\qquad
s_2:=\sum_{i=1}^5 a_i^2.
\]
If the column is constant, say $a_1=\cdots=a_5=a$, then $s_1=5a=0$ and $s_2=5a^2=0$ in characteristic $5$. If the column is injective, then $\{a_1,\dots,a_5\}=\mathbb F_5$, so
\[
s_1=\sum_{t\in \mathbb F_5} t=0,
\qquad
s_2=\sum_{t\in \mathbb F_5} t^2=0.
\]
Thus every element of $\mathcal S_5$ satisfies $s_1=s_2=0$.

Conversely assume $s_1=s_2=0$. Let $e_r$ denote the elementary symmetric polynomials in the $a_i$. Newton's identities give
\[
e_1=s_1=0,
\qquad
2e_2=e_1s_1-s_2=0,
\]
so $e_2=0$ as well. Therefore
\[
p(z):=\prod_{i=1}^5 (z-a_i)=z^5-e_3z^2+e_4z-e_5.
\]
For each $i$ we have $p(a_i)=0$, and since $a_i^5=a_i$ in $\mathbb F_5$,
\[
0=p(a_i)=a_i-e_3a_i^2+e_4a_i-e_5=-e_3a_i^2+(1+e_4)a_i-e_5.
\]
Thus every $a_i$ is a root of the quadratic polynomial
\[
q(z):=e_3z^2-(1+e_4)z+e_5.
\]
If $q\not\equiv 0$, then there are at most two distinct values among the $a_i$. If there are exactly two, say $r\neq s$ with multiplicities $u$ and $5-u$ for some $u\in\{1,2,3,4\}$, then
\[
0=s_1=ur+(5-u)s=u(r-s),
\]
which is impossible in $\mathbb F_5$ because $u\neq 0$ and $r\neq s$. Hence in this case all five entries are equal.

If instead $q\equiv 0$, then $e_3=0$, $e_4=-1$, and $e_5=0$, so
\[
p(z)=z^5-z.
\]
Its roots are exactly the five elements of $\mathbb F_5$, so the column is injective.

Therefore $s_1=s_2=0$ implies that the column is constant or injective, completing the proof.
\end{proof}

\begin{corollary}[Exact reduction of $F_{5,5}(m)$ to a zero-sum problem]\label{cor:555-moment-curve}
Identify $[5]$ with $\mathbb F_5$, and define
\[
P:=\{(t,t^2):t\in \mathbb F_5\}\subset \mathbb F_5^2,
\qquad
\Pi(x_1,\dots,x_m):=((x_1,x_1^2),\dots,(x_m,x_m^2))\in P^m.
\]
Then for $x^{(1)},\dots,x^{(5)}\in [5]^m$,
\[
\text{every coordinate column is constant or injective}
\iff
\Pi(x^{(1)})+\cdots+\Pi(x^{(5)})=0.
\]
Moreover, if $C\subseteq [5]^m$ is $5$-sunflower-free, then the only zero-sum $5$-tuples in $\Pi(C)$ are diagonals. Equivalently,
\[
F_{5,5}(m)=\max\Bigl\{|A|:A\subseteq P^m,\ \forall a_1,\dots,a_5\in A,
\ a_1+\cdots+a_5=0\Rightarrow a_1=\cdots=a_5\Bigr\}.
\]
Finally, the exact local indicator is
\[
\Theta\bigl(x^{(1)},\dots,x^{(5)}\bigr)
:=
\prod_{j=1}^m
\left(1-\left(\sum_{i=1}^5 x^{(i)}_j\right)^4\right)
\left(1-\left(\sum_{i=1}^5 (x^{(i)}_j)^2\right)^4\right),
\]
and for every $5$-sunflower-free code $C$ one has
\[
\Theta\big|_{C^5}=\mathbf 1_{\{x^{(1)}=\cdots=x^{(5)}\}}.
\]
\end{corollary}

\begin{proof}
The coordinatewise equivalence is exactly Theorem~\ref{thm:5column-two-moment}. If a zero-sum $5$-tuple in $\Pi(C)$ has two equal codewords, then in every coordinate the corresponding $5$-column has a repeated value, so it cannot be injective; hence every coordinate is constant, and all five codewords are equal. Thus every non-diagonal zero-sum $5$-tuple in $\Pi(C)$ would give five distinct codewords in $C$ whose coordinates are all equal or all distinct, namely a forbidden $5$-sunflower. The displayed extremal reformulation follows because $\Pi$ is a bijection from $[5]^m$ onto $P^m$. Finally, over $\mathbb F_5$ the factor $1-y^4$ is exactly the indicator of the event $y=0$, so the product formula for $\Theta$ is the exact local indicator given by Theorem~\ref{thm:5column-two-moment}; restricting to $C^5$ leaves only the diagonal.
\end{proof}

\leannote{The proved Lean counterpart records the local support and sunflower-free diagonality components of the \((5,5)\) tensor construction in \texttt{FormalConjectures/Problems/Erdos/E20/Tensors/QuinaryTensor.lean}, notably \texttt{explicitLocalTensor5\_eq\_indicator5} and \texttt{explicitGlobalTensor5Tuple\_eq\_ite\_allEqual\_of\_sunflowerFree}; \texttt{InformalMap.lean} exposes these as \texttt{pasted\_quinary\_local\_tensor\_exact\_support} and \texttt{pasted\_quinary\_global\_tensor\_diagonal\_on\_sunflower\_free}. The zero-sum extremal reformulation is not cited here as a separate completed Lean theorem.}

\subsection*{The exact $(5,5)$ moment-curve tensor: local obstructions and a sharp two-coordinate theorem}
Corollary~\ref{cor:555-moment-curve} reduces the $(5,5)$ model to a diagonal zero-sum problem on the quadratic moment curve. The next results record the strongest rigorous facts currently obtainable from the exact local tensor.

\begin{proposition}[Exact local correction gives the unit tensor]\label{prop:555-unit-tensor}
Let
\[
\mathcal R_5:=\{(t_1,\dots,t_5)\in \mathbb F_5^5:\ \sum_{i=1}^5 t_i=0,\ \sum_{i=1}^5 t_i^2=0\},
\]
and decompose it as
\[
D:=\{(a,a,a,a,a):a\in \mathbb F_5\},
\qquad
J:=\{(\sigma(0),\sigma(1),\sigma(2),\sigma(3),\sigma(4)):\sigma\in S_5\}.
\]
Then $\mathcal R_5=D\sqcup J$. If $z:=\mathbf 1_{\mathcal R_5}$ and
\[
u(t_1,\dots,t_5):=1+\sum_{i=1}^5 t_i^4,
\]
then
\[
z\cdot u=\mathbf 1_D.
\]
Consequently
\[
\Delta_m:=(z\cdot u)^{\otimes m}
\]
is the order-$5$ unit tensor of side length $5^m$, and therefore
\[
\operatorname{slice\text{-}rank}(\Delta_m)=5^m.
\]
In particular, any method that first kills the injective orbit exactly at one coordinate and then takes global slice rank cannot prove an exponential base below $5$ in the $(5,5)$ model.
\end{proposition}

\begin{proof}
The decomposition $\mathcal R_5=D\sqcup J$ is exactly the coordinatewise statement of Theorem~\ref{thm:5column-two-moment}. If $(t_1,\dots,t_5)\in D$, say $t_1=\cdots=t_5=a$, then
\[
\sum_{i=1}^5 t_i^4=5a^4=0
\]
in characteristic $5$, so $u=1$ on $D$. If $(t_1,\dots,t_5)\in J$, then the five values are exactly the elements of $\mathbb F_5$, so
\[
\sum_{i=1}^5 t_i^4=\sum_{a\in\mathbb F_5} a^4=4=-1
\]
and hence $u=0$ on $J$. Thus $z\cdot u$ is exactly the diagonal indicator $\mathbf 1_D$.

Tensoring over $m$ coordinates gives
\[
\Delta_m(x^{(1)},\dots,x^{(5)})=\mathbf 1_{\{x^{(1)}=\cdots=x^{(5)}\}}
\]
on $(\mathbb F_5^m)^5$. This is the order-$5$ unit tensor on a set of size $5^m$, whose slice rank is exactly $5^m$ by the standard diagonal-tensor lemma. The final statement is immediate.
\end{proof}

\begin{proposition}[Raw tensor and separable-weight obstructions in the $(5,5)$ model]\label{prop:555-raw-tensor-obstruction}
Let $Z_m:=z^{\otimes m}$ with $z=\mathbf 1_{\mathcal R_5}$.
\begin{enumerate}[label=(\roman*),leftmargin=2.2em]
\item The matrix flattening of $Z_m$ with the first factor as rows and the remaining four factors as columns has rank at least $5^m$.
\item There do not exist real weights $w_i(a)$, with $i\in\{1,\dots,5\}$ and $a\in\mathbb F_5$, such that all diagonal tuples $(a,a,a,a,a)$ have the same total weight while every injective tuple in $J$ has strictly larger total weight.
\end{enumerate}
Hence neither plain flattening-rank arguments on the raw tensor nor coordinatewise separable weight degenerations can isolate the diagonal from the injective orbit at exponential base below $5$.
\end{proposition}

\begin{proof}
For part (i), flatten $Z_m$ as
\[
\mathbb F_5^m \quad\text{versus}\quad (\mathbb F_5^m)^4.
\]
For each row index $x\in \mathbb F_5^m$, consider the column
\[
c(x):=(x,x,x,x).
\]
If $x'\in \mathbb F_5^m$, then
\[
Z_m(x',c(x))=\prod_{j=1}^m z(x'_j,x_j,x_j,x_j,x_j).
\]
For a single coordinate, the $5$-tuple $(a,b,b,b,b)$ lies in $\mathcal R_5=D\sqcup J$ if and only if $a=b$: it is diagonal when $a=b$, and it is neither diagonal nor injective when $a\neq b$. Therefore
\[
Z_m(x',c(x))=\prod_{j=1}^m \mathbf 1_{\{x'_j=x_j\}}=\mathbf 1_{\{x'=x\}}.
\]
So the flattening contains an identity submatrix of size $5^m$, and its rank is at least $5^m$.

For part (ii), write
\[
d_a:=\sum_{i=1}^5 w_i(a)
\]
for the weight of the diagonal tuple $(a,a,a,a,a)$. If all diagonal tuples have the same weight, then after subtracting $d_a/5$ from each of the five quantities $w_i(a)$ we may normalize to
\[
\sum_{i=1}^5 w_i(a)=0
\qquad\text{for every }a\in \mathbb F_5,
\]
without changing the comparison between diagonal and injective tuples. Now average the weights of the $120$ injective tuples:
\[
\frac{1}{120}\sum_{\sigma\in S_5}\sum_{i=1}^5 w_i(\sigma(i))
=
\frac{1}{120}\sum_{i=1}^5 \sum_{a\in\mathbb F_5} 24\,w_i(a)
=
\frac{1}{5}\sum_{a\in\mathbb F_5}\sum_{i=1}^5 w_i(a)
=
0.
\]
Thus the average injective weight is $0$, the same as the normalized diagonal weight, so at least one injective tuple has weight at most the diagonal weight. Strict separation is impossible.
\end{proof}

\begin{proposition}[One-block monomial counts on the exact $(5,5)$ tensor stall at base $5$]\label{prop:555-local-monomial-obstruction}
Let
\[
\mathscr R:=\mathbb C[P]\cong \mathbb C[u,v]/(v-u^2,\ u^5-u),
\]
graded by $\deg u=\deg v=1$. Then $\mathscr R$ has basis
\[
1,\ u,\ v,\ uv,\ v^2,
\]
so its Hilbert series is
\[
\operatorname{Hilb}_{\mathscr R}(z)=1+2z+2z^2.
\]
Define the exact local tensor on $P^5$ by
\[
\theta\bigl((u_i,v_i)_{i=1}^5\bigr):=
\left(1-\left(\sum_{i=1}^5 u_i\right)^4\right)
\left(1-\left(\sum_{i=1}^5 v_i\right)^4\right),
\]
and set
\[
T_m:=\theta^{\otimes m}.
\]
Then the standard one-block degree argument gives
\[
\operatorname{slice\text{-}rank}(T_m)\le 5\,N_m\!\left(\Bigl\lfloor\frac{8m}{5}\Bigr\rfloor\right),
\qquad
N_m(d):=\sum_{r\le d}[z^r](1+2z+2z^2)^m.
\]
Moreover
\[
N_m\!\left(\frac{8m}{5}\right)=5^{m-o(m)}.
\]
Therefore every unweighted one-block monomial/slice-rank estimate of this type has asymptotic base $5$. The same conclusion holds for every weighted one-block degree scheme obtained by assigning positive weights $\alpha,\beta$ to $u$ and $v$.
\end{proposition}

\begin{proof}
The basis statement is immediate from the relations $v=u^2$ and $u^5=u$. The five basis elements have degrees
\[
0,1,1,2,2,
\]
so the Hilbert series is indeed $1+2z+2z^2$.

Every monomial appearing in the local tensor $\theta$ has total degree at most $8$, and hence every monomial of $T_m=\theta^{\otimes m}$ has total degree at most $8m$ across the five tensor factors. In such a monomial, at least one factor has degree at most $8m/5$. The usual slice-rank pigeonhole decomposition therefore yields
\[
\operatorname{slice\text{-}rank}(T_m)\le 5\,N_m\!\left(\Bigl\lfloor\frac{8m}{5}\Bigr\rfloor\right).
\]

To estimate $N_m(d)$, let $D$ be the random variable taking the values $0,1,2$ with probabilities $1/5,2/5,2/5$. Then
\[
N_m(d)=5^m\,\Pr(D_1+\cdots+D_m\le d),
\]
where $D_1,\dots,D_m$ are independent copies of $D$. Since
\[
\E D=\frac{0+1+1+2+2}{5}=\frac{6}{5}
\qquad\text{and}\qquad
\frac{8}{5}>\frac{6}{5},
\]
the threshold $8m/5$ lies strictly above the mean. By the law of large numbers,
\[
\Pr\!\left(D_1+\cdots+D_m\le \frac{8m}{5}\right)\to 1,
\]
and therefore
\[
N_m\!\left(\frac{8m}{5}\right)=5^{m-o(m)}.
\]
This proves the unweighted statement.

For the weighted version, assign positive weights $\alpha,\beta$ to $u$ and $v$. The five basis functions then have minimal weights
\[
c_0=0,\quad
c_1=\alpha,\quad
c_2=\min(2\alpha,\beta),\quad
c_3=\min(3\alpha,\alpha+\beta),\quad
c_4=\min(4\alpha,2\alpha+\beta,2\beta).
\]
Every monomial of $\theta$ has total weighted degree at most $4(\alpha+\beta)$, so the usual one-factor pigeonhole cutoff occurs at
\[
\frac{4(\alpha+\beta)}{5}\,m.
\]
On the other hand,
\[
c_0+c_1+c_2+c_3+c_4
\le
0+\alpha+\beta+(\alpha+\beta)+2\beta
=
2\alpha+4\beta
\le
4\alpha+4\beta.
\]
Hence the average local basis weight is at most $4(\alpha+\beta)/5$, so the relevant weighted lower-tail space is again cut at or above its mean. The same law-of-large-numbers argument shows that the weighted low-degree space has size $5^{m-o(m)}$. Thus no one-block weighted degree scheme of this type can produce a base below $5$.
\end{proof}



\begin{remark}[Support-only marginals show no local defect]\label{rem:555-support-entropy}
Give equal mass to each of the five diagonal tuples in $D$ and equal mass to each of the $120$ injective tuples in $J$. Then every one-coordinate marginal is uniform on $\mathbb F_5$: each symbol appears once among the diagonal tuples and $24$ times in each coordinate position among the injective tuples. Thus any support-only statistic depending only on one-coordinate marginals already sees the maximal local entropy value. The obstruction in the $(5,5)$ model is therefore genuinely non-marginal.
\end{remark}

\begin{remark}[Recursive form of the remaining tensorization problem]\label{rem:555-recursive-obstacle}
If
\[
A=\bigsqcup_{a\in \mathbb F_5} \{a\}\times B_a
\qquad\text{with }B_a\subseteq \mathbb F_5^{m-1},
\]
then $A$ is $5$-sunflower-free if and only if each $B_a$ is $5$-sunflower-free and there is no tuple
\[
(b_0,\dots,b_4)\in B_0\times\cdots\times B_4
\]
satisfying
\[
b_0+\cdots+b_4=0,
\qquad
b_0^2+\cdots+b_4^2=0
\]
coordinatewise. Thus the real remaining issue is a five-family transversal / Hall-type tensorization theorem rather than any further one-coordinate local calculation.
\end{remark}


\begin{theorem}[Sharp two-coordinate value in the $(5,5)$ model]\label{thm:555-two-coordinate}
One has
\[
F_{5,5}(2)=16.
\]
\end{theorem}

\begin{proof}
Identify a set $A\subseteq [5]^2$ with the bipartite graph $G_A$ on row set $[5]$ and column set $[5]$, where $(r,c)\in A$ corresponds to an edge between row $r$ and column $c$.

A nontrivial bad $5$-tuple in $[5]^2$ can occur in exactly three ways:
\begin{enumerate}[label=(\alph*),leftmargin=2.2em]
\item the first coordinate is constant and the second coordinate is injective, giving a full row of five points;
\item the first coordinate is injective and the second coordinate is constant, giving a full column of five points;
\item both coordinates are injective, giving the graph of a permutation, i.e. a perfect matching of size $5$.
\end{enumerate}
Therefore, if $A$ is $5$-sunflower-free, then every row has size at most $4$, every column has size at most $4$, and the bipartite graph $G_A$ has no perfect matching.

By Hall's theorem, because $G_A$ has no perfect matching, there exists a nonempty set $S$ of rows such that
\[
|N(S)|\le |S|-1.
\]
If $|S|\le 4$, then all edges either lie between $S$ and $N(S)$ or come from rows outside $S$. Hence
\[
|A|\le |S|\,|N(S)| + 4(5-|S|)
\le |S|(|S|-1)+4(5-|S|).
\]
For $|S|=1,2,3,4$, the right-hand side is at most $16$.

If instead $|S|=5$, then $|N(S)|\le 4$, so all edges lie in at most four columns. Since every column has size at most $4$,
\[
|A|\le 4\cdot 4=16.
\]
Thus every $5$-sunflower-free set $A\subseteq [5]^2$ has size at most $16$.

For the lower bound, the grid $[4]^2\subseteq [5]^2$ has size $16$ and is $5$-sunflower-free: it contains neither a full row nor a full column of length $5$, and it cannot contain a perfect matching on all five rows and five columns because only four symbols are available in each coordinate. Therefore $F_{5,5}(2)=16$.
\end{proof}

\begin{corollary}[Two-coordinate consequences for $F_{6,5}$ and for general $m$]\label{cor:555-two-coordinate}
One has
\[
F_{6,5}(2)\le \left(\frac65\right)^2 16 = \frac{576}{25}<25,
\]
hence in fact $F_{6,5}(2)\le 23$. More generally, for every $m\ge 2$,
\[
F_{5,5}(m)\le 16\cdot 5^{m-2}
\qquad\text{and}\qquad
F_{6,5}(m)\le \frac{16}{25}\,6^m.
\]
\end{corollary}

\begin{proof}
The two-coordinate bound for $F_{6,5}(2)$ follows from Theorem~\ref{thm:randomized-deletion} together with Theorem~\ref{thm:555-two-coordinate}. Since $F_{6,5}(2)$ is an integer, the strict inequality $F_{6,5}(2)<25$ implies $F_{6,5}(2)\le 23$.

For the general $F_{5,5}(m)$ bound, fix the last $m-2$ coordinates and view the fiber over those coordinates as a subset of $[5]^2$. Each such fiber is still $5$-sunflower-free, so by Theorem~\ref{thm:555-two-coordinate} it has size at most $16$. Summing over the $5^{m-2}$ choices of the remaining coordinates gives
\[
F_{5,5}(m)\le 16\cdot 5^{m-2}.
\]
Applying Theorem~\ref{thm:randomized-deletion} once more yields
\[
F_{6,5}(m)\le \left(\frac65\right)^m F_{5,5}(m)\le \frac{16}{25}\,6^m.
\]
\end{proof}

\begin{definition}[Five-partite disjoint cross parameter]\label{def:gamma555}
Let $\Gamma_{5}^{\times}(m)$ be the maximum value of
\[
\bigl(|A_1|\cdots |A_5|\bigr)^{1/5}
\]
over all pairwise disjoint sets $A_1,\dots,A_5\subseteq \mathbb F_5^m$ such that there is no tuple
\[
(a_1,\dots,a_5)\in A_1\times\cdots\times A_5
\]
with
\[
a_1+\cdots+a_5=0,
\qquad
a_1^2+\cdots+a_5^2=0
\]
coordinatewise.
\end{definition}

\begin{proposition}[Reduction to the disjoint five-partite cross model]\label{prop:gamma555-reduction}
For every $m\ge 1$,
\[
F_{5,5}(m)\le 5\,\Gamma_{5}^{\times}(m)+4.
\]
Consequently, any bound of the form
\[
\Gamma_{5}^{\times}(m)\le C^m
\qquad\text{with }C<25/6
\]
would imply a strict exponential base below $25/6$ for the ordinary $(5,5)$ problem.
\end{proposition}

\begin{proof}
Let $A\subseteq \mathbb F_5^m$ be a $5$-sunflower-free code of maximum size, so $|A|=F_{5,5}(m)$. Partition $A$ into five disjoint parts
\[
A=A_1\sqcup\cdots\sqcup A_5
\]
as equally as possible, so that
\[
\left\lfloor \frac{|A|}{5}\right\rfloor \le |A_i|\le \left\lceil \frac{|A|}{5}\right\rceil
\qquad\text{for all }i.
\]
Because the parts are pairwise disjoint, any tuple $(a_1,\dots,a_5)\in A_1\times\cdots\times A_5$ is automatically non-diagonal. If it satisfied the two moment equations, Corollary~\ref{cor:555-moment-curve} would produce a forbidden $5$-sunflower in $A$. Thus the five-tuple $(A_1,\dots,A_5)$ is admissible for $\Gamma_{5}^{\times}(m)$, and so
\[
\left\lfloor \frac{|A|}{5}\right\rfloor
\le
\bigl(|A_1|\cdots |A_5|\bigr)^{1/5}
\le
\Gamma_{5}^{\times}(m).
\]
Hence $|A|\le 5\,\Gamma_{5}^{\times}(m)+4$, proving the first claim. The final consequence is immediate.
\end{proof}

\begin{proposition}[The disjoint five-partite cross model has asymptotic base $5$]\label{prop:gamma555-base5}
For every $m\ge 2$,
\[
5^{m-2}\le \Gamma_{5}^{\times}(m)\le 5^{m-1}.
\]
In particular,
\[
\lim_{m\to\infty}\bigl(\Gamma_{5}^{\times}(m)\bigr)^{1/m}=5.
\]
Consequently no estimate of the form
\[
\Gamma_{5}^{\times}(m)\le C^m
\qquad\text{with }C<5
\]
is possible, so Proposition~\ref{prop:gamma555-reduction} cannot by itself yield a base below $25/6$ for the ordinary $(5,5)$ problem.
\end{proposition}

\begin{proof}
We first classify the one-coordinate forbidden relation. Let $(x_1,\dots,x_5)\in \mathbb F_5^5$ satisfy
\[
x_1+\cdots+x_5=0,
\qquad
x_1^2+\cdots+x_5^2=0.
\]
We claim that either all five entries are equal or they are pairwise distinct.

If the multiplicity pattern is $4+1$, write the tuple as $(a,a,a,a,b)$. Then
\[
4a+b=0,
\]
so $b=a$ and the tuple is constant. If the pattern is $3+2$, write it as $(a,a,a,b,b)$. Then
\[
3a+2b=0,
\]
which rearranges to $2(b-a)=0$, hence again $b=a$.

If the pattern is $3+1+1$, write it as $(a,a,a,b,c)$ and set $u:=b-a$, $v:=c-a$. The two equations become
\[
u+v=0,
\qquad
u^2+v^2=0.
\]
Since $v=-u$, the second equation gives $2u^2=0$. Because $2\neq 0$ in $\mathbb F_5$, we get $u=0$ and hence $b=c=a$.

If the pattern is $2+2+1$, write it as $(a,a,b,b,c)$ and set $u:=a-c$, $v:=b-c$. Then the two equations become
\[
2u+2v=0,
\qquad
2u^2+2v^2=0.
\]
Thus $v=-u$, and the second equation gives $4u^2=0$, so again $u=0$ and the tuple is constant.

The only remaining possibilities are the constant pattern $5$ and the pattern $1+1+1+1+1$ of pairwise distinct entries. Conversely, every constant tuple satisfies both equations trivially, and every pairwise distinct tuple is a permutation of $(0,1,2,3,4)$, whose coordinate sum and sum of squares are both $0$ in $\mathbb F_5$. Therefore the forbidden one-coordinate columns are exactly the constant columns and the all-distinct columns.

For the upper bound, if $A_1,\dots,A_5\subseteq \mathbb F_5^m$ are pairwise disjoint, then
\[
|A_1|+\cdots+|A_5|\le 5^m.
\]
By AM-GM,
\[
\bigl(|A_1|\cdots |A_5|\bigr)^{1/5}
\le
\frac{|A_1|+\cdots+|A_5|}{5}
\le
5^{m-1}.
\]
Taking the maximum over all admissible five-tuples gives $\Gamma_{5}^{\times}(m)\le 5^{m-1}$.

For the lower bound, define five signatures in $\mathbb F_5^2$ by
\[
s_1=(0,0),\quad s_2=(0,1),\quad s_3=(1,0),\quad s_4=(1,1),\quad s_5=(2,2),
\]
and for $m\ge 2$ set
\[
A_i:=\{x\in \mathbb F_5^m:(x_1,x_2)=s_i\}.
\]
These cylinder sets are pairwise disjoint and satisfy $|A_i|=5^{m-2}$. For any choice $a_i\in A_i$, the first coordinate column is
\[
(0,0,1,1,2),
\]
which is neither constant nor all-distinct. By the classification above, the local relation already fails in the first coordinate, so the rainbow tuple $(a_1,\dots,a_5)$ is not forbidden. Hence $(A_1,\dots,A_5)$ is admissible and
\[
\Gamma_{5}^{\times}(m)\ge \bigl(5^{m-2}\cdots 5^{m-2}\bigr)^{1/5}=5^{m-2}.
\]
This proves the proposition.
\end{proof}

\begin{proposition}[Perfect-matching data alone cannot tensorize the two-coordinate theorem]\label{prop:555-perfect-matching-lineality}
Let $\mathcal M_5$ be the set of perfect matchings in the grid $[5]\times [5]$. There is no absolute constant $C$ such that every nonnegative weight table $w:[5]^2\to \R_{\ge 0}$ satisfying
\[
\prod_{(x,y)\in M} w(x,y)\le 1
\qquad\text{for every }M\in \mathcal M_5
\]
must also satisfy
\[
\sum_{x,y\in [5]} w(x,y)\le C.
\]
\end{proposition}

\begin{proof}
Choose real numbers $r_1,\dots,r_5,c_1,\dots,c_5$ with
\[
\sum_{x=1}^5 r_x+\sum_{y=1}^5 c_y=0,
\]
and define
\[
w(x,y):=e^{r_x+c_y}.
\]
If $M$ is a perfect matching, then it contains each row and each column exactly once, so
\[
\prod_{(x,y)\in M} w(x,y)
=
\exp\!\left(\sum_{x=1}^5 r_x+\sum_{y=1}^5 c_y\right)=1.
\]
Thus every perfect-matching constraint holds with equality.

On the other hand,
\[
\sum_{x,y\in [5]} w(x,y)
=
\left(\sum_{x=1}^5 e^{r_x}\right)\left(\sum_{y=1}^5 e^{c_y}\right).
\]
Taking
\[
(r_1,r_2,r_3,r_4,r_5)=(t,-t,0,0,0),
\qquad
c_1=\cdots=c_5=0,
\]
we obtain
\[
\sum_{x,y} w(x,y)=5(e^t+e^{-t}+3)\to\infty
\qquad\text{as }t\to\infty.
\]
So no such constant $C$ can exist.
\end{proof}

\begin{remark}[What Proposition~\ref{prop:555-perfect-matching-lineality} does and does not say]\label{rem:555-perfect-matching}
The proposition concerns a globalization scheme that remembers only perfect-matching inequalities on the $5\times 5$ two-block fibers. The actual disjoint-cross model also has row-constant and column-constant rainbow obstructions, so the proposition does \emph{not} refute every possible two-block tensorization. What it does show is that perfect-matching data alone leave a $9$-dimensional row/column-potential lineality and therefore cannot control the total fiber mass.
\end{remark}



\begin{proposition}[The two-block Hall quotient for $(5,5)$ has dimension $16$]\label{prop:555-hall-quotient}
Let $M=(m_{ij})$ be a positive $5\times 5$ fiber matrix.
Quotienting by independent positive row and column scalings gives the local Hall state space
\[
Q_2:=\R^{25}/(\text{row lineality}+\text{column lineality}),
\]
which has dimension
\[
25-(5+5-1)=16.
\]
A concrete coordinate system is given by the cross-ratios
\[
q_{ab}(M):=\log\frac{m_{ab}m_{55}}{m_{a5}m_{5b}},
\qquad
1\le a,b\le 4.
\]
\end{proposition}

\begin{proof}
Independent row and column scalings act by
\[
m_{ij}\longmapsto \lambda_i\,m_{ij}\,\rho_j,
\]
so the space of logarithms is quotiented by the $(5+5-1)$-dimensional subspace generated by row and column potentials, with one overlap coming from simultaneous global scaling.
This leaves dimension $16$.

For $1\le a,b\le 4$, the quantity
\[
\log m_{ab}+\log m_{55}-\log m_{a5}-\log m_{5b}
\]
is invariant under row and column scalings, because the row factor of row $a$ cancels, the row factor of row $5$ cancels, and similarly for the column factors.
There are exactly $4\cdot 4=16$ such independent invariants, so they form a coordinate system on the quotient.
\end{proof}

\begin{proposition}[Bare two-block Hall profiles do not compose]\label{prop:555-hall-noncomposition}
There is no binary operation $\star$ on the bare Hall quotient $Q_2$ such that
\[
q(MN)=q(M)\star q(N)
\]
for every pair of positive $5\times 5$ fiber matrices $M,N$.
Equivalently, the isolated two-block Hall profile is not functorially composable.
\end{proposition}

\begin{proof}
Let $J$ be the $5\times 5$ all-ones matrix, let $E_{11}$ be the matrix with a single $1$ in the $(1,1)$ position and zeros elsewhere, and set
\[
M:=J+E_{11},
\qquad
N:=J+E_{11}.
\]
For $t>0$, define
\[
E_t:=\diag(t,1,1,1,1).
\]
Right multiplication by $E_t$ is a column scaling on the shared block, so the bare Hall profile of $ME_t$ is the same as that of $M$.
The profile of $N$ is independent of $t$.

Now contract across the shared block and write
\[
P_t:=ME_tN.
\]
A direct computation gives
\[
P_t=
\begin{pmatrix}
4t+4 & 2t+4 & 2t+4 & 2t+4 & 2t+4\\
2t+4 & t+4  & t+4  & t+4  & t+4\\
2t+4 & t+4  & t+4  & t+4  & t+4\\
2t+4 & t+4  & t+4  & t+4  & t+4\\
2t+4 & t+4  & t+4  & t+4  & t+4
\end{pmatrix}.
\]
Therefore
\[
q_{11}(P_t)
=
\log\frac{(4t+4)(t+4)}{(2t+4)^2}
=
\log\frac{(t+1)(t+4)}{(t+2)^2}.
\]
This depends nontrivially on $t$.
So the composite Hall profile of $P_t$ changes with $t$, even though the bare input profile of $M$ is unchanged by inserting the shared diagonal gauge $E_t$ and the profile of $N$ is fixed.
Hence no binary operation on bare Hall profiles can recover the composite profile.
\end{proof}

\begin{remark}[Restoring the missing gauge reconstructs the full local fiber]\label{rem:555-hall-gauge}
The obstruction in Proposition~\ref{prop:555-hall-noncomposition} is exactly the missing diagonal gauge on the shared block.
Indeed, once one records
\[
\ell_a:=\log\frac{m_{a5}}{m_{55}},
\qquad
r_b:=\log\frac{m_{5b}}{m_{55}},
\qquad
\sigma:=\log m_{55},
\]
together with the $16$ cross-ratios $q_{ab}(M)$, every matrix entry is recovered from
\[
\log m_{ab}=q_{ab}(M)+\ell_a+r_b+\sigma,
\qquad
\log m_{a5}=\ell_a+\sigma,
\qquad
\log m_{5b}=r_b+\sigma.
\]
So the minimal compositional completion of the bare Hall quotient is the full $25$-parameter local fiber.
\end{remark}

\begin{corollary}[Conditional $4^m$ bound from global Hall flatness]\label{cor:555-hall-flatness}
Assume that extremal objects in the disjoint five-partite cross model admit a simultaneous gauge fixing across a coordinate decomposition, so that two-block Hall profiles compose functorially after the hidden diagonal gauge has been eliminated globally.
Then
\[
\Gamma_5^\times(m)\le 4^m.
\]
Consequently,
\[
F_{5,5}(m)\le 5\cdot 4^m+4.
\]
\end{corollary}

\begin{proof}
For one coordinate, the row/column quotient has dimension
\[
5-1=4.
\]
For two coordinates, Proposition~\ref{prop:555-hall-quotient} gives local dimension $16$.
Write
\[
m=2t+s,
\qquad
s\in\{0,1\}.
\]
Pair the coordinates into $t$ two-blocks and, if necessary, one single block.
Under the assumed global gauge flatness, the corresponding Hall-profile state space has dimension
\[
16^t 4^s = 4^m.
\]
So the disjoint five-partite cross parameter is at most $4^m$.

The reduction from the ordinary $(5,5)$ problem to $\Gamma_5^\times(m)$ was proved earlier in Proposition~\ref{prop:gamma555-reduction}, which gives
\[
F_{5,5}(m)\le 5\,\Gamma_5^\times(m)+4.
\]
Substituting the bound $\Gamma_5^\times(m)\le 4^m$ gives the stated estimate.
\end{proof}

\begin{remark}[What remains in the $(5,5)$ model]\label{rem:555-what-remains}
Propositions~\ref{prop:555-unit-tensor}--\ref{prop:555-local-monomial-obstruction} show that the most obvious coordinatewise tensor routes stall at base $5$ for structural reasons, and Theorem~\ref{thm:555-two-coordinate} shows that the exact local geometry already beats the critical threshold at block length $2$. Proposition~\ref{prop:gamma555-base5} shows, however, that the coarse five-partite disjoint-cross proxy is already too large: its exponential base is $5$. Proposition~\ref{prop:555-perfect-matching-lineality} further shows that a tensorization which remembers only perfect-matching data from the two-block fibers cannot work. Any successful globalization must therefore retain more structure than pairwise disjoint color classes and more information than one scalar per two-block fiber.
\end{remark}

\status{The exact $(5,5)$ moment-curve reduction is rigorous, and the local tensor now comes with several rigorous impossibility results. The five-partite disjoint-cross proxy has asymptotic base $5$, so it cannot by itself deliver a base below $25/6$. Any successful route must retain more same-family structure than pairwise disjoint color classes and more global information than perfect-matching fiber counts.}

\subsection*{Aligned Hall templates and gauge-aware globalization}
The next exact step in the $(5,5)$ lane is to separate the already-rigorous aligned-template tensor law from the still-open global alignment problem.

\begin{definition}[Aligned Hall pieces and template parameters]\label{def:555-aligned-template}
Work in the block model $U=[5]^2$.
For $a,b\in [5]$ define the Hall strips and rectangles
\[
R_a:=([5]\setminus\{a\})\times [5],
\qquad
C_b:=[5]\times([5]\setminus\{b\}),
\qquad
Q_{a,b}:=R_a\cap C_b.
\]
An \emph{aligned Hall template} is a product
\[
T=S_1\times\cdots\times S_m\subseteq U^m
\]
in which each $S_j$ is one of the sets $R_a$, $C_b$, or $Q_{a,b}$.
Write
\[
r(T):=\#\{j:S_j\text{ is a rectangle }Q_{a,b}\},
\qquad
s(T):=\#\{j:S_j\text{ is a proper strip }R_a\text{ or }C_b\},
\]
so that $r(T)+s(T)=m$.
Let $M(T)$ be the maximum size of a good-free family $A\subseteq T$.
For the five-partite additive parameter, write
\[
\Xi_5(n):=
\max\Bigl\{\sum_{i=1}^5 |B_i|:
B_1,\dots,B_5\subseteq [5]^n\text{ are pairwise disjoint and admit no colorful good }5\text{-tuple}\Bigr\},
\]
and define $\Xi(T)$ analogously inside $T$.
\end{definition}

\begin{theorem}[Exact Hall-template tensorization]\label{thm:555-hall-template-tensorization}
With the notation above, every aligned Hall template $T$ satisfies
\[
M(T)=16^{\,r(T)}4^{\,s(T)}F_{5,5}(s(T)),
\]
and
\[
\Xi(T)=16^{\,r(T)}4^{\,s(T)}\Xi_5(s(T)).
\]
\end{theorem}

\begin{proof}
The local good relation is invariant under permuting row symbols, permuting column symbols, and swapping row and column in a given block.
So after blockwise relabeling we may assume every strip is the standard strip
\[
R=[4]\times[5]
\]
and every rectangle is the standard rectangle
\[
Q=[4]\times[4].
\]
Thus there is no loss in writing
\[
T=Q^r\times R^s
\]
with $r=r(T)$ and $s=s(T)$.

Every point of $Q^r\times R^s$ can be written uniquely in the form
\[
((u_1,v_1),\dots,(u_r,v_r),(a_1,b_1),\dots,(a_s,b_s)),
\]
where
\[
u_j,v_j,a_t\in [4],
\qquad
b_t\in [5].
\]
Hence we may identify
\[
T\cong \Lambda\times [5]^s,
\qquad
\Lambda:=[4]^{2r+s},
\]
by the map
\[
((u_1,v_1),\dots,(u_r,v_r),(a_1,b_1),\dots,(a_s,b_s))
\longmapsto
\bigl((u_1,v_1,\dots,u_r,v_r,a_1,\dots,a_s),(b_1,\dots,b_s)\bigr).
\]
In particular,
\[
|\Lambda|=4^{2r+s}=16^r4^s.
\]

We claim that a $5$-tuple
\[
x^1,\dots,x^5\in T\cong \Lambda\times [5]^s,
\qquad
x^\ell=(\lambda^\ell,z^\ell),
\]
is good in the block model if and only if the following two conditions hold:
\begin{enumerate}[label=(\roman*),leftmargin=2.2em]
\item $\lambda^1=\cdots=\lambda^5$;
\item $z^1,\dots,z^5\in [5]^s$ form a good $5$-tuple for the ordinary one-dimensional $(5,5)$ relation.
\end{enumerate}

Indeed, consider first a rectangle block $Q=[4]\times [4]$.
Neither coordinate can be injective on five points, so local goodness forces both the row and the column to be constant.
Thus all five points agree in that block.

Now consider a strip block $R=[4]\times [5]$.
Its first coordinate lives in $[4]$, so it cannot be injective on five points and must therefore be constant.
Its second coordinate lives in $[5]$, so it may be either constant or injective.
Thus the restricted coordinate is frozen, while the free coordinate behaves exactly like one coordinate of the ordinary $[5]^s$ problem.
Doing this block by block proves the claim.

We first prove the formula for $M(T)$.
Let $A\subseteq T$ be good-free.
For each $\lambda\in \Lambda$, define the fiber
\[
A_\lambda:=\{z\in [5]^s:(\lambda,z)\in A\}.
\]
By the claim above, each $A_\lambda$ is good-free in the ordinary $(5,5)$ model.
Therefore
\[
|A_\lambda|\le F_{5,5}(s)
\qquad\text{for every }\lambda\in \Lambda.
\]
Summing over the fibers gives
\[
|A|=\sum_{\lambda\in \Lambda}|A_\lambda|
\le |\Lambda|\,F_{5,5}(s)
=16^r4^sF_{5,5}(s).
\]

To see that this upper bound is sharp, choose an extremal good-free family $B\subseteq [5]^s$ with $|B|=F_{5,5}(s)$, and set
\[
A:=\Lambda\times B\subseteq \Lambda\times [5]^s\cong T.
\]
If $A$ contained a good $5$-tuple, then by the claim all five points would lie in a single fiber $\{\lambda\}\times [5]^s$, and their second coordinates would form a good $5$-tuple in $B$, contradiction.
Hence
\[
M(T)=16^r4^sF_{5,5}(s).
\]

The five-partite additive formula is proved in the same way.
Let $A_1,\dots,A_5\subseteq T$ be pairwise disjoint with no colorful good $5$-tuple.
For each $\lambda\in \Lambda$, define
\[
A_{i,\lambda}:=\{z\in [5]^s:(\lambda,z)\in A_i\}.
\]
The families $A_{1,\lambda},\dots,A_{5,\lambda}$ are pairwise disjoint in $[5]^s$, and by the same claim they admit no colorful good $5$-tuple.
Hence
\[
\sum_{i=1}^5 |A_{i,\lambda}|\le \Xi_5(s)
\qquad\text{for every }\lambda\in \Lambda.
\]
Summing over $\lambda$ yields
\[
\sum_{i=1}^5 |A_i|
=
\sum_{\lambda\in \Lambda}\sum_{i=1}^5 |A_{i,\lambda}|
\le
|\Lambda|\,\Xi_5(s)
=
16^r4^s\Xi_5(s).
\]
This proves
\[
\Xi(T)\le 16^r4^s\Xi_5(s).
\]

For the matching lower bound, choose extremal pairwise disjoint colorful-good-free families
\[
B_1,\dots,B_5\subseteq [5]^s
\]
with
\[
\sum_{i=1}^5 |B_i|=\Xi_5(s),
\]
and set
\[
A_i:=\Lambda\times B_i.
\]
These sets are pairwise disjoint, and the claim again rules out any colorful good tuple.
So
\[
\Xi(T)=16^r4^s\Xi_5(s).
\]
\end{proof}

\begin{corollary}[Aligned Hall covers reduce the global count to one-dimensional residuals]\label{cor:555-hall-template-cover}
If a family $A\subseteq U^m$ is covered by aligned Hall templates $T_1,\dots,T_L$, then
\[
|A|
\le
\sum_{\ell=1}^L 16^{\,r(T_\ell)}4^{\,s(T_\ell)}F_{5,5}(s(T_\ell)).
\]
Likewise, if pairwise disjoint families $A_1,\dots,A_5\subseteq U^m$ are covered by aligned templates $T_1,\dots,T_L$, then
\[
\sum_{i=1}^5 |A_i|
\le
\sum_{\ell=1}^L 16^{\,r(T_\ell)}4^{\,s(T_\ell)}\Xi_5(s(T_\ell)).
\]
\end{corollary}

\begin{proof}
Decompose each family into its intersections with the templates and apply Theorem~\ref{thm:555-hall-template-tensorization} on each template separately.
\end{proof}

\begin{definition}[Gauge-aware Hall relation]\label{def:555-gauge-relation}
Let $\mathcal N$ be the exact two-block Hall normal form coming from the sharp two-coordinate theorem in the $(5,5)$ model.
For a positive $5\times 5$ transfer block $M$, define the relation
\[
R_M\subset \mathbb P^4_{>0}\times \mathbb P^4_{>0}
\]
by
\[
([x],[y])\in R_M
\iff
D(x)^{-1}MD(y)\in \mathcal N,
\]
where $D(x)$ denotes the diagonal matrix with diagonal entries $x_1,\dots,x_4,1$.
\end{definition}

\begin{proposition}[Balanced Hall-flow dual]\label{prop:555-balanced-hall-flow}
Assume the normal form $\mathcal N$ is cut out in logarithmic coordinates by finitely many affine Hall inequalities. Then an admissible chain of restored transfer blocks
\[
M_1,\dots,M_r
\]
admits simultaneous gauge fixing if and only if there is no nonzero balanced Hall-flow with negative defect along that chain.
\end{proposition}

\begin{proof}
Write the defining affine inequalities of $\mathcal N$ in the form
\[
\langle r_\nu,y\rangle-\langle \ell_\nu,x\rangle
\le
c_\nu-\sum_{a,b}\alpha_{ab}^{(\nu)}\log (M_{ab})
\qquad (\nu\in \mathcal I)
\]
for suitable finite data $(\ell_\nu,r_\nu,c_\nu,\alpha^{(\nu)})$.
Fix a chain $M_1,\dots,M_r$.
Choosing projective interface gauges $[x_0],\dots,[x_r]$ with
\[
D(x_{j-1})^{-1}M_jD(x_j)\in \mathcal N
\qquad (1\le j\le r)
\]
is therefore equivalent to the feasibility of the finite linear system
\[
\langle r_\nu,x_j\rangle-\langle \ell_\nu,x_{j-1}\rangle
\le
c_\nu-\sum_{a,b}\alpha_{ab}^{(\nu)}\log (M_j)_{ab}
\qquad (1\le j\le r,\ \nu\in \mathcal I)
\]
in the logarithmic gauge variables.

By Farkas' lemma, this system is infeasible if and only if there exist nonnegative multipliers
\[
\lambda_{j,\nu}\ge 0,
\qquad
\text{not all zero,}
\]
such that the gauge terms balance:
\[
\sum_\nu \lambda_{j,\nu} r_\nu
=
\sum_\nu \lambda_{j+1,\nu}\ell_\nu
\qquad (1\le j<r),
\]
and the total defect is strictly negative:
\[
\sum_{j,\nu}\lambda_{j,\nu}
\left(
c_\nu-\sum_{a,b}\alpha_{ab}^{(\nu)}\log (M_j)_{ab}
\right)
<0.
\]
This is exactly the balanced Hall-flow obstruction.
Thus simultaneous gauge fixing exists if and only if no such negative-defect Hall-flow exists.
\end{proof}

\begin{theorem}[Gauge-aware Hall globalization criterion]\label{thm:555-gauge-aware-globalization}
Assume every admissible chain of restored $5\times 5$ transfer blocks in the disjoint five-partite cross model admits simultaneous gauge fixing; equivalently, by Proposition~\ref{prop:555-balanced-hall-flow}, assume that no admissible chain carries a nonzero balanced Hall-flow with negative defect.
Then
\[
\Gamma_5^\times(m)\le 4^m
\qquad\text{for every }m,
\]
and therefore
\[
F_{5,5}(m)\le 5\cdot 4^m+4.
\]
\end{theorem}

\begin{proof}
Under the assumed simultaneous gauge fixing, every admissible global object is covered by aligned Hall templates obtained by placing each two-block fiber into the normal form $\mathcal N$ and reading off whether the corresponding block is a rectangle or a strip.
The exact counting inside each aligned template is given by Theorem~\ref{thm:555-hall-template-tensorization}.

In particular, a rectangle block contributes a factor $16=4^2$, while a strip block contributes a factor $4$ together with one residual one-dimensional $(5,5)$ coordinate.
After gauge fixing the aligned part therefore contributes exactly
\[
16^r4^s=4^{2r+s}
\]
when the template has $r$ rectangle blocks and $s$ strip blocks.
Since the total number of scalar coordinates represented by such a template is $2r+s=m$, one gets the bound
\[
\Gamma_5^\times(m)\le 4^m.
\]
The reduction from the ordinary $(5,5)$ problem to the five-partite cross parameter was already recorded earlier in Proposition~\ref{prop:gamma555-reduction}, which yields
\[
F_{5,5}(m)\le 5\,\Gamma_5^\times(m)+4.
\]
Substituting the estimate $\Gamma_5^\times(m)\le 4^m$ proves the claim.
\end{proof}

\begin{remark}[What is now exact and what is still missing in the Hall lane]\label{rem:555-hall-template-status}
Theorem~\ref{thm:555-hall-template-tensorization} is fully rigorous and shows that once the Hall defects are aligned coordinatewise there is no further nonlinear gain left to extract:
the count factors exactly into a fiber term $16^{r(T)}4^{s(T)}$ times the residual one-dimensional $(5,5)$ problem.
So the remaining Hall difficulty is genuinely global.
It is either the exact gauge-aware globalization of Theorem~\ref{thm:555-gauge-aware-globalization}, or any weaker sparse-defect variant strong enough to yield a subexponential cover by aligned Hall templates.
\end{remark}




\subsection*{Prescribed-gauge hexagons and apolar rigidity}
The gauge-aware Hall route admits a sharp local classification. The only irreducible projective obstruction is a prescribed-gauge hexagon, and on the quadratic moment curve that obstruction collapses to a single scalar holonomy.

\begin{definition}[Prescribed-gauge hexagon and its holonomy]\label{def:555-hexagon}
Let $K$ be a field and let
\[
H=(h_{ij})_{i\neq j,\ i,j\in\{1,2,3\}}
\in (K^\times)^6
\]
be the off-diagonal edge data of a $3\times 3$ branch patch. A gauge change rescales rows and columns by
\[
h_{ij}\longmapsto h'_{ij}=a_i h_{ij} b_j^{-1}
\qquad (a_i,b_j\in K^\times).
\]
The \emph{holonomy} of $H$ is
\[
\Xi(H):=\frac{h_{12}h_{23}h_{31}}{h_{21}h_{32}h_{13}}.
\]
\end{definition}

\begin{proposition}[Holonomy classification of prescribed-gauge hexagons]\label{prop:555-hexagon-holonomy}
Let $H$ be a prescribed-gauge hexagon in the sense of Definition~\ref{def:555-hexagon}. Then:
\begin{enumerate}[label=(\roman*),leftmargin=2.2em]
\item $\Xi(H)$ is gauge-invariant;
\item the following are equivalent:
\begin{enumerate}[label=(\alph*),leftmargin=2.2em]
\item $\Xi(H)=1$;
\item there exist nonzero row and column potentials $\rho_i,\kappa_j$ such that
\[
h_{ij}=\rho_i\kappa_j^{-1}
\qquad (i\neq j);
\]
\item $H$ is gauge-equivalent to the unit hexagon $h_{ij}\equiv 1$;
\end{enumerate}
\item every gauge class has a canonical representative of the form
\[
T_\lambda:\qquad
h_{12}=h_{21}=h_{23}=h_{32}=h_{31}=1,
\qquad
h_{13}=\lambda,
\]
and then
\[
\Xi(T_\lambda)=\lambda^{-1}.
\]
\end{enumerate}
Thus the minimal genuine gauge obstruction is one-parameter.
\end{proposition}

\begin{proof}
Gauge invariance is immediate:
\[
\frac{(a_1h_{12}b_2^{-1})(a_2h_{23}b_3^{-1})(a_3h_{31}b_1^{-1})}
{(a_2h_{21}b_1^{-1})(a_3h_{32}b_2^{-1})(a_1h_{13}b_3^{-1})}
=
\Xi(H).
\]

Now choose
\[
a_3=1,\qquad
b_1=h_{31},\qquad
b_2=h_{32},\qquad
a_2=\frac{h_{31}}{h_{21}},\qquad
b_3=\frac{h_{31}h_{23}}{h_{21}},\qquad
a_1=\frac{h_{32}}{h_{12}}.
\]
The transformed edge data satisfy
\[
h'_{12}=h'_{23}=h'_{31}=h'_{21}=h'_{32}=1,
\]
while
\[
h'_{13}
=
\frac{a_1h_{13}}{b_3}
=
\frac{h_{32}h_{21}h_{13}}{h_{12}h_{31}h_{23}}
=
\Xi(H)^{-1}.
\]
Hence every gauge class has the displayed canonical representative. In particular, $\Xi(H)=1$ if and only if $H$ gauges to the unit hexagon. Pulling the gauge back from the unit hexagon gives
\[
h_{ij}=a_i^{-1}b_j,
\]
which is the required row/column factorization.
\end{proof}

\begin{theorem}[Prescribed-gauge hexagon criterion]\label{thm:555-hexagon-criterion}
Assume the projective Hall-profile category has the following $2\times 2$ local positivity property: every admissible lift of a single $2\times 2$ Hall face carries nonnegative projective defect. Then balanced Hall-flow exclusion is equivalent to the following local extension statement:

\begin{quote}
For every ordered triple of distinct rows $(r_1,r_2,r_3)$, every ordered triple of distinct columns $(c_1,c_2,c_3)$, and every admissible lift of the first face
\[
Q(r_1,r_2;c_1,c_2),
\]
the lift extends across the adjacent face
\[
Q(r_1,r_3;c_2,c_3)
\]
with the same projective state on the shared anchor edge $(r_1,c_2)$.
\end{quote}
\end{theorem}

\begin{proof}
If balanced Hall-flow exclusion already holds, then the displayed local extension property is automatic, because any failure of extension produces a balanced projective defect on the $6$-cycle
\[
(r_1,c_1,r_2,c_2,r_3,c_3,r_1).
\]

Conversely, assume the local extension property. Let $\mathfrak C$ be an extreme projective defect circuit, and let its scalar shadow be the balanced table
\[
x=(x_{ij})_{i,j\in[5]}
\]
obtained by forgetting the projective anchor states. Standard transportation-lattice theory implies that the support of $x$ is a simple alternating cycle
\[
C=(r_1,c_1,r_2,c_2,\dots,r_\ell,c_\ell,r_1),
\]
with shadow vector
\[
\chi(C)=\sum_{t=1}^{\ell}\bigl(E_{r_t c_t}-E_{r_{t+1}c_t}\bigr),
\qquad r_{\ell+1}=r_1.
\]
Anchoring at $r_1$ gives the exact decomposition
\[
\chi(C)=Q_1+\cdots+Q_{\ell-1},
\]
where
\[
Q_t=
E_{r_1c_t}-E_{r_{t+1}c_t}+E_{r_{t+1}c_{t+1}}-E_{r_1c_{t+1}}.
\]

By the assumed $2\times 2$ positivity theorem, every $Q_t$ admits an admissible lift of nonnegative defect. Start with any admissible lift of $Q_1$. Inductively, suppose aligned lifts of $Q_1,\dots,Q_t$ have been chosen. The current lift determines a projective state on the shared anchor $(r_1,c_{t+1})$. By the local extension statement, the lift extends across $Q_{t+1}$ while preserving that anchor state. Continuing inductively yields an aligned lift of the entire shadow cycle.

Because adjacent lifts agree on their common anchor state, the internal anchors cancel projectively, so the total defect of the lifted cycle is the sum of the defects of the individual face lifts. Each summand is nonnegative. Therefore the whole circuit has nonnegative defect. Since this holds for every extreme projective defect circuit, balanced Hall-flow exclusion follows.
\end{proof}

\begin{proposition}[Factor-line vanishing of hexagon holonomy]\label{prop:555-factor-line}
Assume that every actual branch-realized $3\times 3$ patch on the quadratic moment curve comes with row factor lines $L_i$, column factor lines $M_j$, and quadratic edge lines
\[
Q_{ij}=L_iM_j
\qquad (i\neq j),
\]
and assume that the prescribed Hall gauge on each edge is the scalar induced from the tensor-product gauge on
\[
L_i\otimes M_j \xrightarrow{\sim} Q_{ij}.
\]
Then every actual prescribed-gauge hexagon has trivial holonomy:
\[
\Xi(H)=1.
\]
\end{proposition}

\begin{proof}
Choose nonzero generators
\[
\ell_i=\alpha_i \ell_i^\circ\in L_i,
\qquad
m_j=\beta_j m_j^\circ\in M_j,
\]
where $\ell_i^\circ$ and $m_j^\circ$ are any fixed reference generators. Under the tensor-product identification,
\[
\ell_i^\circ m_j^\circ=(\alpha_i\beta_j)^{-1}\,\ell_i m_j.
\]
Therefore the prescribed Hall scalar on the edge $(i,j)$ is
\[
h_{ij}=(\alpha_i\beta_j)^{-1}.
\]
Substituting into the holonomy gives
\[
\Xi(H)
=
\frac{(\alpha_1\beta_2)^{-1}(\alpha_2\beta_3)^{-1}(\alpha_3\beta_1)^{-1}}
{(\alpha_2\beta_1)^{-1}(\alpha_3\beta_2)^{-1}(\alpha_1\beta_3)^{-1}}
=1.
\]
So under the factor-line/tensor-product interpretation, every actual branch-realized hexagon is gauge-trivial.
\end{proof}

\begin{proposition}[Explicit apolar pairing on the quadratic moment curve]\label{prop:555-apolar}
Let $V\cong K^2$ and identify the quadratic moment curve with the pure squares
\[
q_t:=(X-tY)^2\in \operatorname{Sym}^2(V).
\]
Define the bilinear form
\[
\langle aX^2+bXY+cY^2,\ a'X^2+b'XY+c'Y^2\rangle
:=
ac'-\frac12 bb'+ca'.
\]
Then
\[
\langle q_s,q_t\rangle=(s-t)^2.
\]
Equivalently, under the coordinate model $x(t)=(1,t,t^2)$ one has
\[
x(s)^\top
\begin{pmatrix}
0&0&1\\
0&-2&0\\
1&0&0
\end{pmatrix}
x(t)
=
(s-t)^2.
\]
\end{proposition}

\begin{proof}
Expand
\[
q_t=(X-tY)^2=X^2-2t\,XY+t^2Y^2.
\]
Substituting
\[
(a,b,c)=(1,-2s,s^2),
\qquad
(a',b',c')=(1,-2t,t^2)
\]
into the definition of $\langle\cdot,\cdot\rangle$ gives
\[
\langle q_s,q_t\rangle
=
1\cdot t^2-\frac12(-2s)(-2t)+s^2\cdot 1
=
t^2-2st+s^2
=
(s-t)^2.
\]
The matrix identity is the same calculation in coordinates.
\end{proof}

\begin{corollary}[Apolar rigidity closes the Hall lane]\label{cor:555-apolar-rigidity}
Assume that for every actual branch-realized pair in the exact quadratic lane, the prescribed-gauge edge class is projectively equal to the apolar value $(t_i-t_j)^2$. Then every actual prescribed-gauge hexagon has trivial holonomy, and consequently
\[
\Gamma_5^\times(m)\le 4^m
\qquad\text{and}\qquad
F_{5,5}(m)\le 5\cdot 4^m+4.
\]
\end{corollary}

\begin{proof}
Under the stated rigidity, every edge gauge has the form
\[
h_{ij}=r_i c_j (t_i-t_j)^2
\]
for suitable nonzero row and column factors $r_i,c_j$. Therefore
\[
\Xi(H)
=
\frac{(t_1-t_2)^2(t_2-t_3)^2(t_3-t_1)^2}
{(t_2-t_1)^2(t_3-t_2)^2(t_1-t_3)^2}
=1.
\]
So every actual prescribed-gauge hexagon is gauge-trivial. Theorem~\ref{thm:555-gauge-aware-globalization} then gives $\Gamma_5^\times(m)\le 4^m$, and the reduction in Proposition~\ref{prop:gamma555-reduction} yields
\[
F_{5,5}(m)\le 5\cdot 4^m+4.
\]
\end{proof}

\begin{remark}[The remaining local Hall gap]\label{rem:555-hexagon-gap}
The Hall route is now localized to one scalar invariant. Propositions~\ref{prop:555-hexagon-holonomy} and \ref{prop:555-factor-line} show that every prescribed-gauge hexagon is classified by its holonomy $\Xi(H)$ and that any factor-line realization kills that holonomy automatically. Corollary~\ref{cor:555-apolar-rigidity} shows that, on the quadratic moment curve, it would even be enough to prove the two-branch apolar rigidity statement. Thus the first genuine unresolved Hall theorem is either:
\begin{enumerate}[label=(\roman*),leftmargin=2.2em]
\item the prescribed-gauge hexagon extension theorem of Theorem~\ref{thm:555-hexagon-criterion}, or
\item the stronger two-branch apolar rigidity statement used in Corollary~\ref{cor:555-apolar-rigidity}.
\end{enumerate}
\end{remark}

\status{The Hall lane now has an exact local obstruction theory. Up to gauge, every primitive $3\times 3$ defect is a single twisted hexagon $T_\lambda$. What remains open is not a higher-dimensional globalization mystery, but the elimination of that one scalar holonomy in actual branch-realized hexagons.}


\begin{proposition}[Exact local scalar diagonalizer for the $(4,4)$ model]\label{prop:444-local-diagonalizer}
Identify $[4]$ with $G=\mathbb F_2^2$. Define
\[
T(a,b,c,d)
:=\mathbf 1_{a+b+c+d=0}
-\mathbf 1_{a=b}\mathbf 1_{c=d}
-\mathbf 1_{a=c}\mathbf 1_{b=d}
-\mathbf 1_{a=d}\mathbf 1_{b=c}
+3\mathbf 1_{a=b=c=d}.
\]
Then $T(a,b,c,d)=1$ exactly when the $4$-tuple $(a,b,c,d)$ is constant or injective, and $T(a,b,c,d)=0$ otherwise. Consequently, if
\[
T_m(x,y,z,w):=\prod_{j=1}^m T(x_j,y_j,z_j,w_j),
\]
and $C\subseteq G^m$ is $4$-sunflower-free, then
\[
T_m(x,y,z,w)=\mathbf 1_{x=y=z=w}
\qquad\text{for all }x,y,z,w\in C.
\]
\end{proposition}

\begin{proof}
For a $4$-tuple in $G^4$ there are only five collision types: $4$, $3+1$, $2+2$, $2+1+1$, and $1+1+1+1$. If all four entries are equal, then each of the three pairing indicators is $1$, the diagonal indicator is $1$, and $a+b+c+d=0$, so $T=1-1-1-1+3=1$. If the tuple has type $2+2$, then $a+b+c+d=0$, exactly one of the three pairings occurs, and the diagonal term vanishes, so $T=1-1=0$. If the tuple has type $3+1$ or $2+1+1$, then $a+b+c+d\neq 0$ in $\mathbb F_2^2$, and the correction terms do not change this, so $T=0$. If the tuple is injective, then its entries are exactly the four elements of $G$, whose sum is $0$; none of the pairing terms occurs, so $T=1$. This proves the local support statement.

Now let $C\subseteq G^m$ be $4$-sunflower-free. If $x=y=z=w$, then every coordinate is constant and $T_m(x,y,z,w)=1$. If $x,y,z,w$ are distinct and $T_m(x,y,z,w)\neq 0$, then every coordinate is constant or injective, so the four codewords form a forbidden sunflower. Finally, if not all four codewords are equal but some repetition occurs, choose two distinct codewords among them and a coordinate where they differ. At that coordinate the $4$-tuple is neither constant nor injective, because a repeated codeword forces a repeated value. Hence that coordinate contributes a factor $0$, so again $T_m(x,y,z,w)=0$. Therefore $T_m$ restricts to the diagonal tensor on $C^4$.
\end{proof}

\leannote{This local support and global diagonality statement is formalized and proved in \texttt{FormalConjectures/Problems/Erdos/E20/Tensors/CharacterTensors.lean}, including \texttt{localCharacterTensor\_eq\_indicator}, \texttt{explicitLocalTensor\_eq\_indicatorF2}, \texttt{globalTensorTuple\_eq\_ite\_allEqual\_of\_sunflowerFree}, and \texttt{explicitGlobalTensorTuple\_eq\_ite\_allEqual\_of\_sunflowerFree}. The wrappers \texttt{pasted\_local\_character\_tensor\_exact\_support}, \texttt{pasted\_explicit\_local\_tensor\_exact\_support}, \texttt{pasted\_global\_tensor\_diagonal\_on\_sunflower\_free}, and \texttt{pasted\_explicit\_global\_tensor\_diagonal\_on\_sunflower\_free} appear in \texttt{InformalMap.lean}.}

\begin{remark}[The proposed $(16/(3\sqrt 3))^m$ bound is still not certified]\label{rem:444-bound-not-certified}
Proposition~\ref{prop:444-local-diagonalizer} gives a completely rigorous local diagonalizer and exact restriction-to-diagonal theorem. What it does \emph{not} yet provide is a certified global slice-rank bound of the form
\[
F_{4,4}(m)\le \left(\frac{16}{3\sqrt 3}\right)^m.
\]
In particular, the specific count
\[
\operatorname{slice-rank}(T_m)\le \sum_{j\le m/2}\binom{2m}{j}
\]
cannot be correct as stated: at $m=1$ the right-hand side equals $1$, whereas already $F_{4,4}(1)=3$ and the restriction of $T_1$ to any $3$-point sunflower-free subset of $G$ is the $3$-point diagonal tensor, whose slice rank is $3$. So the exact local object is rigorous, but the global rank/counting step still needs a different argument.
\end{remark}

\begin{proposition}[Honest global rank repair for the $(4,4)$ diagonalizer]\label{prop:444-honest-rank}
In the notation of Proposition~\ref{prop:444-local-diagonalizer}, the matrix flattening of $T_m$ with $x$ as the row variable and $(y,z,w)$ as the column variable has rank exactly $4^m$. Consequently
\[
\operatorname{slice-rank}(T_m)\le 4^m.
\]
Hence every $4$-sunflower-free set $C\subseteq G^m$ satisfies
\[
|C|\le 4^m.
\]
Equivalently,
\[
F_{4,4}(m)\le 4^m.
\]
\end{proposition}

\begin{proof}
For each $x\in G^m$, consider the column indexed by $(x,x,x)$. If $x'\in G^m$, then
\[
T_m(x',x,x,x)=\prod_{j=1}^m T(x'_j,x_j,x_j,x_j).
\]
By the local description of $T$, one has $T(a,b,b,b)=1$ if and only if $a=b$: if $a=b$, the $4$-tuple $(a,b,b,b)$ is constant, while if $a\neq b$, then $(a,b,b,b)$ is neither constant nor injective. Therefore
\[
T_m(x',x,x,x)=\prod_{j=1}^m \mathbf 1_{x'_j=x_j}=\mathbf 1_{x'=x}.
\]
So the $x\mid (y,z,w)$ flattening contains an identity submatrix of size $4^m$, indexed by all rows $x'$ and by the columns $(x,x,x)$. Since the flattening has exactly $4^m$ rows, its rank is exactly $4^m$. The slice rank of a tensor is at most the matrix rank of any flattening, so $\operatorname{slice-rank}(T_m)\le 4^m$.

Now let $C\subseteq G^m$ be $4$-sunflower-free. By Proposition~\ref{prop:444-local-diagonalizer},
\[
T_m\big|_{C^4}=\mathbf 1_{\{x=y=z=w\}}.
\]
The diagonal tensor on $C^4$ has slice rank $|C|$, hence
\[
|C|\le \operatorname{slice-rank}(T_m)\le 4^m.
\]
This is exactly the claimed bound for $F_{4,4}(m)$.
\end{proof}

\begin{proposition}[Balanced local support of the $(4,4)$ tensor]\label{prop:444-balanced}
Let
\[
S:=\{(a,b,c,d)\in G^4:T(a,b,c,d)=1\}.
\]
Then $S$ is exactly the union of the $4$ constant tuples and the $24$ injective tuples, so $|S|=28$. Moreover each one-coordinate marginal of the uniform measure on $S$ is uniform on $G$. Equivalently, if $\lambda_1,\lambda_2,\lambda_3,\lambda_4:G\to \R$ each have mean zero, then
\[
\sum_{(a,b,c,d)\in S}\bigl(\lambda_1(a)+\lambda_2(b)+\lambda_3(c)+\lambda_4(d)\bigr)=0.
\]
\end{proposition}

\begin{proof}
The description of $S$ follows from Proposition~\ref{prop:444-local-diagonalizer}. Fix a coordinate, say the first. Each symbol $u\in G$ appears exactly once in that coordinate among the constant tuples, and exactly six times among the injective tuples, because after fixing $a=u$ the remaining entries must be the three elements of $G\setminus\{u\}$ in arbitrary order. Thus each symbol appears $7$ times among the $28$ members of $S$, so the first-coordinate marginal is uniform. By symmetry the same is true in every coordinate.

For the displayed identity, use the uniform marginals coordinate by coordinate. One has
\[
\sum_{(a,b,c,d)\in S}\lambda_1(a)=7\sum_{u\in G}\lambda_1(u)=0,
\]
and similarly for $\lambda_2,\lambda_3$, and $\lambda_4$. Summing the four equalities gives the claim.
\end{proof}

\begin{remark}[Representation-theoretic interpretation]\label{rem:444-polystable}
For each fixed first coordinate $a\in G$, the corresponding row of the mode-$1$ flattening of $T$ is supported on the single constant triple $(a,a,a)$ and on the six permutations of $G\setminus\{a\}$. Hence the four rows are pairwise orthogonal and each has squared norm $7$; the same statement holds in every mode. Thus every one-body reduced density matrix is the scalar matrix $7I_4$. In the usual Kempf--Ness language, the local tensor is already balanced/polystable. This is the representation-theoretic form of Proposition~\ref{prop:444-balanced}: the previous lower-tail instability count had no genuine local defect to exploit.
\end{remark}

\status{The current exact global output of the $(4,4)$ diagonalizer is therefore only the ambient base $4$. Any genuine base below $4$ must come from a different global rank count, not from the previous lower-tail generating-function step.}

\subsection{Exploratory local tensor program (quarantined)}\label{subsec:exploratory-tensors}
We now record the exact local constructions that were found for several diagonal fixed-alphabet models. The local support and global diagonalization parts are rigorous. The global rank counts, however, currently contain proof gaps in several cases. We record them here for completeness but do not use them elsewhere.

\begin{remark}[The recurring proof gap]
In multiple slice-rank notes, the following invalid maneuver appears in one form or another:
\[
\sum_{j\le L} a_j \le x^{-L}\sum_j a_j x^j
\quad\text{with }x>1.
\]
This is an \emph{upper-tail} estimate being misapplied to a \emph{lower-tail} count. Therefore the current numerical global bounds for $(4,4)$, $(5,5)$, $(6,6)$, and $(8,8)$ obtained by that route should be treated as provisional until the counting step is repaired.
\end{remark}

\paragraph{(4,4) local tensor.}
Proposition~\ref{prop:444-local-diagonalizer} gives an exact local scalar diagonalizer and exact restriction-to-diagonal theorem for sunflower-free codes in $[4]^m\cong (\mathbb F_2^2)^m$. What remains missing is a correct global rank count. In particular, Remark~\ref{rem:444-bound-not-certified} explains why the proposed bound $F_{4,4}(m)\le (16/(3\sqrt 3))^m$ is still not certified.

\paragraph{(5,5) local tensor / moment-curve reduction.}
Theorem~\ref{thm:5column-two-moment} and Corollary~\ref{cor:555-moment-curve} show that the local support is already cut out exactly by the two equations $\sum a_i=0$ and $\sum a_i^2=0$, giving an exact reduction to diagonal zero-sum problems on the quadratic moment curve $P=\{(t,t^2):t\in \mathbb F_5\}$. Propositions~\ref{prop:555-unit-tensor}, \ref{prop:555-raw-tensor-obstruction}, and \ref{prop:555-local-monomial-obstruction} now make rigorous that the most obvious coordinatewise tensor counts stall at base $5$, while Theorem~\ref{thm:555-two-coordinate} shows that block length $2$ already beats the critical $25/6$ threshold. What remains missing is a genuine tensorization / globalization theorem for this reduced model. Earlier numerical bases such as $(19/4)^m$ or $(3.37118\ldots)^m$ remain provisional.

\paragraph{(6,6) local tensor.}
A scalar Fourier-twisted diagonal selector for constant columns plus a Vandermonde antisymmetrizer gives exact support on constant-or-injective 6-columns. The current global base $<6$ again relies on an invalid lower-tail generating-function estimate and remains provisional.

\paragraph{Power-of-two diagonalizer.}
For $q=2^d$, there is a remarkably clean matrix-valued local diagonalizer based on $-I$ on constant columns and a trace-zero square-zero nilpotent $N$ on injective columns. This exact local object really diagonalizes every sunflower-free code. Counterexample~\ref{ce:444-power-of-two-false} shows, however, that the naive scalar bound $F_{2^d,2^d}(m)\le (2^d-1)^m$ is already false at $d=2$. What is missing is a rigorous global rank estimate for the resulting cyclic tensor network. A Boolean-polynomial substitute for the case $q=8$ also currently has the same lower-tail counting issue.

\status{The exploratory local-tensor program remains the most promising hard-model lane in the project, but at present only the exact local support and exact global diagonalization steps are fully secure. The final numerical exponential constants remain unverified until the global rank step is repaired.}

\section{Counterexamples to false routes}\label{sec:counterexamples}
This section collects the strongest rigorous counterexamples discovered in the project.

\subsection{Broad support-piece dichotomies are fake}
Proposition~\ref{prop:one-block-universality} shows that one-block empty-core support pieces can encode arbitrary smaller coreless sunflower-free families. Thus any broad support-piece dichotomy of the form ``balanced sparse versus heavy trace'' is not a genuine simplification of the global problem.

\subsection{Low-order projected-branch LP inverse theorems are false}
Counterexample~\ref{ce:parity-projective} shows that parity/code slices are exact dual obstructions for bounded-order LPs while staying globally far from every block-product measure. The same phenomenon occurs with Latin-square slices and more general OA/code slices. Thus no inverse theorem based only on bounded-order inclusion statistics can classify the true obstruction class.

\subsection{Low-$\sigma$ rigidity is false}
Inside the $q$-partite transversal model with $q=k-1$, let
\[
C_{\mathrm{par}}:=\{x\in G^r:x_1+\cdots+x_r=0\}
\]
for an abelian group $G$ of order $q$, and let $H(C_{\mathrm{par}})$ be its transversal family. Then every coordinate marginal is uniform, so
\[
\sigma(H(C_{\mathrm{par}}))=\frac{1}{q}=\frac{1}{k-1},
\]
which is the same value as for the full width-$q$ product. But $C_{\mathrm{par}}$ is a constant edit distance away from every genuine width-$q$ block product of the same size. Thus low $\sigma$ controls only one-coordinate marginals, not global shape.

\subsection{Positive-density microscopic box extraction is false}
The first failure of several predictive-state and dense-box routes already appears inside the exact width-$q$ product, where $q:=k-1$.

\begin{counterexample}[Balanced multislice blocks positive-density box extraction]\label{ce:balanced-multislice}
Fix $k\ge 3$ and write $q:=k-1$. Assume $q\mid n$. On the ground set
\[
U=[n]\times [q]
\]
encode each word $x=(x_1,\dots,x_n)\in [q]^n$ by the $n$-set
\[
A_x:=\{(i,x_i):1\le i\le n\}.
\]
Let
\[
X_n:=\Bigl\{x\in [q]^n:\ |\{i:x_i=a\}|=n/q\text{ for every }a\in [q]\Bigr\},
\qquad
\mathcal B_n:=\{A_x:x\in X_n\}.
\]
Then:
\begin{enumerate}[label=(\roman*),leftmargin=2.2em]
\item $\mathcal B_n$ is $k$-sunflower-free;
\item for all sufficiently large $n$ divisible by $q$, $\mathcal B_n$ is already a terminal kernel under leaf stripping;
\item
\[
|\mathcal B_n|=\frac{n!}{((n/q)!)^q}
=
q^n\cdot \Theta_q\bigl(n^{-(q-1)/2}\bigr);
\]
\item for every support box $B=\prod_{i=1}^n S_i\subseteq [q]^n$, if
\[
r(B):=\bigl|\{i:|S_i|\ge 2\}\bigr|,
\]
then
\[
\frac{|X_n\cap B|}{|B|}\le \frac{C_q}{\sqrt{r(B)}}
\]
for a constant $C_q$ depending only on $q$.
\end{enumerate}
In particular, every support box carrying density at least $\delta>0$ inside $X_n$ has
\[
r(B)\le \frac{C_q^2}{\delta^2}.
\]
\end{counterexample}

\begin{proof}
For (i), take $k=q+1$ distinct members $A_{x^{(1)}},\dots,A_{x^{(k)}}$ of $\mathcal B_n$. In coordinate $i$, the chosen elements are
\[
(i,x_i^{(1)}),\dots,(i,x_i^{(k)}).
\]
For a sunflower, each coordinate column would have to be either constant or pairwise disjoint. But there are only $q=k-1$ possible symbols, so $k$ distinct symbols cannot occur in one coordinate. Therefore every coordinate must be constant, which forces
\[
x^{(1)}=\cdots=x^{(k)},
\]
contrary to distinctness. So $\mathcal B_n$ is $k$-sunflower-free.

For (ii), fix a vertex $(i,a)\in U$. The number of balanced words with $x_i=a$ is
\[
d_{\mathcal B_n}(i,a)=
\frac{(n-1)!}{((n/q)-1)!\,((n/q)!)^{q-1}},
\]
which tends to infinity with $n$ for fixed $q$. Hence for all sufficiently large $n$ every vertex degree exceeds $1$, so leaf stripping removes nothing and $\mathcal B_n$ is already terminal.

For (iii), the exact counting formula is immediate from the definition of $X_n$, and Stirling's formula gives
\[
\frac{n!}{((n/q)!)^q}
=
q^n\cdot \Theta_q\bigl(n^{-(q-1)/2}\bigr).
\]

For (iv), let $L:=\{i:|S_i|\ge 2\}$, so $|L|=r(B)$. For each symbol $a\in [q]$, write
\[
m_a:=|\{i\in L:a\in S_i\}|.
\]
Because every live coordinate contributes at least two symbol incidences,
\[
\sum_{a=1}^q m_a=\sum_{i\in L}|S_i|\ge 2r(B).
\]
Hence some symbol $a$ satisfies
\[
m_a\ge \frac{2r(B)}{q}.
\]
For $s\in \{2,\dots,q\}$, let
\[
I_s:=\{i\in L:\ a\in S_i,\ |S_i|=s\}.
\]
Since the sets $I_s$ partition the $m_a$ live coordinates containing $a$, there exists some $s$ with
\[
M:=|I_s|\ge \frac{m_a}{q-1}\ge \frac{2r(B)}{q(q-1)}.
\]

Under the uniform product measure on $B$, define
\[
Z:=\sum_{i=1}^n \mathbf 1_{\{x_i=a\}}.
\]
Write
\[
Z=Y+W,
\qquad
Y:=\sum_{i\in I_s}\mathbf 1_{\{x_i=a\}},
\]
so $Y\sim \mathrm{Bin}(M,1/s)$ and $W$ is independent of $Y$. For any integer $t$,
\[
\Pr_B(Z=t)=\sum_u \Pr(W=u)\Pr(Y=t-u)\le \max_v \Pr(Y=v).
\]
By the standard Stirling estimate for binomial masses, since $2\le s\le q$ is fixed,
\[
\max_v \Pr(Y=v)\le \frac{C'_q}{\sqrt{M}}\le \frac{C_q}{\sqrt{r(B)}}.
\]
But every $x\in X_n$ has exactly $n/q$ occurrences of the symbol $a$, so
\[
\frac{|X_n\cap B|}{|B|}
=
\Pr_B(x\in X_n)
\le
\Pr_B\!\left(Z=\frac{n}{q}\right)
\le
\frac{C_q}{\sqrt{r(B)}}.
\]
This proves (iv).
\end{proof}

\begin{corollary}[No positive-density extraction from arbitrary large terminal kernels]\label{cor:no-positive-density-extraction}
There is no theorem forcing every sufficiently large terminal kernel of near-neutral size to contain a support box on $\Omega(n)$ live coordinates with density bounded below by a positive constant. The balanced multislice family of Counterexample~\ref{ce:balanced-multislice} already has size
\[
q^n n^{-O_q(1)},
\]
is terminal for all large $n$, and every support box of positive density has only $O_q(1)$ live coordinates.
\end{corollary}

\begin{proof}
This is immediate from parts (ii)--(iv) of Counterexample~\ref{ce:balanced-multislice}.
\end{proof}

\begin{proposition}[Mesoscopic slice products show the sharp density scale]\label{prop:mesoscopic-slice-density}
Fix $q\ge 2$. Let $n=tm$ with $q\mid m$, partition the coordinates into $t$ consecutive blocks of size $m$, and set
\[
Y_{t,m}:=X_m\times \cdots \times X_m\subseteq [q]^n.
\]
Then
\[
|Y_{t,m}|=
q^n\exp\!\Bigl(-\frac{q-1}{2}t\log m+O_q(t)\Bigr).
\]
If $B=\prod_{i=1}^n S_i$ has at least $\alpha n$ live coordinates, then
\[
\frac{|Y_{t,m}\cap B|}{|B|}
\le
\exp\bigl(-\Omega_{\alpha,q}(t\log m)\bigr).
\]
Consequently, whenever $m\to\infty$, $t\to\infty$, and $t\log m=o(n)$, these examples have total density $e^{-o(n)}$ inside the ambient width-$q$ product, but every linear-width support box has density at most $e^{-\Omega_{\alpha,q}(t\log m)}$. Thus the best possible universal microscopic density scale in such extraction theorems is only $e^{-o(n)}$, not positive density.
\end{proposition}

\begin{proof}
The size formula follows by multiplying the estimate from Counterexample~\ref{ce:balanced-multislice}(iii) across the $t$ blocks.

For the density claim, let $r_j$ be the number of live coordinates of $B$ inside block $j$, so
\[
\sum_{j=1}^t r_j\ge \alpha tm.
\]
Let $H$ be the set of blocks with $r_j\ge \alpha m/2$. If $|H|<\alpha t/2$, then
\[
\sum_{j=1}^t r_j
<
\frac{\alpha t}{2}\,m+\left(t-\frac{\alpha t}{2}\right)\frac{\alpha m}{2}
=
\alpha tm\left(1-\frac{\alpha}{4}\right)
<
\alpha tm,
\]
a contradiction. Hence
\[
|H|\ge \frac{\alpha t}{2}.
\]

Write $B=\prod_{j=1}^t B_j$ according to the block decomposition. Since $Y_{t,m}$ is a direct product,
\[
\frac{|Y_{t,m}\cap B|}{|B|}
=
\prod_{j=1}^t \frac{|X_m\cap B_j|}{|B_j|}.
\]
For each heavy block $j\in H$, Counterexample~\ref{ce:balanced-multislice}(iv) gives
\[
\frac{|X_m\cap B_j|}{|B_j|}
\le
\frac{C_q}{\sqrt{\alpha m/2}}
=
C_{\alpha,q}m^{-1/2}.
\]
For the remaining blocks we use only the trivial bound $1$. Therefore
\[
\frac{|Y_{t,m}\cap B|}{|B|}
\le
\bigl(C_{\alpha,q}m^{-1/2}\bigr)^{|H|}
\le
\bigl(C_{\alpha,q}m^{-1/2}\bigr)^{\alpha t/2}
=
\exp\bigl(-\Omega_{\alpha,q}(t\log m)\bigr).
\]
The final conclusion follows by comparing this with the displayed size formula.
\end{proof}

\subsection{The naive A/B bridge is false; the right replacement is a trichotomy}
Consider the full $(k-1)$-ary transversal product
\[
\mathcal T_{n,k-1}:=\{\{v_1,\dots,v_n\}: v_i\in V_i,\ |V_i|=k-1\}.
\]
It is terminal, $k$-sunflower-free, has no profitable prefix block beyond the neutral rate $(k-1)^{-t}$, and contains no exponentially thick $(4,4)$ or $(5,5)$ solved trace piece when $k\ge 7$. Thus the dichotomy ``solved piece or profitable prefix'' is false. The correct large-scale picture is a trichotomy:
\begin{center}
\emph{solved local tensor piece / profitable-prefix piece / neutral obstruction piece.}
\end{center}
The neutral piece is realized already by exact or near $(k-1)$-ary products and their finite-memory analogues.

\section{Recursive closure of solved trace classes}\label{sec:union-classes}
The recursion itself is exact once a solved trace class is identified. This is the right formal closure theorem for all later routing programs.

\begin{definition}[Solved trace class]
A trace class $\mathcal C$ is called solved with base function $B_{\mathcal C}(r,k)$ if, whenever $G$ is $k$-sunflower-free, $P$ is a support piece of rank $r$, and $\Tr_P(G)\in \mathcal C$, one has
\[
|\Tr_P(G)|\le B_{\mathcal C}(r,k).
\]
\end{definition}

\begin{theorem}[Exact one-step recursion over a solved trace class]\label{thm:unionclass-closure}
If $G$ has residual rank $m$ and $P$ is a support piece of rank $r$ with $\Tr_P(G)\in \mathcal C$, then
\[
|G|\le B_{\mathcal C}(r,k)\,h(m-r,k),
\]
where $h(m,k)$ is the extremal size among reduced kernel links of residual rank $m$.
Consequently, if
\[
C_{\mathcal C}^*(k):=\sup_{r\ge 1} B_{\mathcal C}(r,k)^{1/r},
\]
then
\[
|G|\le \bigl(C_{\mathcal C}^*(k)\bigr)^r h(m-r,k).
\]
Moreover, if a reduced kernel link is recursively decomposable through a union $\mathcal U$ of solved trace classes, then its size is at most
\[
\bigl(C_{\mathcal U}^*(k)\bigr)^m,
\qquad
C_{\mathcal U}^*(k):=\sup_{r\ge 1} B_{\mathcal U}(r,k)^{1/r}.
\]
\end{theorem}

\begin{proof}
The one-step inequality is the exact trace recursion together with the trace-class bound. The recursive closure follows by induction on the residual rank.
\end{proof}

\leannote{This recursive-closure argument is formalized and proved in \texttt{FormalConjectures/Problems/Erdos/E20/Compendium/Recurrences/RecursiveClosure.lean} as \texttt{one\_step\_recursion\_over\_solved\_class} and \texttt{recursive\_closure\_bound}.}

\begin{remark}[What can be plugged in rigorously]
Without using any provisional local-tensor constants, the theorem applies immediately to exact block products (base $k-1$), positive-density support boxes (base $\lfloor (k-1)/\delta\rfloor$), and fixed-memory finite-state classes (base equal to their transfer root or entropy-gap bound). If one later proves rigorous fixed-$q$ local-tensor bounds, they plug in seamlessly here.
\end{remark}

\section{Conditional and programmatic statements}\label{sec:conditional}
This section records the most useful exact reductions that still depend on a missing structural theorem.

\begin{conjecture}[Universal routing/profitable-prefix theorem for terminal kernels]
For each fixed $k$, every sufficiently large terminal kernel either already lies in one of the theorem-producing model classes of Sections~\ref{sec:block-products} and \ref{sec:finite-state}, or else has a support piece / reduced link with a profitable prefix block whose descendants recursively do.
\end{conjecture}

\begin{remark}
Together with Theorem~\ref{thm:adaptive-kernel-prefix} and Theorem~\ref{thm:unionclass-closure}, this conjecture would immediately imply the full exponential sunflower bound.
\end{remark}

\subsection*{Forced trichotomy and neutral terminals}
In the routing statements below, we work inside a reduced-link recursion domain that is closed under taking nonempty links and under passing to terminal kernels of traces and projected branches. We write $h(m,k)$ for the corresponding extremal function in residual rank $m$.

\begin{definition}[Solved closure, profitable class, and neutral terminals]\label{def:neutral-terminal}
Fix $k$ and $\lambda>1$. Let $\mathcal S_k$ be any class of reduced kernel links such that:
\begin{enumerate}[label=(\roman*),leftmargin=2.2em]
\item $\mathcal S_k$ contains the theorem-producing model classes from Sections~\ref{sec:block-products} and \ref{sec:finite-state}, together with the bounded-degree reduced kernel links from Section~\ref{sec:correlation};
\item $\mathcal S_k$ is closed under exact support-piece recursion: whenever $H$ has a proper support piece $P$ such that $K(\Tr_P(H))\in \mathcal S_k$ and $K(H_\tau)\in \mathcal S_k$ for every $\tau\in \Tr_P(H)$, one has $H\in \mathcal S_k$.
\end{enumerate}
Let $\mathcal P_\lambda$ be the class of reduced kernel links $H$ for which
\[
\prod_{i=0}^{t-1}\beta_i(H)\ge \lambda^{-t}
\qquad\text{for some }1\le t\le \rank(H).
\]
For a proper support piece $P$ of $H$, define its child set by
\[
\mathcal C_P(H):=\{K(\Tr_P(H))\}\cup \{K(H_\tau):\tau\in \Tr_P(H)\}.
\]
A reduced kernel link $T$ is called \emph{$(k,\lambda)$-terminal} if $T\notin \mathcal S_k$ and
\[
\mathcal C_P(T)\subseteq \mathcal S_k\cup \mathcal P_\lambda
\qquad\text{for every proper support piece }P.
\]
If moreover $T\notin \mathcal P_\lambda$, we call $T$ a \emph{neutral terminal}.
\end{definition}

\begin{theorem}[Forced residual trichotomy]\label{thm:forced-trichotomy}
For every nontrivial reduced kernel link $H$, at least one of the following holds:
\begin{enumerate}[label=(\roman*),leftmargin=2.2em]
\item $H\in \mathcal S_k$;
\item $H\in \mathcal P_\lambda$;
\item $H$ reaches a neutral terminal kernel after finitely many support-piece reductions.
\end{enumerate}
Equivalently, every reduced kernel link outside $\mathcal S_k\cup \mathcal P_\lambda$ has a neutral terminal descendant.
\end{theorem}

\begin{proof}
Assume $H\notin \mathcal S_k\cup \mathcal P_\lambda$. Consider all descendants of $H$ obtainable by repeatedly choosing a proper support piece $P$ and passing to one of the children in $\mathcal C_P(H)$. Each step strictly decreases the residual rank, so the search tree is finite.

Choose, among all descendants lying outside $\mathcal S_k\cup \mathcal P_\lambda$, one of minimal residual rank and call it $N$. Then $N\notin \mathcal S_k\cup \mathcal P_\lambda$, and every child of every proper support piece of $N$ has smaller residual rank. By the minimality of $N$, no such child can remain outside $\mathcal S_k\cup \mathcal P_\lambda$. Thus
\[
\mathcal C_P(N)\subseteq \mathcal S_k\cup \mathcal P_\lambda
\qquad\text{for every proper support piece }P,
\]
so $N$ is $(k,\lambda)$-terminal. Since $N\notin \mathcal P_\lambda$, it is neutral terminal.
\end{proof}

\begin{theorem}[Neutral terminals force a profitable child]\label{thm:neutral-terminal-child}
Let $T$ be a neutral terminal kernel. Then for every proper support piece $P$ of $T$,
\[
\mathcal C_P(T)\cap \mathcal P_\lambda\neq \varnothing.
\]
In words: every proper support split of a neutral terminal immediately produces a profitable smaller child.
\end{theorem}

\begin{proof}
Fix a proper support piece $P$ of $T$. Since $T$ is terminal,
\[
\mathcal C_P(T)\subseteq \mathcal S_k\cup \mathcal P_\lambda.
\]
Suppose for contradiction that $\mathcal C_P(T)\cap \mathcal P_\lambda=\varnothing$. Then every child in $\mathcal C_P(T)$ belongs to $\mathcal S_k$. By the recursive closure property in Definition~\ref{def:neutral-terminal}, this implies $T\in \mathcal S_k$, contradicting that $T$ is neutral terminal. Therefore $\mathcal C_P(T)\cap \mathcal P_\lambda\neq \varnothing$.
\end{proof}

\begin{proposition}[Neutral terminals are the smallest forced residual class]\label{prop:neutral-terminal-minimal}
Let $\mathcal R$ be any class of reduced kernel links such that every nontrivial reduced kernel link belongs to
\[
\mathcal S_k\cup \mathcal P_\lambda\cup \mathcal R.
\]
Then every neutral terminal kernel belongs to $\mathcal R$.
\end{proposition}

\begin{proof}
If $N$ is neutral terminal, then by definition $N\notin \mathcal S_k\cup \mathcal P_\lambda$. Hence any universal trichotomy of the displayed form must place $N$ in $\mathcal R$.
\end{proof}

\begin{corollary}[If neutral terminals do not exist, the routing step closes]\label{cor:no-neutral-terminals}
Fix $k$, and suppose there is $\lambda_k>1$ such that every member of $\mathcal S_k$ of residual rank $m$ has size at most $\lambda_k^m$. If there are no neutral terminal reduced kernel links for $\lambda_k$, then
\[
h(m,k)\le \lambda_k^m
\qquad\text{for all }m.
\]
Consequently,
\[
M(n,k)\le ((k-1)\lambda_k)^n,
\qquad
f(n,k)<((k-1)\lambda_k+1)^n.
\]
\end{corollary}

\begin{proof}
We prove $h(m,k)\le \lambda_k^m$ by induction on $m$. Let $H$ be a reduced kernel link of residual rank $m$. If $H\in \mathcal S_k$, the claimed bound is part of the hypothesis. Otherwise Theorem~\ref{thm:forced-trichotomy} and the nonexistence of neutral terminals imply that $H\in \mathcal P_{\lambda_k}$. Hence for some $t\in \{1,\dots,m\}$ one has
\[
\prod_{i=0}^{t-1}\beta_i(H)\ge \lambda_k^{-t}.
\]
By Theorem~\ref{thm:heavy-set}, some $t$-link $H_T$ satisfies $|H_T|\ge |H|\lambda_k^{-t}$. Since $H_T$ has smaller residual rank, induction gives $|H_T|\le h(m-t,k)\le \lambda_k^{m-t}$. Therefore $|H|\le \lambda_k^t|H_T|\le \lambda_k^m$. The bound for ordinary sunflower-free families now follows from Theorem~\ref{thm:exact-reduction}.
\end{proof}

\status{The trichotomy itself is formal and rigorous. The unresolved issue is not whether a neutral terminal class exists as a logical possibility, but whether every neutral terminal can be eliminated or lifted to a genuine profitable prefix.}

\subsection*{Lift coherence for profitable child witnesses}
Theorem~\ref{thm:neutral-terminal-child} is only a pointwise child-profit statement. The remaining issue is whether the profitable children produced by different support pieces align to a \emph{single} profitable prefix upstairs.

\begin{definition}[Child-profit lift hypergraph]\label{def:child-profit-lift}
Let $T$ be a terminal kernel in the routing domain. Write $\Pi(T)$ for the set of admissible prefix blocks on $T$. For a proper support piece $P$ of $T$, define the child family
\[
\mathcal C_T(P):=\{K(\Tr_P(T))\}\cup \{K(T_\tau):\tau\in \Tr_P(T)\}.
\]
Define the corresponding \emph{lift set}
\[
\Lambda_T(P):=\{\rho\in \Pi(T): \rho \text{ lifts a profitable prefix from some child in }\mathcal C_T(P)\cap \mathcal P_\lambda\}.
\]
For $\rho\in \Pi(T)$, let
\[
\Gamma_T(\rho):=\{P: P \text{ is a proper support piece of }T \text{ and } \rho\in \Lambda_T(P)\}.
\]
The incidence system
\[
H_T:=\bigl(\Pi(T),\{\Lambda_T(P): P \text{ a proper support piece of }T\}\bigr)
\]
is called the \emph{child-profit lift hypergraph} of $T$.
\end{definition}

\begin{definition}[$\beta$-coherent prefix]\label{def:beta-coherent-prefix}
Let $T$ be a terminal kernel. A prefix $\rho\in \Pi(T)$ is called \emph{$\beta$-coherent} if the profitable child witnesses lifted through $\rho$ satisfy the profitable-prefix hypothesis from Theorem~\ref{thm:adaptive-kernel-prefix} for $T$. Equivalently, the family of support pieces in $\Gamma_T(\rho)$ is large enough to certify that $T\in \mathcal P_\lambda$.
\end{definition}

\begin{proposition}[Coherent-lift reformulation of the missing terminal step]\label{prop:coherent-lift}
Let $T$ be a terminal kernel. Then the following are equivalent:
\begin{enumerate}[label=(\roman*),leftmargin=2.2em]
\item $T$ has a profitable prefix, equivalently $T\in \mathcal P_\lambda$;
\item some prefix $\rho\in \Pi(T)$ is $\beta$-coherent.
\end{enumerate}
\end{proposition}

\begin{proof}
By Definition~\ref{def:beta-coherent-prefix}, condition (ii) is exactly the statement that the lifted child witnesses supported by $\rho$ verify the profitable-prefix hypothesis of Theorem~\ref{thm:adaptive-kernel-prefix}; this immediately implies (i).

Conversely, suppose $T\in \mathcal P_\lambda$, and let $\rho$ be a profitable prefix witnessing this. For every proper support piece $P$ whose descendants are counted in the profitable-prefix witness, exact trace recursion and exact terminal-kernel reduction identify a profitable child in $\mathcal C_T(P)$ whose profitable prefix pulls back to $\rho$. Hence $P\in \Gamma_T(\rho)$ for every support piece used by the witness, so the lifted family through $\rho$ satisfies the same profitable-prefix inequality. Therefore $\rho$ is $\beta$-coherent.
\end{proof}

\begin{conjecture}[$\beta$-coherent lift hypothesis]\label{conj:beta-coherent-lift}
For every fixed $k$ and every neutral terminal kernel $T$, the child-profit lift hypergraph $H_T$ has a $\beta$-coherent vertex. Equivalently, some prefix $\rho\in \Pi(T)$ lifts enough profitable child witnesses to force $T\in \mathcal P_\lambda$.
\end{conjecture}

\begin{corollary}[Conditional consequence of the $\beta$-coherent lift hypothesis]\label{cor:beta-coherent-lift}
Fix $k$. Suppose the solved class $\mathcal S_k$ admits an exponential bound $|\Hh|\le \lambda_k^{\er(\Hh)}$, and suppose Conjecture~\ref{conj:beta-coherent-lift} holds for the same value of $\lambda_k$. Then there exists $c_k<\infty$ such that
\[
f(n,k)<c_k^n
\qquad\text{for all }n.
\]
\end{corollary}

\begin{proof}
Apply Theorem~\ref{thm:forced-trichotomy}. If a reduced kernel link lies in $\mathcal S_k$, the solved-class bound applies. If it is profitable, Theorem~\ref{thm:adaptive-kernel-prefix} gives the required exponential estimate. Otherwise it has a neutral terminal descendant $T$. By Conjecture~\ref{conj:beta-coherent-lift}, the lift hypergraph $H_T$ has a $\beta$-coherent vertex. Proposition~\ref{prop:coherent-lift} then implies $T\in \mathcal P_{\lambda_k}$, contradicting neutrality. Hence neutral terminals do not exist, and Corollary~\ref{cor:no-neutral-terminals} closes the recursion.
\end{proof}

\begin{remark}[Why Theorem~\ref{thm:neutral-terminal-child} does not by itself close the routing step]\label{rem:neutral-child-not-enough}
Theorem~\ref{thm:neutral-terminal-child} only says that every edge $\Lambda_T(P)$ of the child-profit lift hypergraph is nonempty. That does \emph{not} force a single parent prefix to lie in enough edges to trigger Theorem~\ref{thm:adaptive-kernel-prefix}. The full width-$(k-1)$ transversal product already exhibits the failure mode: different support pieces can produce profitable children whose parent lifts are mutually incompatible. Thus the missing ingredient is a \emph{coherence theorem}, not another pointwise child-profit statement.
\end{remark}

\begin{example}[Triangle obstruction to an unweighted lift]
At the hypergraph level the smallest nondegenerate obstruction already appears on three support pieces:
\[
\Lambda_T(P_1)=\{a,b\},\qquad
\Lambda_T(P_2)=\{b,c\},\qquad
\Lambda_T(P_3)=\{c,a\}.
\]
Every edge is nonempty, but no single vertex lies in all three edges. Thus pointwise child profitability cannot be upgraded to a parent lift without an additional overlap hypothesis.
\end{example}

\begin{definition}[Weighted coherence of a lift hypergraph]\label{def:lift-hypergraph-coherence}
Let $H=(V,E)$ be a finite hypergraph and let $\mu$ be a probability measure on $E$. For $v\in V$ define its $\mu$-degree by
\[
d_\mu(v):=\mu(\{e\in E:v\in e\}).
\]
Define the $\mu$-coherence of $H$ by
\[
\operatorname{Coh}_\mu(H):=
\frac{\sum_{v\in V} d_\mu(v)^2}{\sum_{v\in V} d_\mu(v)},
\]
provided the denominator is nonzero. Equivalently,
\[
\operatorname{Coh}_\mu(H)=
\frac{\mathbb E_{e,e'\sim \mu}|e\cap e'|}{\mathbb E_{e\sim \mu}|e|}.
\]
\end{definition}

\begin{proposition}[Coherence threshold forces a heavy vertex]\label{prop:coherence-heavy-vertex}
Let $H=(V,E)$ be a finite hypergraph with an edge-probability measure $\mu$. If
\[
\operatorname{Coh}_\mu(H)\ge \beta,
\]
then some vertex $v\in V$ satisfies
\[
d_\mu(v)\ge \beta.
\]
\end{proposition}

\begin{proof}
Set
\[
M:=\sum_{v\in V} d_\mu(v)=\mathbb E_{e\sim \mu}|e|.
\]
If $M=0$, every edge is empty and the claim is vacuous. Otherwise
\[
\operatorname{Coh}_\mu(H)
=
\sum_{v\in V}\frac{d_\mu(v)}{M}\,d_\mu(v),
\]
so $\operatorname{Coh}_\mu(H)$ is a weighted average of the values $d_\mu(v)$. Hence
\[
\operatorname{Coh}_\mu(H)\le \max_{v\in V} d_\mu(v).
\]
If $\operatorname{Coh}_\mu(H)\ge \beta$, some vertex has $\mu$-degree at least $\beta$.
\end{proof}

\begin{corollary}[Coherence threshold forces a profitable parent prefix]\label{cor:terminal-coherence-threshold}
Let $T$ be a terminal kernel, and endow its child-profit lift hypergraph $H_T$ with a probability measure $\mu$ on the edges. Suppose that every prefix $\rho\in \Pi(T)$ with
\[
d_\mu(\rho)\ge \beta
\]
is $\beta$-coherent in the sense of Definition~\ref{def:beta-coherent-prefix}. If
\[
\operatorname{Coh}_\mu(H_T)\ge \beta,
\]
then $T\in \mathcal P_\lambda$.
\end{corollary}

\begin{proof}
By Proposition~\ref{prop:coherence-heavy-vertex}, some prefix $\rho\in \Pi(T)$ has $d_\mu(\rho)\ge \beta$. By hypothesis, $\rho$ is then $\beta$-coherent. Proposition~\ref{prop:coherent-lift} now gives $T\in \mathcal P_\lambda$.
\end{proof}

\begin{corollary}[Neutral terminals lie in a low-coherence / large-fractional-matching regime]\label{cor:neutral-low-coherence}
In the setting of Corollary~\ref{cor:terminal-coherence-threshold}, if $T$ is a genuine neutral terminal then
\[
\operatorname{Coh}_\mu(H_T)<\beta.
\]
Moreover the normalized edge weights
\[
x_P:=\frac{\mu(P)}{\beta}
\]
satisfy the fractional-matching inequalities
\[
\sum_{P:\,\rho\in \Lambda_T(P)} x_P<1
\qquad\text{for every }\rho\in \Pi(T).
\]
\end{corollary}

\begin{proof}
If $\operatorname{Coh}_\mu(H_T)\ge \beta$, Corollary~\ref{cor:terminal-coherence-threshold} would imply $T\in \mathcal P_\lambda$, contradicting neutrality. Hence $\operatorname{Coh}_\mu(H_T)<\beta$. Since $T$ is neutral and any prefix with $d_\mu(\rho)\ge \beta$ would be $\beta$-coherent, in fact
\[
d_\mu(\rho)<\beta
\qquad\text{for every }\rho\in \Pi(T).
\]
Therefore
\[
\sum_{P:\,\rho\in \Lambda_T(P)} x_P
=
\frac{d_\mu(\rho)}{\beta}
<1,
\]
which is exactly the claimed fractional-matching bound.
\end{proof}

\begin{remark}[What remains missing after the coherence threshold lemma]\label{rem:low-coherence-barrier}
Corollary~\ref{cor:terminal-coherence-threshold} does \emph{not} prove Conjecture~\ref{conj:beta-coherent-lift}. It only shows that any genuine neutral terminal must lie in the complementary low-coherence regime described by Corollary~\ref{cor:neutral-low-coherence}. Thus the unresolved structural statement is no longer another pointwise child-profit lemma; it is a classification or stability theorem for low-coherence lift hypergraphs, with the full width-$(k-1)$ product as the model obstruction.
\end{remark}

\begin{remark}[Why the direct Farkas-separation closure attempt is not yet rigorous]\label{rem:farkas-gap}
A tempting LP/Farkas argument tries to assign affine scores to profitable children and then telescope the resulting boundary terms back to the parent. At present this does not produce a proof of Conjecture~\ref{conj:beta-coherent-lift}: the required cancellation of boundary discrepancies is exactly the missing global compatibility statement. In other words, dual separation explains why a weighted coherence theorem would be natural, but it does not by itself certify that every neutral terminal already has such a coherent weighting.
\end{remark}



\subsection*{Exact convex lift theorem for weighted profitable children}
The previous remark identifies the correct place where the naive Farkas closure attempt overreaches.
The LP part can in fact be made exact; what remains missing is the structural step that turns a dual separator into a forbidden solved witness.

\begin{definition}[Boundary discrepancy space for a neutral terminal]\label{def:boundary-space}
Let $T$ be a neutral terminal.
Define $\Omega_T$ to be the finite set of oriented interface-trace atoms of $T$, meaning that an atom records a boundary slot, an orientation, and the local trace state carried by the exact trace recursion.
Set
\[
U_T:=\R^{\Omega_T}.
\]
Let $R_T\subseteq U_T$ be the subspace generated by the local gluing identities:
opposite orientations of the same internal interface with the same trace state cancel, and any further linear identities already built into the exact trace recursion are also imposed.
The \emph{boundary discrepancy space} of $T$ is the finite-dimensional quotient
\[
V_T:=U_T/R_T.
\]
\end{definition}

\begin{definition}[Weighted child-lift data]\label{def:weighted-child-lift}
Fix a neutral terminal $T$ and a probability mass vector $(\mu_P)_{P\in \mathcal P(T)}$ on its proper support pieces.
For each proper support piece $P$, choose a nonempty set $C(P)$ of profitable children.
A \emph{weighted child-lift datum} consists of:
\begin{enumerate}[label=(\roman*),leftmargin=2.2em]
\item a positive gain $g_{P,c}>0$ for each $c\in C(P)$;
\item a boundary class $\beta_{P,c}\in V_T$ for each $c\in C(P)$;
\item a target class $t_T\in V_T$ for the desired profitable parent lift.
\end{enumerate}
\end{definition}

\begin{proposition}[Assembly lemma for weighted profitable children]\label{prop:assembly-weighted-children}
Fix a neutral terminal $T$ together with weighted child-lift data as in Definition~\ref{def:weighted-child-lift}.
Suppose there exist nonnegative weights $x_{P,c}$ such that
\[
\sum_{c\in C(P)} x_{P,c}=\mu_P
\qquad\text{for every }P\in \mathcal P(T),
\]
and
\[
\sum_{P\in \mathcal P(T)}\sum_{c\in C(P)} x_{P,c}\,\beta_{P,c}=t_T
\qquad\text{in }V_T.
\]
Then $T$ has a profitable parent prefix.
\end{proposition}

\begin{proof}
By construction, each profitable child $c\in C(P)$ comes with a strict gain $g_{P,c}>0$ and a boundary contribution class $\beta_{P,c}$.
Summing the exact child-to-parent lift identities with weights $x_{P,c}$ is legitimate because exact trace recursion is linear.
The weighted boundary remainder is
\[
\sum_{P,c} x_{P,c}\,\beta_{P,c}-t_T,
\]
which is zero in $V_T$.
Hence this remainder lies in the relation subspace $R_T$, so all leftover interface terms cancel after gluing.

Let
\[
\eta_T:=\min\{g_{P,c}:P\in \mathcal P(T),\ c\in C(P)\}.
\]
Because $T$ and all sets $C(P)$ are finite, one has $\eta_T>0$.
The total gain of the assembled parent lift is therefore
\[
\sum_{P,c} x_{P,c}g_{P,c}
\ge
\eta_T \sum_{P\in \mathcal P(T)} \sum_{c\in C(P)} x_{P,c}
=
\eta_T \sum_{P\in \mathcal P(T)} \mu_P
=
\eta_T
>
0.
\]
So the assembled lift is a genuine profitable parent prefix.
\end{proof}

\begin{theorem}[Exact convex lift theorem]\label{thm:exact-convex-lift}
Fix a neutral terminal $T$ together with weighted child-lift data as in Definition~\ref{def:weighted-child-lift}.
Define the weighted Minkowski sum
\[
K_T(\mu):=
\sum_{P\in \mathcal P(T)} \mu_P\,\operatorname{conv}\{\beta_{P,c}:c\in C(P)\}
\subseteq V_T.
\]
Then exactly one of the following holds:
\begin{enumerate}[label=(\roman*),leftmargin=2.2em]
\item $t_T\in K_T(\mu)$, in which case $T$ has a profitable parent prefix;
\item $t_T\notin K_T(\mu)$, in which case there exists a linear functional $\phi\in V_T^\ast$ such that
\[
\phi(t_T)
<
\sum_{P\in \mathcal P(T)} \mu_P\,\min_{c\in C(P)} \phi(\beta_{P,c}).
\]
\end{enumerate}
\end{theorem}

\begin{proof}
A point $y\in K_T(\mu)$ has the form
\[
y=\sum_{P\in \mathcal P(T)} \mu_P y_P,
\qquad
y_P\in \operatorname{conv}\{\beta_{P,c}:c\in C(P)\}.
\]
Unpacking each convex combination, this is equivalent to the existence of nonnegative numbers $x_{P,c}$ with
\[
\sum_{c\in C(P)} x_{P,c}=\mu_P
\qquad\text{for every }P,
\]
such that
\[
y=\sum_{P,c} x_{P,c}\,\beta_{P,c}.
\]
Therefore $t_T\in K_T(\mu)$ is exactly the feasibility condition in Proposition~\ref{prop:assembly-weighted-children}, and case (i) implies that $T$ has a profitable parent prefix.

Now assume $t_T\notin K_T(\mu)$.
Since each set $\operatorname{conv}\{\beta_{P,c}:c\in C(P)\}$ is compact and convex, the Minkowski sum $K_T(\mu)$ is also compact and convex.
By the separating-hyperplane theorem, there exist $\phi\in V_T^\ast$ and $\alpha\in \R$ such that
\[
\phi(t_T)<\alpha\le \phi(y)
\qquad\text{for every }y\in K_T(\mu).
\]
Taking the infimum over $K_T(\mu)$ gives
\[
\phi(t_T)
<
\inf_{y\in K_T(\mu)} \phi(y).
\]
Because $\phi$ is linear and each summand is a scaled convex hull,
\[
\inf_{y\in K_T(\mu)} \phi(y)
=
\sum_{P\in \mathcal P(T)} \mu_P
\inf_{y_P\in \operatorname{conv}\{\beta_{P,c}:c\in C(P)\}} \phi(y_P).
\]
A linear functional attains its minimum on a finite convex hull at an extreme point, so
\[
\inf_{y_P\in \operatorname{conv}\{\beta_{P,c}:c\in C(P)\}} \phi(y_P)
=
\min_{c\in C(P)} \phi(\beta_{P,c}).
\]
Substituting this identity gives
\[
\phi(t_T)
<
\sum_{P\in \mathcal P(T)} \mu_P\,\min_{c\in C(P)} \phi(\beta_{P,c}),
\]
which is exactly case (ii).

The two cases are mutually exclusive because the inequality in (ii) cannot hold when $t_T\in K_T(\mu)$.
\end{proof}

\begin{corollary}[Discrepancy form of the convex alternative]\label{cor:exact-convex-lift-discrepancy}
Choose vectors $\theta_P\in V_T$ with
\[
t_T=\sum_{P\in \mathcal P(T)} \mu_P\,\theta_P,
\]
and define
\[
\partial_{P,c}:=\beta_{P,c}-\theta_P.
\]
Then exactly one of the following holds:
\begin{enumerate}[label=(\roman*),leftmargin=2.2em]
\item
\[
0\in \sum_{P\in \mathcal P(T)} \mu_P\,\operatorname{conv}\{\partial_{P,c}:c\in C(P)\},
\]
in which case $T$ has a profitable parent prefix;
\item there exists $\phi\in V_T^\ast$ such that
\[
\sum_{P\in \mathcal P(T)} \mu_P\,\min_{c\in C(P)} \phi(\partial_{P,c})>0.
\]
\end{enumerate}
\end{corollary}

\begin{proof}
Subtract $\sum_P \mu_P\theta_P=t_T$ from every point of $K_T(\mu)$.
This translates the convex body by $-t_T$ and replaces each $\beta_{P,c}$ by $\partial_{P,c}$.
Applying Theorem~\ref{thm:exact-convex-lift} to the translated body yields the stated alternative.
\end{proof}

\begin{remark}[What the convex lift theorem does \emph{not} yet do]\label{rem:exact-convex-lift-gap}
Theorem~\ref{thm:exact-convex-lift} repairs the LP/Farkas part of the weighted child-to-parent lift completely.
What it still does \emph{not} provide is a way to convert the separating functional in case~(ii) into a forbidden fixed-memory strict-gap class for an arbitrary neutral terminal.
That conversion requires extra structure, typically some support-box or exact product decomposition on which every relevant functional is realized by a finite-memory observable.
Thus the remaining gap is structural, not linear-programming.
\end{remark}

\status{The profitable-child theorem and the coherence-threshold lemma are rigorous. The LP/Farkas part of the weighted child-to-parent lift can be repaired into the exact convex alternative of Theorem~\ref{thm:exact-convex-lift}. What remains missing is the structural theorem ruling out dual obstruction functionals on the complementary low-coherence / large-fractional-matching regime.}

\subsection*{Exposed simplicial cores and the local splice gap}
The convex part of the terminal obstruction theory can be pushed a little further.
What emerges is a canonical finite-support core.
The first genuine missing input is then a local splice principle on that core.

\begin{definition}[Support-irreducible affine witness model]\label{def:affine-witness-model}
Let $T$ be a low-coherence neutral terminal.
An \emph{affine witness model} for $T$ consists of a finite-dimensional real vector space $V_T$, an affine slice $A_T\subseteq V_T$, and a compact convex body $K_T\subseteq V_T$ such that the weighted child-to-parent lift problem is equivalent to
\[
A_T\cap K_T\neq \varnothing.
\]
We say that the model is \emph{support-irreducible} if every proper active support restriction of the failed system is feasible; equivalently, every obstruction requires the full active support of the chosen witness.
\end{definition}

\begin{proposition}[Canonical exposed simplicial core]\label{prop:canonical-exposed-core}
Assume $T$ is a support-irreducible low-coherence neutral terminal with an affine witness model $(V_T,A_T,K_T)$ in the sense of Definition~\ref{def:affine-witness-model}.
Assume furthermore that the lift fails, so $A_T\cap K_T=\varnothing$, and let
\[
F=\operatorname{conv}(x_0,\dots,x_m)\subseteq K_T
\]
be a rational exposed simplicial strict-gap witness, where $m\le \dim(V_T)$.
Then:
\begin{enumerate}[label=(\roman*),leftmargin=2.2em]
\item $A_T\cap \operatorname{aff}(F)$ is a single point $q$;
\item $q\notin F$;
\item $q$ lies in the affine hull of no proper face of $F$;
\item writing
\[
q=\sum_{i=0}^m \beta_i x_i,
\qquad
\sum_{i=0}^m \beta_i=1,
\]
all coefficients $\beta_i$ are nonzero, and at least one of them is negative;
\item therefore the augmented support
\[
\{q,x_0,\dots,x_m\}
\]
is a rational signed affine circuit of size at most $\dim(V_T)+2$.
\end{enumerate}
\end{proposition}

\begin{proof}
Set
\[
L:=A_T\cap \operatorname{aff}(F).
\]
We first show that $L\neq \varnothing$.
If $L$ were empty, then any singleton active support $\{x_i\}$ would already fail:
\[
A_T\cap \operatorname{conv}\{x_i\}=A_T\cap \{x_i\}=\varnothing.
\]
This would contradict support-irreducibility, because $\{x_i\}$ is a proper active support whenever the witness uses more than one vertex.

We next show that $L$ cannot have positive dimension.
Assume for contradiction that $\dim L>0$.
Then $L$ contains an affine line
\[
q_0+t v
\qquad (t\in \R)
\]
with $v\neq 0$.
Because $F$ is a simplex, the barycentric coordinate functions
\[
\lambda_0,\dots,\lambda_m:\operatorname{aff}(F)\to \R
\]
are affine, satisfy $\sum_i \lambda_i\equiv 1$, and obey
\[
y\in F \iff \lambda_i(y)\ge 0 \text{ for all }i.
\]
Along the line $q_0+t v$ each function has the form
\[
\lambda_i(q_0+t v)=a_i+t b_i.
\]
Not all $b_i$ vanish, because otherwise the line would be constant.
Since $\sum_i b_i=0$, some $b_j$ is nonzero.
Choose $t_*$ so that
\[
\lambda_j(q_0+t_*v)=0.
\]
Then
\[
q_*:=q_0+t_*v\in L\cap \operatorname{aff}(F_j),
\]
where $F_j$ is the proper face of $F$ opposite $x_j$.

Because $F$ is a strict-gap witness, $A_T\cap F=\varnothing$.
Hence $q_*\notin F_j$.
So the proper face $F_j$ already carries a failed support restriction, contradicting support-irreducibility.
Therefore $\dim L=0$, and so
\[
L=\{q\}
\]
for a unique point $q$.

By construction,
\[
q\in A_T\cap \operatorname{aff}(F).
\]
Because $A_T\cap F=\varnothing$, one has $q\notin F$, proving (ii).

Now suppose that $q$ lies in the affine hull of a proper face $F'\subsetneq F$.
Then $q\notin F'$, again because $q\notin F$.
But then the support restriction to the vertices of $F'$ already fails, contradicting support-irreducibility.
So $q$ lies in the affine hull of no proper face of $F$.

Since $F$ is simplicial, $q$ has unique barycentric coordinates
\[
q=\sum_{i=0}^m \beta_i x_i,
\qquad
\sum_{i=0}^m \beta_i=1.
\]
If some $\beta_i=0$, then $q$ lies in the affine hull of the proper face opposite $x_i$, contradiction.
Thus every $\beta_i$ is nonzero.
If all $\beta_i$ were nonnegative, then $q\in F$, again a contradiction.
Hence at least one coefficient is negative.
This is exactly a rational signed affine circuit on the set $\{q,x_0,\dots,x_m\}$.
Because $m\le \dim(V_T)$, its size is at most $\dim(V_T)+2$.
\end{proof}

\begin{definition}[Primitive exposed core]\label{def:primitive-exposed-core}
Let $T$ be a low-coherence support-irreducible neutral terminal.
After quotienting by the solved-trace lineality, let $V_T$ be the resulting finite-dimensional space and let $C_T\subseteq V_T$ be the closed cone of good one-step witness directions.
A \emph{primitive exposed core} consists of data
\[
\mathfrak C=\bigl(T;\,x_1,\dots,x_m;\,\alpha_1,\dots,\alpha_m;\,d;\,F;\,\phi_1,\dots,\phi_m\bigr)
\]
such that:
\begin{enumerate}[label=(\roman*),leftmargin=2.2em]
\item $m\ge 2$ and $\alpha_i>0$ for every $i$;
\item
\[
d=\sum_{i=1}^m \alpha_i x_i\in \operatorname{ri}(C_T),
\]
so $d$ is a full-support one-step sibling defect;
\item no single $x_i$ is already a genuine one-step witness, i.e.
\[
x_i\notin C_T\setminus \{0\}
\qquad\text{for all }i;
\]
\item the family is deletion-minimal:
\[
\operatorname{cone}(x_j:j\neq i)\cap \operatorname{ri}(C_T)=\varnothing
\qquad\text{for all }i;
\]
\item $F$ is the minimal exposed face of $\operatorname{cone}(x_1,\dots,x_m)$ containing $d$, and each $x_i$ spans an extreme ray of $F$;
\item each $\phi_i\in V_T^*$ is an essential separator satisfying
\[
\phi_i(\operatorname{ri}(C_T))<0,
\qquad
\phi_i(x_i)<0,
\qquad
\phi_i(x_j)\ge 0 \text{ for }j\neq i;
\]
\item after exact trace normalization, each $\phi_i$ localizes to a proper support piece of $T$ carrying a profitable child.
\end{enumerate}
\end{definition}

\begin{proposition}[Failure of local uncrossing compresses to a primitive exposed core]\label{prop:primitive-core-reduction}
Assume the one-step uncrossing step fails somewhere in the exact convex lift model.
Then some counterexample compresses to a primitive exposed core with
\[
2\le m\le \dim(V_T).
\]
\end{proposition}

\begin{proof}
Choose a counterexample $T$ that is lexicographically minimal in support size, then in the number of active children, then in the dimension of the face generated by those children.
Because the uncrossing step fails, there is a full-support sibling defect
\[
d=\sum_{i=1}^m \alpha_i x_i\in \operatorname{ri}(C_T),
\qquad
\alpha_i>0,
\]
with no genuine single-child witness among the active child vectors $x_i$.

By conic Caratheodory, after discarding redundant rays we may assume
\[
m\le \dim(V_T).
\]
Let
\[
K:=\operatorname{cone}(x_1,\dots,x_m),
\]
and let $F$ be the minimal face of $K$ containing $d$.
Replacing the active family by the extreme rays of $F$, we may assume that each $x_i$ spans an extreme ray of $F$ and that $d\in \operatorname{ri}(F)$.

Fix $i$ and set
\[
K_{-i}:=\operatorname{cone}(x_j:j\neq i).
\]
By minimality of the active family,
\[
K_{-i}\cap \operatorname{ri}(C_T)=\varnothing.
\]
Since $K_{-i}$ is closed and $\operatorname{ri}(C_T)$ is open and convex, strict separation gives a nonzero functional $\phi_i\in V_T^*$ such that
\[
\phi_i(K_{-i})\ge 0
\qquad\text{and}\qquad
\phi_i(\operatorname{ri}(C_T))<0.
\]
Writing
\[
d=\alpha_i x_i+y_i
\qquad\text{with } y_i\in K_{-i},
\]
we obtain
\[
0>\phi_i(d)=\alpha_i\phi_i(x_i)+\phi_i(y_i).
\]
Because $\phi_i(y_i)\ge 0$, it follows that
\[
\phi_i(x_i)<0.
\]
Thus every active child is indispensable, and every indispensable child has its own exposed dual separator.

Apply the exact trace-normalization / adaptive kernel-prefix criterion to each $\phi_i$.
The parent-prefix branch is excluded because $T$ is low-coherence neutral.
The global solved branch is excluded because $T$ was chosen as the first unresolved object.
Hence each $\phi_i$ localizes to a proper support piece of $T$.
By Theorem~\ref{thm:neutral-terminal-child}, every proper support piece of a neutral terminal has a profitable child.
So each localized separator indeed lands on a proper support piece carrying a profitable child.

Collecting these data gives a primitive exposed core in the sense of Definition~\ref{def:primitive-exposed-core}.
\end{proof}

\begin{remark}[The first real gap in the exposed-core lane]\label{rem:primitive-core-gap}
The reduction above does not prove uncrossing.
What it isolates is the exact missing local statement:
one needs a splice principle showing that a profitable child on one of the proper exposed support pieces can be inserted back into the ambient one-step lift without destroying the global boundary constraints.
Without that local splice principle, the correct endpoint of the present argument is the primitive exposed core, not a full uncrossing theorem.
\end{remark}


\subsection*{The residual low-coherence class, one-step witnesses, and visible split forcing}
The coherence threshold isolates the remaining obstruction class quite sharply.
What is still missing is a local witness theorem for that class.

\begin{proposition}[Residual support-irreducible low-coherence class]\label{prop:low-coherence-residual-class}
Let $T$ be any unresolved terminal kernel in the sense of this section.
Then:
\begin{enumerate}[label=(\roman*),leftmargin=2.2em]
\item $T$ is neutral terminal;
\item
\[
\operatorname{Coh}_\mu(H_T)<\beta;
\]
\item every proper support piece of $T$ has a profitable child.
\end{enumerate}
Equivalently, every remaining obstruction is already reduced to a support-irreducible low-coherence neutral terminal.
\end{proposition}

\begin{proof}
By Theorem~\ref{thm:forced-trichotomy}, every unresolved reduced-kernel link reaches a neutral terminal descendant.
Applying Theorem~\ref{thm:neutral-terminal-child} to that descendant gives profitable children on every proper support piece.
Finally, Corollary~\ref{cor:neutral-low-coherence} shows that a genuine neutral terminal must satisfy $\operatorname{Coh}_\mu(H_T)<\beta$.
\end{proof}

\begin{conjecture}[One-step witness theorem]\label{conj:one-step-witness}
For each fixed $k$, every support-irreducible low-coherence neutral terminal $T$ admits at least one of the following local witnesses:
\begin{enumerate}[label=(\roman*),leftmargin=2.2em]
\item a solved witness already covered by the proved machinery, such as a profitable prefix block or a support-box decomposition into solved strict-gap classes;
\item a thin neutral obstruction witness, such as a balanced slice, a mesoscopic product, or an explicit interface to a downstream hard model.
\end{enumerate}
\end{conjecture}

\begin{remark}[Why Conjecture~\ref{conj:one-step-witness} is the real next theorem]\label{rem:one-step-witness-minimal}
The proved recursion already handles everything \emph{after} a usable local witness has been extracted.
Thus the first genuinely missing theorem is not another recursive closure statement.
It is a one-step witness theorem for the residual low-coherence class of Proposition~\ref{prop:low-coherence-residual-class}.
Any stronger multiscale extraction theorem must factor through such a local witness step.
\end{remark}

\begin{definition}[Finite visible quotient of lifted child witnesses]\label{def:visible-witness-quotient}
Fix $k$.
A \emph{finite visible quotient} of lifted child witnesses consists of a finite set $\Sigma_k$ and a map from lifted child witnesses to $\Sigma_k$ such that each class records the visible next-step support set and the visible child-transition data needed to continue the exact trace recursion.
\end{definition}

\begin{conjecture}[Visible split forcing]\label{conj:visible-split-forcing}
Fix $k$ and a finite visible quotient $\Sigma_k$ as above.
At any conditioned node of a neutral terminal kernel, if two lifted child witnesses induce different visible next-step support sets, then that node already has a profitable prefix.
\end{conjecture}

\begin{proposition}[One-disagreement forcing from visible split forcing]\label{prop:visible-split-to-coherence}
Assume Conjecture~\ref{conj:visible-split-forcing}.
Let $T$ be a neutral terminal kernel with no profitable prefix.
Then every conditioned node of $T$ has a unique visible lifted-child state in $\Sigma_k$.
\end{proposition}

\begin{proof}
Assume for contradiction that some conditioned node admits two lifted child witnesses with different visible states.
Choose such a counterexample with minimal remaining trace depth.
Let the two visible states be $\sigma_1,\sigma_2\in \Sigma_k$.

If their visible next-step support sets differ, Conjecture~\ref{conj:visible-split-forcing} gives a profitable prefix immediately, contradicting the hypothesis.
So the support sets must agree.
Because $\sigma_1\ne \sigma_2$, there is some visible next symbol $a$ for which the visible child states differ.
Condition on that symbol.
Exact trace recursion passes to a strictly smaller-depth conditioned node.
If the child had a profitable prefix, lifting it back would make the parent profitable, contradiction.
Thus the child is a smaller counterexample, contradicting minimality.
Therefore no such pair of visible states can exist.
\end{proof}

\begin{proposition}[Automaton reconstruction from unique visible witnesses]\label{prop:visible-automaton-reconstruction}
Suppose a neutral terminal kernel $T$ has the property that every conditioned node has a unique visible lifted-child state in some finite quotient $\Sigma_k$.
Then the exact trace tree of $T$ is generated by a finite deterministic recursive system.
If one of the visible states already encodes an exact block product or a positive-density support box, then $T$ is solved by the theorem-producing classes already proved in this paper.
Otherwise $T$ lies in a fixed-memory finite-state branch class.
\end{proposition}

\begin{proof}
For each conditioned node $u$, let $\sigma(u)\in \Sigma_k$ denote its unique visible state.
Because $\Sigma_k$ is finite, the map $u\mapsto \sigma(u)$ takes only finitely many values.
Uniqueness means that whenever a visible symbol $a$ is admissible at $u$, the next visible state is determined:
\[
\delta(\sigma(u),a)=\sigma(ua).
\]
Hence the exact trace tree is recognized by a finite deterministic labeled graph on the state set $\Sigma_k$.

If one of the states is already recognized as an exact block-product or positive-density support-box state, then the corresponding descendants are solved by Section~\ref{sec:block-products}.
If no such decomposition occurs, the deterministic graph itself is a fixed-memory finite-state branch presentation, so Section~\ref{sec:finite-state} applies.
\end{proof}

\begin{theorem}[Conditional reduction to visible split forcing]\label{thm:visible-split-reduction}
Assume that for each fixed $k$:
\begin{enumerate}[label=(\roman*),leftmargin=2.2em]
\item lifted child witnesses admit a finite visible quotient $\Sigma_k$;
\item Conjecture~\ref{conj:visible-split-forcing} holds.
\end{enumerate}
Then every neutral terminal kernel is recursively solved.
Consequently, once the solved classes are given their certified exponential bounds, one obtains
\[
f(n,k)<c_k^n
\]
for some constant $c_k$ depending only on $k$.
\end{theorem}

\begin{proof}
Take a neutral terminal kernel $T$.
If $T$ already has a profitable prefix, it is handled by Theorem~\ref{thm:adaptive-kernel-prefix}.
Otherwise Proposition~\ref{prop:visible-split-to-coherence} shows that every conditioned node has a unique visible lifted-child state.
Proposition~\ref{prop:visible-automaton-reconstruction} then shows that $T$ is recursively solved, either because a product/box state appears or because the entire trace tree is finite-state.

Thus neutral terminals do not contribute any new obstruction beyond the already solved classes.
Combining this with Theorem~\ref{thm:forced-trichotomy}, Theorem~\ref{thm:adaptive-kernel-prefix}, and Theorem~\ref{thm:unionclass-closure} closes the recursion and yields the stated exponential bound.
\end{proof}

\begin{remark}[Relation between the two local conjectures]\label{rem:witness-vsf}
Conjecture~\ref{conj:one-step-witness} is the broad local-witness formulation.
Conjecture~\ref{conj:visible-split-forcing} is a sharper, more automaton-oriented version targeted at the exact point where two lifted child witnesses first differ.
Either conjecture would supply the missing local input.
The visible-split formulation is smaller and better aligned with the finite-state endgame.
\end{remark}

\subsection*{Singleton traces, separator descent, and the sharp coherence threshold}
The visible-split route has a clean hypergraph consequence once one isolates the correct finite witness configuration.

\begin{definition}[$k$-singleton trace]\label{def:k-singleton-trace}
Let $T$ be a neutral terminal kernel.
For each proper support piece $\sigma$ of $T$, write
\[
E_\sigma:=\Lambda_T(\sigma)\subseteq \Pi(T).
\]
A set
\[
Q=\{\rho_1,\dots,\rho_k\}\subseteq \Pi(T)
\]
is called a \emph{$k$-singleton trace} if there exist pairwise distinct proper support pieces $\sigma_1,\dots,\sigma_k$ such that
\[
E_{\sigma_i}\cap Q=\{\rho_i\}
\qquad (1\le i\le k).
\]
\end{definition}

\begin{proposition}[Residual functoriality plus visible split forcing imply singleton-trace forcing]\label{prop:vsf-to-singleton}
Assume the residual functoriality identity
\[
\Lambda_{T|u}(\sigma|u)\cap (Q/u)
=
\{\rho/u:\rho\in \Lambda_T(\sigma)\cap Q\}
\]
whenever $u$ is a common visible prefix of every $\rho\in Q$.
Assume also Conjecture~\ref{conj:visible-split-forcing}.
Then every $k$-singleton trace in a neutral terminal forces a profitable parent prefix.
\end{proposition}

\begin{proof}
Let
\[
Q=\{\rho_1,\dots,\rho_k\}\subseteq \Pi(T)
\]
be a $k$-singleton trace witnessed by pairwise distinct support pieces $\sigma_1,\dots,\sigma_k$, so that
\[
E_{\sigma_i}\cap Q=\{\rho_i\}
\qquad (1\le i\le k).
\]

Adjoin a terminal stop symbol $\bot$ to the visible alphabet and view the prefixes $\rho_i$ as leaves of the visible trie in the extended alphabet.
Let $u$ be the longest common visible prefix of the family $Q$.
Equivalently, $u$ is the first visible branching node of the trie of $Q$.
Write
\[
\eta_i:=\rho_i/u.
\]
Because the $\rho_i$ are distinct, the visible next-step supports of the residual suffixes $\eta_i$ are not all equal.
So there exist $i\neq j$ for which $\eta_i$ and $\eta_j$ have different visible next-step supports at the conditioned node $u$.

By the stated residual functoriality identity,
\[
\Lambda_{T|u}(\sigma_i|u)\cap (Q/u)
=
\{\rho/u:\rho\in \Lambda_T(\sigma_i)\cap Q\}
=
\{\eta_i\}.
\]
Hence the conditioned family $Q/u$ is again a $k$-singleton trace in the conditioned terminal $T|u$.
In particular, $\eta_i$ and $\eta_j$ are lifted profitable child witnesses at $u$ with different visible next-step supports.
Conjecture~\ref{conj:visible-split-forcing} therefore gives a profitable prefix at $u$.
Lifting this profitable prefix back through exact trace recursion yields a profitable parent prefix of $T$.
\end{proof}

\begin{lemma}[Private-edge lemma]\label{lem:private-edge}
Let $X\subseteq \Pi(T)$ be an inclusion-minimal transversal of the hypergraph
\[
H_T=\bigl(\Pi(T),\{E_\sigma:\sigma\text{ a proper support piece of }T\}\bigr).
\]
Then for every $x\in X$ there exists an edge $E_x$ such that
\[
E_x\cap X=\{x\}.
\]
\end{lemma}

\begin{proof}
Fix $x\in X$.
If no edge met $X$ only in $x$, then every edge containing $x$ would also meet $X\setminus\{x\}$.
Every edge not containing $x$ already meets $X\setminus\{x\}$.
Thus $X\setminus\{x\}$ would still be a transversal, contradicting the inclusion-minimality of $X$.
\end{proof}

\begin{theorem}[Singleton traces force a sharp heavy prefix]\label{thm:singleton-heavy-prefix}
Assume every $k$-singleton trace in a neutral terminal produces a profitable parent prefix.
Let $T$ be a neutral terminal kernel, and let $\mu_T$ be any probability measure on the proper support pieces of $T$.
Define
\[
d_T(\rho):=\sum_{\sigma:\,\rho\in E_\sigma}\mu_T(\sigma),
\qquad
E_\sigma:=\Lambda_T(\sigma)\subseteq \Pi(T).
\]
Then there exists a prefix $\rho\in \Pi(T)$ satisfying
\[
d_T(\rho)\ge \frac{1}{k-1}.
\]
\end{theorem}

\begin{proof}
If the lift hypergraph $H_T$ already contains a $k$-singleton trace, the conclusion about profitable parent prefixes is part of the hypothesis.
So suppose that $H_T$ contains no $k$-singleton trace.

Let $X\subseteq \Pi(T)$ be an inclusion-minimal transversal of $H_T$.
We claim that $|X|\le k-1$.
Assume instead that $|X|\ge k$.
Choose distinct points $x_1,\dots,x_k\in X$.
By Lemma~\ref{lem:private-edge}, for each $i$ there exists an edge $E_i$ such that
\[
E_i\cap X=\{x_i\}.
\]
Each $E_i$ comes from some proper support piece $\sigma_i$.
The edges are distinct because their intersections with $X$ are distinct, so the support pieces $\sigma_i$ are distinct.
Setting
\[
Q:=\{x_1,\dots,x_k\},
\]
we obtain
\[
E_{\sigma_i}\cap Q=(E_{\sigma_i}\cap X)\cap Q=\{x_i\}.
\]
Thus $Q$ is a $k$-singleton trace, contradicting the assumption that no such trace exists.
Therefore $|X|\le k-1$.

Now average the weighted degrees over $X$:
\[
\sum_{\rho\in X} d_T(\rho)
=
\sum_{\rho\in X}\sum_{\sigma:\,\rho\in E_\sigma}\mu_T(\sigma)
=
\sum_\sigma \mu_T(\sigma)\,|E_\sigma\cap X|.
\]
Because $X$ is a transversal, every edge meets $X$, so $|E_\sigma\cap X|\ge 1$ for every proper support piece $\sigma$.
Hence
\[
\sum_{\rho\in X} d_T(\rho)\ge \sum_\sigma \mu_T(\sigma)=1.
\]
Therefore some $\rho\in X$ satisfies
\[
d_T(\rho)\ge \frac{1}{|X|}\ge \frac{1}{k-1}.
\]
\end{proof}

\begin{corollary}[Conditional sharp coherent-prefix theorem]\label{cor:singleton-sharp-coherence}
Assume the residual functoriality identity of Proposition~\ref{prop:vsf-to-singleton} and Conjecture~\ref{conj:visible-split-forcing}.
Then every neutral terminal kernel has a prefix $\rho\in \Pi(T)$ with
\[
d_T(\rho)\ge \frac{1}{k-1}.
\]
In particular, if every prefix of terminal edge-degree at least $1/(k-1)$ is already $(1/(k-1))$-coherent in the sense of Definition~\ref{def:beta-coherent-prefix}, then every neutral terminal has a profitable prefix.
\end{corollary}

\begin{proof}
The first claim follows by combining Proposition~\ref{prop:vsf-to-singleton} with Theorem~\ref{thm:singleton-heavy-prefix}.
If the displayed degree threshold implies $(1/(k-1))$-coherence, then Proposition~\ref{prop:coherent-lift} turns that prefix into a profitable parent prefix.
\end{proof}

\begin{conjecture}[Obstruction-witness realization at the first unresolved descendant]\label{conj:owr}
Let $u$ be a conditioned node of a neutral terminal, let $P$ be a proper visible support piece at $u$, and fix a profitable descendant trace inside a profitable $P$-child.
Assume that no lift of that trace is already a profitable prefix at $u$.
Choose a lexicographically minimal width-$(k-1)$ obstruction to lifting this trace, and project it by exact cone-state recursion to the first unresolved descendant $v$.
If the projected obstruction remains mixed across the projected separator $P_v$, then there exist lifted child witnesses at $v$, one seeing $P_v$ and one missing $P_v$.
\end{conjecture}

\begin{proposition}[Obstruction-witness realization implies visible split forcing]\label{prop:owr-implies-vsf}
Assume Conjecture~\ref{conj:owr}.
Then Conjecture~\ref{conj:visible-split-forcing} holds.
\end{proposition}

\begin{proof}
Suppose Conjecture~\ref{conj:visible-split-forcing} fails.
Choose a counterexample $u$ of minimal conditioned depth.
So $u$ is a conditioned node of a neutral terminal, there are lifted child witnesses with different visible next-step supports at $u$, but $u$ has no profitable prefix.

Choose a proper visible support piece $P$ that is seen by one witness and missed by another.
By Theorem~\ref{thm:neutral-terminal-child}, the support piece $P$ has a profitable child.
Fix a profitable descendant trace inside that child.
Because $u$ has no profitable prefix, no lift of this trace can already be profitable at $u$.
Choose a lexicographically minimal width-$(k-1)$ obstruction to lifting the trace.

This obstruction cannot lie entirely on the $P$-side, because then exact recursion would place the obstruction wholly inside the chosen profitable child.
It cannot lie entirely off the $P$-side either, because then it would not interact with the chosen child trace at all.
So the obstruction is mixed across the separator $P$.

Project the obstruction by exact cone-state recursion until the first unresolved descendant $v$.
Minimality of $v$ ensures that the projected obstruction is still mixed across the projected separator $P_v$.
Conjecture~\ref{conj:owr} therefore produces lifted child witnesses at $v$, one seeing $P_v$ and one missing $P_v$.
Thus $v$ is a strictly smaller counterexample to visible split forcing, contradicting the minimal choice of $u$.
Hence visible split forcing must hold.
\end{proof}

\begin{remark}[Proof audit for the separator route]\label{rem:separator-route-audit}
A tempting but invalid shortcut is the implication
\[
\text{``a profitable child on a separator piece }P\text{''}
\Longrightarrow
\text{``a profitable parent prefix''}.
\]
The full width-$(k-1)$ obstruction family can remain mixed across the separator, and pointwise child-profit does not by itself re-realize that mixed obstruction as an actual visible child disagreement downstairs.
Conjecture~\ref{conj:owr} isolates exactly the missing local statement needed to repair that gap.
\end{remark}





\subsection*{Visible boundary classes, separator locks, and primitive common-ray circuits}
The primitive-core reduction can be sharpened further once one records the unique separator profile already visible in a full-support cancellation witness.

\begin{definition}[Visible boundary class and splice gap]\label{def:visible-boundary-splice}
Suppose an unresolved one-step obstruction is represented on a finite active core $C$ by:
\begin{enumerate}[label=(\roman*),leftmargin=2.2em]
\item a matrix $A_C$ recording the non-target separator defects of the active child rays;
\item a target defect row $b_C$;
\item a positive full-support cancellation vector
\[
\omega\in \R_{>0}^{C}
\qquad\text{with}\qquad
A_C\omega=0
\quad\text{and}\quad
b_C\cdot \omega>0.
\]
\end{enumerate}
For an exposed proper subset $P\subsetneq C$, write $P^c:=C\setminus P$ and define the \emph{visible boundary class}
\[
w_P:=A_P(\omega|_P)=A_{P^c}(\omega|_{P^c}).
\]
The \emph{splice gap} of $P$ at a boundary class $w$ is
\[
\Psi_P(w):=
\sup\Bigl\{
b_P\cdot x-b_{P^c}\cdot y:
x,y\ge 0,\ A_Px=A_{P^c}y=w
\Bigr\}.
\]
\end{definition}

\begin{proposition}[Positive splice gap gives a visible split]\label{prop:positive-splice-visible-split}
In the setting of Definition~\ref{def:visible-boundary-splice}, suppose some exposed proper piece $P\subsetneq C$ satisfies
\[
\Psi_P(w_P)>0.
\]
Then there exist nonnegative vectors $x$ on $P$ and $y$ on $P^c$ such that
\[
A_Px=A_{P^c}y=w_P
\qquad\text{and}\qquad
b_P\cdot x>b_{P^c}\cdot y.
\]
Consequently the unresolved witness admits a separator-preserving profitable one-step split across $P$.
\end{proposition}

\begin{proof}
By definition of the supremum, the inequality $\Psi_P(w_P)>0$ yields feasible nonnegative $x,y$ with common boundary class $w_P$ and strictly larger target value on the $P$-side than on the $P^c$-side. The vector $x-y$ therefore preserves all non-target separator coordinates and strictly improves the target coordinate. This is exactly a separator-preserving profitable one-step split across $P$.
\end{proof}

\begin{definition}[Separator-locked primitive exposed core]\label{def:separator-locked-core}
In the situation of Definition~\ref{def:visible-boundary-splice}, the active core $C$ is called \emph{separator-locked} if
\[
\Psi_P(w_P)\le 0
\qquad\text{for every exposed proper piece }P\subsetneq C.
\]
It is called \emph{primitive} if every active ray is essential, equivalently if no proper active exposed sub-support carries the same unresolved failure.
\end{definition}

\begin{theorem}[Failure of local uncrossing compresses to a separator-locked primitive exposed core]\label{thm:separator-locked-core-reduction}
Suppose an unresolved one-step obstruction has already been compressed to a canonical exposed simplicial core in the sense of Proposition~\ref{prop:primitive-core-reduction}, and suppose that on this core the active child data admit the split into non-target separator coordinates $A_C$, target coordinate $b_C$, and full-support cancellation vector $\omega$ from Definition~\ref{def:visible-boundary-splice}. Then exactly one of the following holds:
\begin{enumerate}[label=(\roman*),leftmargin=2.2em]
\item some exposed proper piece $P\subsetneq C$ has $\Psi_P(w_P)>0$, so a separator-preserving visible split exists by Proposition~\ref{prop:positive-splice-visible-split};
\item the obstruction is supported on a separator-locked primitive exposed core.
\end{enumerate}
\end{theorem}

\begin{proof}
If alternative~(i) holds there is nothing more to prove.

Assume instead that $\Psi_P(w_P)\le 0$ for every exposed proper piece $P\subsetneq C$. Then $C$ is separator-locked by definition. The underlying exposed simplicial witness is already primitive after deleting redundant active rays: if a proper active exposed sub-support still carried the same unresolved obstruction, the chosen compressed core would not have been support-minimal. Thus $C$ is a separator-locked primitive exposed core.
\end{proof}

\begin{proposition}[Explicit augmented-circuit form of a separator lock]\label{prop:separator-lock-augmented-circuit}
Assume the hypotheses of Theorem~\ref{thm:separator-locked-core-reduction}, and suppose furthermore that:
\begin{enumerate}[label=(\roman*),leftmargin=2.2em]
\item for some exposed proper piece $P\subsetneq C$, the splice problem at $w_P$ attains its supremum and
\[
\Psi_P(w_P)=0;
\]
\item the active supports on the two sides of $P$ are simplicial.
\end{enumerate}
Write
\[
\widetilde a_e:=(A_e,b_e),
\qquad
\widetilde A:=\binom{A_C}{b_C}.
\]
Then there exists a primitive signed affine circuit
\[
z=z^+-z^-\in \ker \widetilde A
\]
such that
\[
\supp(z^+)\subseteq P,\qquad
\supp(z^-)\subseteq P^c,\qquad
z^\pm>0,
\]
and the common augmented image
\[
\widetilde A z^+=\widetilde A z^-=(w_P,\theta)
\]
spans a common exposed ray of the two simplicial side cones
\[
K^+:=\operatorname{cone}\{\widetilde a_e:e\in \supp(z^+)\},
\qquad
K^-:=\operatorname{cone}\{\widetilde a_e:e\in \supp(z^-)\}.
\]
\end{proposition}

\begin{proof}
Because the splice supremum at $w_P$ is attained and equals $0$, there exist nonnegative vectors $x$ on $P$ and $y$ on $P^c$ such that
\[
A_Px=A_{P^c}y=w_P,
\qquad
b_P\cdot x=b_{P^c}\cdot y=:\theta.
\]
Equivalently,
\[
\widetilde A_Px=\widetilde A_{P^c}y=(w_P,\theta).
\]
Choose such a pair $(x,y)$ with inclusion-minimal active support. Then
\[
z:=x-y
\]
is a nonzero vector in $\ker \widetilde A$ with positive part supported on $P$ and negative part supported on $P^c$.

Support-minimality makes $z$ primitive as a signed affine dependence: if a proper sub-support carried another nonzero signed dependence, then the same common augmented point would already be realized on a smaller active support, contradicting the choice of $(x,y)$. Thus $z$ is a primitive signed affine circuit.

Finally, the point $(w_P,\theta)$ lies in both $K^+$ and $K^-$. If the intersection contained more than one ray, an extreme ray of the intersection would lie on a proper face of at least one simplicial side cone, hence on a strict sub-support of the circuit. That would again contradict primitivity. Therefore
\[
K^+\cap K^-=\R_{\ge 0}(w_P,\theta),
\]
so the lock is a one-dimensional common-ray circuit.
\end{proof}

\begin{corollary}[The first admissible separator lock is a balanced $2{:}2$ circuit]\label{cor:separator-lock-22}
Assume the setting of Proposition~\ref{prop:separator-lock-augmented-circuit}, and assume moreover that no single child already carries the one-step defect. Then every separator-locked primitive circuit has at least two positive and two negative coordinates. In particular, the smallest admissible lock is a balanced $2{:}2$ circuit
\[
\alpha_1\widetilde a_1+\alpha_2\widetilde a_2
=
\beta_1\widetilde a_3+\beta_2\widetilde a_4,
\qquad
\alpha_i,\beta_j>0.
\]
\end{corollary}

\begin{proof}
If the positive support had size $1$, then one active child on the $P$-side would realize the entire positive part of the circuit, contradicting the assumption that no single child already carries the one-step defect. The same argument applies to the negative support. Hence both signs occur at least twice, and the smallest possible primitive lock is of type $2{:}2$.
\end{proof}

\begin{remark}[Audit of the explicit circuit reduction]\label{rem:separator-lock-audit}
The preceding proposition is completely rigorous once the local splice problem is known to attain on the compressed core and the active sides are simplicial. What is \emph{not} yet proved elsewhere in the document is the structural theorem guaranteeing that every remaining lock enters this common-ray normal form. Thus the argument isolates the local obstruction exactly, but it does not yet eliminate that obstruction.
\end{remark}

\begin{proposition}[Quadrilateral discriminant under common-ray star reduction]\label{prop:quadrilateral-discriminant}
Assume a separator-locked primitive common-ray circuit has been reduced to a two-dimensional local star in which the only remaining primitive core is the alternating quadrilateral
\[
Q_4:\qquad u_1+u_3=u_2+u_4,
\]
with normalized circuit vector
\[
w=e_1-e_2+e_3-e_4.
\]
Assume furthermore that all visible target data lie on a common oriented ray $\R_{\ge 0}\rho$, so that
\[
b_i=\beta_i\rho
\qquad (1\le i\le 4),
\]
and that the lock relation is
\[
\beta_1+\beta_3=\beta_2+\beta_4.
\]
Write
\[
\beta_1=m+s,\qquad \beta_3=m-s,\qquad
\beta_2=m+t,\qquad \beta_4=m-t
\]
and
\[
\sigma_{13}:=b_1-b_3=2s\,\rho,
\qquad
\sigma_{24}:=b_2-b_4=2t\,\rho.
\]
Then the quadrilateral is support-minimal primitive if and only if
\[
s\neq 0,\qquad t\neq 0,\qquad s\neq t,\qquad s\neq -t.
\]
Equivalently,
\[
\sigma_{13}\neq 0,\qquad \sigma_{24}\neq 0,\qquad
\sigma_{13}\neq \sigma_{24},\qquad
\sigma_{13}\neq -\sigma_{24}.
\]
\end{proposition}

\begin{proof}
Because $Q_4$ is a convex quadrilateral, no $3$-point subset is an affine circuit: the only minimal $3$-point affine circuit would be a collinear triple, which cannot occur in a convex quadrilateral. Thus any proper-support counterwitness must have support $2$.

On the common ray, a support-$2$ invisible defect occurs exactly when some pair $b_i=b_j$, equivalently $\beta_i=\beta_j$. The displayed parametrization gives
\[
\beta_1=\beta_3 \iff s=0 \iff \sigma_{13}=0,
\]
\[
\beta_2=\beta_4 \iff t=0 \iff \sigma_{24}=0,
\]
\[
\beta_1=\beta_2 \ \text{and}\ \beta_3=\beta_4
\iff s=t \iff \sigma_{13}=\sigma_{24},
\]
and
\[
\beta_1=\beta_4 \ \text{and}\ \beta_2=\beta_3
\iff s=-t \iff \sigma_{13}=-\sigma_{24}.
\]

Each of these equalities destroys primitiveness. If $s=0$, then $b_1=b_3$, so $e_1-e_3$ is already invisible. If $t=0$, then $b_2=b_4$, so $e_2-e_4$ is already invisible. If $s=t$, then $b_1=b_2$ and $b_3=b_4$, and
\[
w=(e_1-e_2)+(e_3-e_4)
\]
splits into two support-$2$ invisible differences. If $s=-t$, then $b_1=b_4$ and $b_2=b_3$, and
\[
w=(e_1-e_4)+(e_3-e_2)
\]
again splits into two support-$2$ invisible differences.

Conversely, if none of
\[
s=0,\qquad t=0,\qquad s=t,\qquad s=-t
\]
holds, then no pair $b_i=b_j$ exists, so there is no support-$2$ invisible defect. Since no support-$3$ affine circuit exists in a convex quadrilateral, the $4$-point circuit is support-minimal primitive.
\end{proof}

\begin{remark}[What remains in the separator lane]\label{rem:quadrilateral-what-remains}
Proposition~\ref{prop:quadrilateral-discriminant} isolates the exact local discriminant
\[
\Delta_{Q_4}:=s\,t\,(s^2-t^2).
\]
Once a separator lock has been reduced to the common-ray quadrilateral normal form, the only remaining task is to show that every geometric instance satisfies $\Delta_{Q_4}=0$. On the genuinely full-support branch this reduces to the slope-lock identity
\[
\sigma_{13}=\pm \sigma_{24}.
\]
Thus the separator route has been reduced from a diffuse uncrossing problem to the elimination of one explicit common-ray quadrilateral lock.
\end{remark}

\subsection*{Separator-compatible lifting as the earliest unresolved theorem}
The separator route also yields a clean proof-dag statement: once the already-rigorous reductions are accepted, the first genuinely missing theorem is a local lifting theorem on bounded exposed witnesses.

\begin{definition}[Local separator theorem $L_{\mathrm{sep}}(k)$]\label{def:lsep}
For fixed $k$, let $L_{\mathrm{sep}}(k)$ denote the statement that every bounded exposed witness, equivalently every bounded signed affine circuit, on a trace-atomic low-coherence neutral terminal either
\begin{enumerate}[label=(\roman*),leftmargin=2.2em]
\item admits separator-compatible lifting to the support side, or
\item forces a visible split.
\end{enumerate}
\end{definition}

\begin{theorem}[Conditional closure from $L_{\mathrm{sep}}(k)$]\label{thm:lsep-implies-exp}
Fix $k$. Assume:
\begin{enumerate}[label=(\roman*),leftmargin=2.2em]
\item the repaired low-excess extraction package together with descendant re-entry;
\item the reduction to trace-atomic neutral leaves;
\item the low-coherence threshold theorem for neutral terminals;
\item the exact convex/Farkas lift alternative and the exposed-core compression to bounded signed affine circuits;
\item the already-recorded finite base-window reductions of the continuation-cone, OA, and $(5,5)$ lanes to bounded exposed witnesses.
\end{enumerate}
If $L_{\mathrm{sep}}(k)$ holds, then there exists $c_k<\infty$ such that
\[
f(n,k)<c_k^n
\qquad\text{for all }n.
\]
\end{theorem}

\begin{proof}
Assume toward contradiction that the exponential sunflower bound fails for this fixed $k$. By the repaired low-excess extraction and descendant re-entry, the obstruction reduces, up to an $e^{o(n)}$ loss, to irreducible neutral leaves. By the trace-atomic reduction, one may take a minimal bad leaf to be trace-atomic.

The coherence-threshold theorem excludes the high-coherence regime: if a surviving bad leaf had coherence at least the threshold $\beta$, it would already carry a profitable parent prefix. Hence every remaining bad leaf is low-coherence.

Now inspect assembly from proper support pieces. Every proper support piece of a neutral terminal has a profitable child, so the only way a low-coherence trace-atomic leaf can survive is that these profitable children fail to assemble to a profitable parent. The exact convex/Farkas alternative then produces a strict-gap dual witness, and the exposed-core compression reduces that witness to a bounded exposed simplicial core, equivalently a bounded signed affine circuit.

By the finite base-window reductions in the continuation-cone, OA, and Hall lanes, every remaining universal obstruction already lives on such a bounded exposed witness. Apply $L_{\mathrm{sep}}(k)$ to that witness.

If $L_{\mathrm{sep}}(k)$ yields a visible split, then the witness exposes a genuine proper support split. Each side has a profitable child, and the exact convex lift theorem assembles those child improvements into a profitable parent, contradicting neutrality.

If $L_{\mathrm{sep}}(k)$ yields separator-compatible lifting, then the exposed separator data are realized on the support side. The same exact support-piece bookkeeping again produces profitable children on the resulting proper pieces, and the exact convex/Farkas alternative assembles them to a profitable parent, again contradicting neutrality.

So no trace-atomic low-coherence neutral leaf can survive. Undoing the $e^{o(n)}$ loss from the repaired extraction shows that the original bad family cannot exist. Hence $f(n,k)<c_k^n$ for some constant $c_k$.
\end{proof}

\begin{definition}[Residual low-coherence circuit class]\label{def:residual-class}
Fix $k$. Let $\mathcal C_k(n)$ be the class of residual-rank-$n$ objects $T$ such that:
\begin{enumerate}[label=(\roman*),leftmargin=2.2em]
\item $T$ is an irreducible neutral leaf;
\item $T$ is trace-atomic;
\item $T$ is low-coherence;
\item the weighted child-to-parent lift fails on $T$;
\item the failure is represented by a support-minimal canonical exposed simplicial core, equivalently by a bounded signed affine circuit;
\item $T$ is not already in a solved trace, support-box, fixed-memory, continuation-cone, OA, or local diagonal class handled elsewhere in this document.
\end{enumerate}
\end{definition}

\begin{proposition}[Residual-class reformulation of the global problem]\label{prop:residual-class-equivalence}
Assume the same reduction package as in Theorem~\ref{thm:lsep-implies-exp}. Then, up to the $e^{o(n)}$ loss in the repaired low-excess extraction, the fixed-$k$ exponential sunflower conjecture is equivalent to the eventual emptiness of the residual class $\mathcal C_k(n)$.
\end{proposition}

\begin{proof}
Take a minimal bad family. The repaired extraction reduces it to an irreducible neutral leaf, the trace-atomic reduction sharpens it to a trace-atomic leaf, and the coherence threshold forces low coherence. Failure of the weighted child-to-parent lift yields a dual witness, which the exposed-core compression turns into a bounded signed affine circuit. Any leaf already entering a solved trace, support-box, fixed-memory, continuation-cone, OA, or local diagonal class would be handled by the corresponding theorem or reduction recorded earlier. Hence every genuine survivor lies in $\mathcal C_k(n)$.

Conversely, if $\mathcal C_k(n)$ is empty for all sufficiently large $n$, then after the $e^{o(n)}$ reduction there are no large-rank bad leaves. The finitely many small ranks are absorbed into the base of an exponential bound. So eventual emptiness of $\mathcal C_k(n)$ is equivalent to the conjectured exponential bound.
\end{proof}

\begin{remark}[The earliest non-derivable separator sentence]\label{rem:lsep-earliest-gap}
The current conditional package does \emph{not} already imply the full conjecture. The earliest unresolved sentence is exactly $L_{\mathrm{sep}}(k)$. A stronger statement,
\[
\text{``the canonical exposed core is realized on an actual support subset in the same boundary metric''},
\]
would certainly imply $L_{\mathrm{sep}}(k)$, but it is not earlier in the proof DAG: it contains both separator lifting and an additional rigidity step that would reconstruct the whole exposed core on the support side. The latter is precisely the sort of inverse/local-coordinate statement ruled out by the counterexamples recorded in Section~\ref{sec:counterexamples}. Thus the first honest unresolved theorem is separator-compatible lifting or visible split forcing, not full same-boundary realization.
\end{remark}

\status{The local separator notes now separate three layers cleanly. The splice-gap and quadrilateral-discriminant statements are rigorous once the local model has been reduced to the stated common-ray normal forms. The global proof still stops at the same place as before: the earliest unresolved theorem is $L_{\mathrm{sep}}(k)$, namely separator-compatible lifting or visible split forcing on bounded exposed witnesses.}


\subsection*{Cone states, exact links, and the real continuation-cone bottleneck}
The continuation-cone route becomes cleaner once one distinguishes genuine cone states from raw exact-support classes.

\begin{definition}[Cone-reduced residual trace state]\label{def:cone-state}
Fix an active coordinate set $P$.
A \emph{cone-reduced residual trace state} on $P$ is an up-set
\[
T\subseteq 2^P
\]
of the Boolean lattice on $P$.
Its set of \emph{minimal continuation cones} is
\[
\mathcal C:=\min(T),
\]
and one has
\[
T=\uparrow \mathcal C:=\{X\subseteq P:\exists\, C\in \mathcal C \text{ with } C\subseteq X\}.
\]
The \emph{cone rank} is
\[
r(T):=\max_{C\in \mathcal C} |C|.
\]
\end{definition}

\begin{proposition}[Projection formula for exact conditioned cone children]\label{prop:cone-child-projection}
Let $T\subseteq 2^P$ be a cone-reduced residual trace state with minimal cone family $\mathcal C=\min(T)$.
Let $H\subseteq P$ be a transversal of $\mathcal C$, and let $S\subseteq H$.
Define the exact $H$-conditioned child state on $P\setminus H$ by
\[
T_S^{-H}:=\{R\subseteq P\setminus H : R\cup S\in T\},
\qquad
\mathcal C^+:=\min(T_S^{-H}).
\]
Then
\[
\mathcal C^+
=
\min\{\, C\setminus H : C\in \mathcal C,\ C\cap H\subseteq S\,\}.
\]
\end{proposition}

\begin{proof}
Let $R\subseteq P\setminus H$.
By definition,
\[
R\in T_S^{-H}
\iff
R\cup S\in T.
\]
Because $T=\uparrow \mathcal C$, this is equivalent to the existence of some $C\in \mathcal C$ with
\[
C\subseteq R\cup S.
\]
Since $R$ is disjoint from $H$, the containment $C\subseteq R\cup S$ is equivalent to the two conditions
\[
C\cap H\subseteq S
\qquad\text{and}\qquad
C\setminus H\subseteq R.
\]
Therefore
\[
T_S^{-H}
=
\uparrow\{\, C\setminus H : C\in \mathcal C,\ C\cap H\subseteq S\,\}.
\]
Taking minimal elements on both sides gives exactly the stated formula for $\mathcal C^+$.
\end{proof}

\begin{corollary}[Strict rank-drop for transversal conditioning of cone states]\label{cor:cone-rank-drop}
With the notation of Proposition~\ref{prop:cone-child-projection}, suppose that $H$ is a transversal of $\mathcal C$.
Then
\[
r(T_S^{-H})\le r(T)-1
\qquad\text{for every }S\subseteq H.
\]
\end{corollary}

\begin{proof}
Because $H$ is a transversal, every $C\in \mathcal C$ satisfies $C\cap H\neq \varnothing$.
Hence
\[
|C\setminus H|\le |C|-1\le r(T)-1
\qquad\text{for every }C\in \mathcal C.
\]
By Proposition~\ref{prop:cone-child-projection}, every child minimal cone belongs to the minimalization of the family $\{C\setminus H: C\in \mathcal C,\ C\cap H\subseteq S\}$.
So each child minimal cone has size at most $r(T)-1$, proving the claim.
\end{proof}

\begin{counterexample}[Raw exact-support conditioning need not drop rank]\label{ce:raw-cone-rank}
The strict rank-drop conclusion of Corollary~\ref{cor:cone-rank-drop} fails if one first conditions the \emph{raw exact-support family} and only afterwards cone-reduces.
\end{counterexample}

\begin{proof}
Let
\[
P=\{h_1,h_2,x\},
\qquad
T_{\mathrm{raw}}=\bigl\{\{h_1\},\{h_2\},\{h_1,h_2,x\}\bigr\}.
\]
The inclusion-minimal exact supports are
\[
\mathcal C_{\mathrm{raw}}=\bigl\{\{h_1\},\{h_2\}\bigr\},
\]
so the raw minimal-support rank is $1$.
Take
\[
H=\{h_1,h_2\},
\qquad
S=H.
\]
After conditioning on the exact $H$-pattern $S$, the only surviving raw support is $\{h_1,h_2,x\}$.
Deleting $H$ leaves the child family $\{\{x\}\}$, whose minimal-support rank is still $1$.
Thus the strict drop fails at the raw exact-support level.

The point is that $\{h_1,h_2,x\}$ was not minimal before conditioning, but becomes minimal afterwards.
Proposition~\ref{prop:cone-child-projection} avoids this problem because a genuine cone state is the up-set generated by its minimal cones from the start.
\end{proof}

\begin{remark}[Audit correction for the continuation-cone route]\label{rem:cone-route-audit}
An earlier audit of the continuation-cone route identified the possibility that new minimal cones could appear after branching and spoil strict rank-drop.
Counterexample~\ref{ce:raw-cone-rank} shows that this concern is genuine for raw exact-support families.
Corollary~\ref{cor:cone-rank-drop} shows, however, that the issue disappears for genuine cone states.
So the first real obstruction in the continuation-cone lane is not the rank-drop step.
It is the positive-mass bounded-transversal capture problem described below.
\end{remark}

\begin{definition}[Exact link and transversal number]\label{def:exact-link-transversal}
Let $C$ be a continuation cone equipped with normalized counting measure $\mu_C$.
For an atom $x$, define the \emph{exact link} $L_x(C)$ to be the cone obtained by restricting to those members of $C$ that contain $x$, deleting $x$, and cone-reducing.
The \emph{transversal number} of a family $\mathcal A$ is
\[
\tau(\mathcal A):=\min\{|T|:T\cap A\neq \varnothing\ \text{for every }A\in \mathcal A\}.
\]
\end{definition}

\begin{proposition}[Heavy exact link from a bounded-transversal subcone]\label{prop:heavy-exact-link}
Let $C$ be a continuation cone with normalized counting measure $\mu_C$.
Suppose $C'\subseteq C$ is an exact subcone such that
\[
\mu_C(C')\ge \delta
\qquad\text{and}\qquad
\tau(C')\le t.
\]
Then there exists an atom $x$ such that
\[
\mu_C\bigl(L_x(C')\bigr)\ge \frac{\delta}{t}.
\]
Moreover any $k$-sunflower in $L_x(C')$ would lift, after reattaching $x$, to a $k$-sunflower in $C'$.
\end{proposition}

\begin{proof}
Choose a transversal $T$ of $C'$ with $|T|\le t$.
Every member of $C'$ contains at least one atom from $T$, so
\[
\sum_{x\in T} \mu_C(\{A\in C':x\in A\})\ge \mu_C(C')\ge \delta.
\]
Hence some $x\in T$ satisfies
\[
\mu_C(\{A\in C':x\in A\})\ge \frac{\delta}{t}.
\]
Passing from these sets to the exact link $L_x(C')$ just deletes the common atom $x$ and cone-reduces, so the same lower bound holds for $\mu_C(L_x(C'))$.

If $L_x(C')$ contained a $k$-sunflower, then reattaching $x$ to every petal would produce a $k$-sunflower in $C'$.
Therefore $L_x(C')$ is a legitimate recursive child in the sunflower-free setting.
\end{proof}

\begin{proposition}[Bounded active rank gives bounded transversal]\label{prop:bounded-rank-transversal}
Let $C$ be a continuation cone all of whose members have size at most $r$.
If $C$ contains no $k$-matching, then
\[
\tau(C)\le r(k-1).
\]
\end{proposition}

\begin{proof}
Take a maximal matching $M\subseteq C$.
Because $C$ has no $k$-matching, one has $|M|\le k-1$.
The union
\[
X:=\bigcup_{A\in M} A
\]
meets every member of $C$, else some disjoint member could be added to the matching.
Hence $X$ is a transversal.
Since every matched set has size at most $r$,
\[
|X|\le r|M|\le r(k-1).
\]
Therefore $\tau(C)\le r(k-1)$.
\end{proof}

\begin{theorem}[Positive-mass bounded-transversal capture would force an exponential bound]\label{thm:bounded-transversal-capture}
Fix $k\ge 3$.
Assume there exist constants $\delta_k>0$ and $t_k<\infty$ with the following property:
every unsolved continuation cone $C$ contains an exact subcone $C'\subseteq C$ such that
\[
\mu_C(C')\ge \delta_k
\qquad\text{and}\qquad
\tau(C')\le t_k.
\]
Then there exists $c_k<\infty$ such that
\[
f(n,k)<c_k^n
\qquad\text{for all }n.
\]
More precisely, if $h(m,k)$ denotes the extremal size in residual rank $m$, then
\[
h(m,k)\le \left(\frac{t_k}{\delta_k}\right)^m
\]
up to the fixed base cases of the solved classes.
\end{theorem}

\begin{proof}
Let $C$ be an unsolved continuation cone of residual rank $m\ge 1$, and equip it with normalized counting measure.
By hypothesis, $C$ contains an exact subcone $C'$ with $\mu_C(C')\ge \delta_k$ and $\tau(C')\le t_k$.
Proposition~\ref{prop:heavy-exact-link} gives an atom $x$ with
\[
\mu_C\bigl(L_x(C')\bigr)\ge \frac{\delta_k}{t_k}.
\]
Equivalently,
\[
|C|\le \frac{t_k}{\delta_k}\,|L_x(C')|.
\]
The exact link $L_x(C')$ has residual rank at most $m-1$ because the common atom $x$ has been removed.
Applying the same argument recursively yields
\[
h(m,k)\le \frac{t_k}{\delta_k}\,h(m-1,k)
\le
\left(\frac{t_k}{\delta_k}\right)^m h(0,k).
\]
Absorbing the solved base cases into the constant shows that $h(m,k)\le c_k^m$ for some $c_k$ depending only on $k$.
Finally, Theorem~\ref{thm:exact-reduction} converts this residual-rank exponential bound into the ordinary exponential bound for $f(n,k)$.
\end{proof}

\begin{counterexample}[No-matching does not imply bounded transversal in general]\label{ce:no-matching-no-transversal}
For each $m\ge 1$, the uniform family
\[
\binom{[2m-1]}{m}
\]
is intersecting, so it has matching number $1$, but its transversal number is $m$.
Thus the step
\[
\text{``no }k\text{-matching''}\Longrightarrow \text{``bounded transversal''}
\]
is false without extra structure.
\end{counterexample}

\begin{proof}
Two $m$-subsets of $[2m-1]$ cannot be disjoint, because their union would then have size $2m$, larger than $2m-1$.
So the family is intersecting.

On the other hand, any set $T\subseteq [2m-1]$ of size at most $m-1$ fails to hit some $m$-subset:
simply take an $m$-subset of the complement $[2m-1]\setminus T$, which has size at least
\[
(2m-1)-(m-1)=m.
\]
Therefore every transversal has size at least $m$.
Since any fixed $m$-subset itself is a transversal, the transversal number is exactly $m$.
\end{proof}

\begin{remark}[The real local gap in the continuation-cone route]\label{rem:bounded-transversal-gap}
Theorem~\ref{thm:bounded-transversal-capture} isolates the actual missing local statement for the continuation-cone lane:
one needs a \emph{positive-mass} theorem placing a definite fraction of every unsolved cone inside a bounded-transversal, or bounded-active-rank, regime.
The heavy-link step itself is completely rigorous once such a regime is reached.
\end{remark}

\begin{theorem}[First-hit profiled solved cover]\label{thm:first-hit-profiled-cover}
Let $X$ be an unsolved continuation-cone state on a finite ground set $P$ in the sense of Definition~\ref{def:cone-state}.
Choose a minimal continuation cone
\[
U_0\in \min(X)
\]
and a transversal $T\subseteq P$ of $U_0$.
Then there exists a finite hereditary profiled decomposition tree $\mathfrak D(X,T)$ with the following properties:
\begin{enumerate}[label=(\roman*),leftmargin=2.2em]
\item each node is a pair $(\pi,A)$, where $\pi$ is an exact support-pattern prefix in the terminal exact-trace tree and $A\subseteq T$ is the active transversal profile;
\item the root is $(\varnothing,T)$;
\item if $U_\pi$ denotes the projected child cone of $U_0$ under the exact prefix $\pi$, then for every internal node $(\pi,A)$ the set $A$ is a transversal of $U_\pi$;
\item when $(\pi,A)$ branches on the next exact support-pattern child $\sigma$ carried by a trace piece $P_\sigma$, there are two cases:
\begin{enumerate}[label=(\alph*),leftmargin=2.2em]
\item if $\sigma\cap A=\varnothing$, the child is $(\pi\sigma,A\setminus P_\sigma)$ and the new active profile is again a transversal of $U_{\pi\sigma}$;
\item if $\sigma\cap A\neq \varnothing$, the branch stops at that child, and the corresponding projected child cone has strictly smaller cone rank than $U_0$.
\end{enumerate}
\item every exact support-pattern child of every internal node appears in the tree, so the cover is hereditary under exact support-pattern branching;
\item every leaf is either empty, already solved, or a strict rank-drop leaf.
\end{enumerate}
\end{theorem}

\begin{proof}
For every exact support-pattern prefix $\pi$, let $U_\pi$ be the projected child cone obtained from $U_0$ by exact recursion.
We build the tree recursively.

At the root $(\varnothing,T)$ the active profile $T$ is a transversal of $U_0$ by assumption.
Now assume inductively that $(\pi,A)$ is an internal node and that $A$ is a transversal of $U_\pi$.
Let $\sigma$ be an exact support-pattern child at the next step, supported on a trace piece $P_\sigma$.

First suppose that $\sigma\cap A=\varnothing$.
Define
\[
A':=A\setminus P_\sigma.
\]
We claim that $A'$ is a transversal of $U_{\pi\sigma}$.
Take any member $R\in U_{\pi\sigma}$.
By exact child-cone recursion, $R$ extends to some member $S\in U_\pi$ compatible with the support-pattern child $\sigma$.
Because $A$ is a transversal of $U_\pi$, the set $S$ meets $A$.
If $S\cap A$ were contained entirely in $P_\sigma$, then $\sigma$ would meet $A$, contrary to assumption.
Hence
\[
S\cap A\subseteq A\setminus P_\sigma=A',
\]
and therefore $R$ meets $A'$.
So $A'$ is indeed a transversal of $U_{\pi\sigma}$.

Now suppose that $\sigma\cap A\neq \varnothing$.
Choose any $t\in \sigma\cap A$.
Conditioning on $t$ removes one atom from every minimal cone of $U_\pi$ because $A$ is a transversal.
By Corollary~\ref{cor:cone-rank-drop}, the exact child obtained by conditioning on $t$ has strictly smaller cone rank than $U_\pi$, and hence strictly smaller cone rank than the root cone $U_0$.
So this branch may be terminated as a strict rank-drop leaf.

Because every exact support-pattern child is included at each internal node, the tree is hereditary by construction.
The finite exact trace tree ensures that the process terminates after finitely many steps.
Every leaf is either empty, already solved by the ambient recursion, or a first-hit leaf where strict rank-drop has occurred.
\end{proof}

\begin{remark}[Why the profiled cover repairs the failed greedy branch]\label{rem:first-hit-profiled-cover}
The counterexample to the unprofiled greedy heaviest-child branch shows that raw child selection is not hereditary finite-state.
Theorem~\ref{thm:first-hit-profiled-cover} identifies the missing state variable:
one must carry the active transversal profile $A\subseteq T$.
Once that profile is remembered, the exact branch recursion becomes hereditary and every first hit forces strict rank-drop.
\end{remark}

\begin{definition}[Bounded-transversal mass profile and witness pressure]\label{def:witness-pressure}
For a continuation cone $C$ with normalized counting measure $\mu_C$, define
\[
\kappa_t(C):=
\sup\bigl\{\mu_C(D): D\subseteq C\text{ is an exact subcone and }\tau(D)\le t\bigr\}.
\]
The \emph{witness pressure} of $C$ is
\[
w(C):=\inf_{t\ge 1}\left(t+\left\lceil \log_2\frac{1}{\kappa_t(C)}\right\rceil\right),
\]
with the convention $w(C)=\infty$ if every $\kappa_t(C)$ vanishes.
\end{definition}

\begin{lemma}[Hereditary pressure drop]\label{lem:hereditary-pressure-drop}
Let $D\subseteq C$ be an exact subcone with $\mu_C(D)=\eta>0$.
Then for every exact subcone $E\subseteq D$,
\[
\tau(E)+\left\lceil \log_2\frac{1}{\mu_D(E)}\right\rceil
>
w(C)-\left\lceil \log_2\frac{1}{\eta}\right\rceil-1.
\]
In particular,
\[
w(D)\ge w(C)-\left\lceil \log_2\frac{1}{\eta}\right\rceil-1.
\]
\end{lemma}

\begin{proof}
For every exact $E\subseteq D$ one has
\[
\mu_C(E)=\eta\,\mu_D(E).
\]
Therefore
\[
\left\lceil \log_2\frac{1}{\mu_C(E)}\right\rceil
\le
\left\lceil \log_2\frac{1}{\eta}\right\rceil
+
\left\lceil \log_2\frac{1}{\mu_D(E)}\right\rceil.
\]
If for some exact $E\subseteq D$ we had
\[
\tau(E)+\left\lceil \log_2\frac{1}{\mu_D(E)}\right\rceil
\le
w(C)-\left\lceil \log_2\frac{1}{\eta}\right\rceil-1,
\]
then
\[
\tau(E)+\left\lceil \log_2\frac{1}{\mu_C(E)}\right\rceil
\le
w(C)-1,
\]
contradicting the definition of $w(C)$ as the infimum over all exact subcones of $C$.
This proves the displayed inequality and hence the bound on $w(D)$.
\end{proof}

\begin{definition}[$(a,b)$-neutral continuation cone]\label{def:ab-neutral-cone}
Fix integers $a,b\ge 1$.
An unsolved continuation cone $C$ is called \emph{$(a,b)$-neutral} if every exact subcone $D\subseteq C$ with
\[
\mu_C(D)\ge 2^{-a}
\]
satisfies
\[
\kappa_b(D)<2^{-b}.
\]
\end{definition}

\begin{theorem}[Failure of a uniform witness-pressure bound produces neutral scales]\label{thm:witness-pressure-neutral-scales}
Assume that the witness pressure is not uniformly bounded on unsolved continuation cones.
Then for every pair of integers $a,b\ge 1$ there exists an unsolved $(a,b)$-neutral continuation cone.
\end{theorem}

\begin{proof}
Because the witness pressure is not uniformly bounded, for every integer $m$ there exists an unsolved continuation cone $C_m$ with
\[
w(C_m)>m.
\]
Fix $a,b\ge 1$, and choose $m>a+2b+1$.
We claim that $C_m$ is $(a,b)$-neutral.

Let $D\subseteq C_m$ be an exact subcone with
\[
\mu_{C_m}(D)\ge 2^{-a}.
\]
By Lemma~\ref{lem:hereditary-pressure-drop},
\[
w(D)\ge w(C_m)-a-1>2b.
\]
If $\kappa_b(D)\ge 2^{-b}$, then some exact subcone $E\subseteq D$ satisfies
\[
\tau(E)\le b
\qquad\text{and}\qquad
\mu_D(E)\ge 2^{-b}.
\]
Hence
\[
\tau(E)+\left\lceil \log_2\frac{1}{\mu_D(E)}\right\rceil
\le b+b=2b,
\]
which would imply $w(D)\le 2b$, a contradiction.
So every exact $D\subseteq C_m$ of mass at least $2^{-a}$ satisfies $\kappa_b(D)<2^{-b}$, exactly as required.
\end{proof}

\begin{theorem}[Eliminating one finite neutral scale yields a uniform pressure bound]\label{thm:one-scale-elimination}
Fix integers $a,b\ge 1$.
Assume that no unsolved continuation cone is $(a,b)$-neutral.
Then every unsolved continuation cone $C$ satisfies
\[
w(C)\le a+2b.
\]
Consequently, if the solved trace classes carry their proved exponential bounds, the continuation-cone route closes with an exponential sunflower bound.
\end{theorem}

\begin{proof}
Let $C$ be an unsolved continuation cone.
Because $C$ is not $(a,b)$-neutral, there exists an exact subcone $D\subseteq C$ with
\[
\mu_C(D)\ge 2^{-a}
\qquad\text{and}\qquad
\kappa_b(D)\ge 2^{-b}.
\]
Choose an exact subcone $E\subseteq D$ such that
\[
\tau(E)\le b
\qquad\text{and}\qquad
\mu_D(E)\ge 2^{-b}.
\]
Then
\[
\mu_C(E)=\mu_C(D)\mu_D(E)\ge 2^{-(a+b)}.
\]
Therefore
\[
w(C)\le \tau(E)+\left\lceil \log_2\frac{1}{\mu_C(E)}\right\rceil
\le
b+(a+b)=a+2b.
\]
This gives a uniform pressure bound on all unsolved continuation cones.
Once the solved classes are inserted with their certified exponential estimates, Theorem~\ref{thm:bounded-transversal-capture} and the exact reduction machinery convert this bound into an exponential sunflower estimate.
\end{proof}

\begin{remark}[What the witness-pressure reformulation isolates]\label{rem:witness-pressure-isolates-gap}
The witness-pressure language sharpens Theorem~\ref{thm:bounded-transversal-capture}.
The exact missing local statement is now:
eliminate one finite neutral scale.
Equivalently, prove that every unsolved continuation cone contains a positive-mass exact descendant carrying a bounded-transversal exact witness.
The heavy-link step and the recursive closure are already rigorous once that finite-scale theorem is available.
\end{remark}



\subsection*{Finite base-window concentration and BridgeLock}
The witness-pressure reformulation can be sharpened further: once the finite base window is reached, every remaining continuation-cone obstruction concentrates on a single local bridge class.

\begin{definition}[Finite base-window local-bridge package]\label{def:bridgelock-package}
Fix $k$. Assume that the continuation-cone lane comes equipped with:
\begin{enumerate}[label=(\roman*),leftmargin=2.2em]
\item a finite residue set $\{0,\dots,\sigma_k-1\}$;
\item classes $\mathsf U_k$ of unsolved continuation cones and $\mathsf B_{a,k}(b)\subseteq \mathsf U_k$ of residue-$a$ exact continuation cones carrying witness pressure at least $b$;
\item canonical base-window collapse maps
\[
\Phi_a:\mathsf B_{a,k}\longrightarrow \mathsf{LB}_k
\]
into a common local-bridge class $\mathsf{LB}_k$;
\item constants $\lambda_{a,k}=O_k(1)$ such that every exact witness reflects across the collapse with
\[
w(C)\le w(\Phi_a(C))+\lambda_{a,k}
\qquad\text{for every }C\in \mathsf B_{a,k}(b).
\]
\end{enumerate}
\end{definition}

\begin{proposition}[Unbounded witness pressure concentrates on the local-bridge class]\label{prop:bridgelock-concentration}
Assume the finite base-window local-bridge package of Definition~\ref{def:bridgelock-package}. If the witness pressure is not uniformly bounded on $\mathsf U_k$, then it is not uniformly bounded on the common local-bridge class $\mathsf{LB}_k$.
\end{proposition}

\begin{proof}
Assume that $w$ is not uniformly bounded on $\mathsf U_k$. Then for every integer $b$ there exists an exact continuation cone
\[
C_b\in \mathsf U_k
\]
with witness pressure at least $b$.

By the finite base-window collapse, each such cone belongs to some residue class $a_b\in\{0,\dots,\sigma_k-1\}$, and after collapsing we obtain
\[
L_b:=\Phi_{a_b}(C_b)\in \mathsf{LB}_k.
\]
Because only finitely many residues occur, one residue $a_*$ appears infinitely often. Restrict to a sequence $b_j\to\infty$ with
\[
C_j\in \mathsf B_{a_*,k}(b_j).
\]
The reflection inequality gives
\[
w(L_j)\ge w(C_j)-\lambda_{a_*,k}\ge b_j-\lambda_{a_*,k}\to\infty.
\]
Hence witness pressure is unbounded on $\mathsf{LB}_k$.
\end{proof}

\begin{definition}[BridgeLock theorem]\label{def:bridgelock}
For fixed $k$, write $\mathbf{BridgeLock}_k$ for the assertion that there exists $R_k$ such that every object $L\in \mathsf{LB}_k$ contains an exact subcone $E\subseteq L$ with
\[
\mu_L(E)\ge 2^{-R_k}
\]
and either
\[
\tau(E)\le R_k
\qquad\text{or}\qquad
\operatorname{arank}(E)\le R_k.
\]
\end{definition}

\begin{theorem}[BridgeLock closes the continuation-cone route]\label{thm:bridgelock-closes}
Assume the finite base-window local-bridge package of Definition~\ref{def:bridgelock-package} and assume $\mathbf{BridgeLock}_k$. Then the witness pressure is uniformly bounded on unsolved continuation cones. Consequently, if the solved classes entering Theorem~\ref{thm:bounded-transversal-capture} carry their certified exponential bounds, then there exists $c_k<\infty$ such that
\[
f(n,k)<c_k^n
\qquad\text{for all }n.
\]
\end{theorem}

\begin{proof}
Let $L\in \mathsf{LB}_k$. By $\mathbf{BridgeLock}_k$, there exists an exact subcone $E\subseteq L$ with mass at least $2^{-R_k}$ and either bounded transversal or bounded active rank.

If $\tau(E)\le R_k$, then by definition of witness pressure,
\[
w(L)\le R_k+\left\lceil \log_2\frac{1}{\mu_L(E)}\right\rceil\le 2R_k.
\]

If instead $\operatorname{arank}(E)\le R_k$, then Proposition~\ref{prop:bounded-rank-transversal} gives
\[
\tau(E)\le R_k(k-1),
\]
so again
\[
w(L)\le R_k(k-1)+\left\lceil \log_2\frac{1}{\mu_L(E)}\right\rceil
\le R_k(k-1)+R_k.
\]
Thus the witness pressure is uniformly bounded on $\mathsf{LB}_k$.

If witness pressure were still unbounded on $\mathsf U_k$, Proposition~\ref{prop:bridgelock-concentration} would force it to be unbounded on $\mathsf{LB}_k$, contradiction. So the witness pressure is uniformly bounded on all unsolved continuation cones. Theorem~\ref{thm:bounded-transversal-capture} then converts this uniform bound into an exponential sunflower estimate once the solved local classes are inserted with their certified exponential bases.
\end{proof}

\begin{proposition}[Route-disjointness implies BridgeLock]\label{prop:route-disjointness-bridgelock}
Assume that every $L\in \mathsf{LB}_k$ comes with a family of heavy routed children $\xi$, each carrying a route support $R(\xi)$, and suppose there exists a constant $\rho_k$ such that:
\begin{enumerate}[label=(\roman*),leftmargin=2.2em]
\item $|R(\xi)|\le \rho_k$ for every heavy routed child $\xi$;
\item any $k$ pairwise route-disjoint heavy routed children lift to a genuine $k$-sunflower.
\end{enumerate}
Then $\mathbf{BridgeLock}_k$ holds.
\end{proposition}

\begin{proof}
Fix $L\in \mathsf{LB}_k$ and choose a maximal family $\Xi$ of pairwise route-disjoint heavy routed children. By assumption~(ii), one has
\[
|\Xi|\le k-1,
\]
for otherwise the corresponding $k$ children would already lift to a sunflower.

Let
\[
S:=\bigcup_{\xi\in \Xi} R(\xi).
\]
By assumption~(i),
\[
|S|\le (k-1)\rho_k.
\]
Maximality of $\Xi$ implies that every heavy routed child meets $S$; otherwise a child disjoint from $S$ would be route-disjoint from every member of $\Xi$, contradicting maximality.

Hence $S$ is a bounded separator lock for the heavy routed part of $L$. The exact subcone carried by heavy routed children therefore has bounded transversal at most $(k-1)\rho_k$, and positive mass by definition of heaviness. This is exactly the content of $\mathbf{BridgeLock}_k$.
\end{proof}

\begin{remark}[What BridgeLock isolates]\label{rem:bridgelock-isolates}
The continuation-cone lane has therefore been reduced to one genuinely local theorem on the common bridge class $\mathsf{LB}_k$. The residue is no longer a zoo of unrelated base-window obstructions. Everything concentrates on a single local bridge problem, and a route-disjointness lemma inside that local normal form would already suffice to finish the continuation-cone route.
\end{remark}

\status{The witness-pressure part of the continuation-cone route is now split cleanly into a rigorous concentration theorem and one local open statement, namely $\mathbf{BridgeLock}_k$. The extra route-disjointness criterion shows that any future bounded-width routing theorem would close this lane immediately.}


\begin{conjecture}[Bounded sunflower-trace piece lemma]
For each fixed $k$, there exists $R_k$ such that every cone-reduced $k$-sunflower-free uniform family contains a support piece $P$ of some rank $1\le r\le R_k$ whose trace family $\Tr_P(F)$ is itself $k$-sunflower-free.
\end{conjecture}

\begin{remark}
If true, this conjecture immediately implies $f(n,k)\le c_k^n$ by exact trace recursion and induction on uniformity. It encapsulates the guiding ``bounded-rank local gadget'' philosophy of the project. At present it remains completely open.
\end{remark}

\subsection*{Corrected bounded trace-decomposition statements}
The preceding conjecture is stronger than necessary. What the recursion actually needs is bounded rank together with bounded decomposition complexity of the trace.

\begin{definition}[Sunflower-chromatic number of a trace family]
For a uniform family $\mathcal A$, define
\[
\chi_k^\odot(\mathcal A):=\min\Bigl\{q: \mathcal A=\bigsqcup_{i=1}^q \mathcal A_i,\ \mathcal A_i\text{ is }k\text{-sunflower-free}\Bigr\}.
\]
\end{definition}

\begin{theorem}[Corrected bounded sunflower-trace decomposition lemma]\label{thm:corrected-trace-piece}
Fix $k$, and set $g(m,k):=f(m,k)-1$ for $m\ge 0$, with $g(0,k):=1$. Suppose there exist constants $R_k,Q_k$ such that every $k$-sunflower-free $m$-uniform family $G$ contains a support piece $P$ of rank $r$, with $1\le r\le R_k$, and
\[
\chi_k^\odot(\Tr_P(G))\le Q_k.
\]
Then for every $m\ge 1$,
\[
g(m,k)\le \max_{1\le r\le R_k} Q_k\,g(r,k)\,g(m-r,k).
\]
Consequently there exists $c_k<\infty$ such that
\[
g(n,k)\le c_k^n
\qquad\text{and}\qquad
f(n,k)<(c_k+1)^n
\]
for all $n$.
\end{theorem}

\begin{proof}
Fix a $k$-sunflower-free $m$-uniform family $G$ with $|G|=g(m,k)$. Choose $P$ as in the hypothesis, and write
\[
\Tr_P(G)=T_1\sqcup\cdots\sqcup T_q,
\qquad q\le Q_k,
\]
with each $T_j$ $k$-sunflower-free. Let $G^{(j)}\subseteq G$ be the subfamily whose trace on $P$ lies in $T_j$. By exact trace recursion,
\[
|G^{(j)}|=\sum_{\tau\in T_j}|G_\tau|.
\]
Each branch $G_\tau$ is $(m-r)$-uniform and $k$-sunflower-free, so $|G_\tau|\le g(m-r,k)$. Also $T_j$ is an $r$-uniform $k$-sunflower-free family, so $|T_j|\le g(r,k)$. Therefore
\[
|G^{(j)}|\le |T_j|\,g(m-r,k)\le g(r,k)g(m-r,k).
\]
Summing over $j$ gives
\[
g(m,k)=|G|\le q\,g(r,k)g(m-r,k)\le Q_k\,g(r,k)g(m-r,k).
\]
Taking the maximum over $1\le r\le R_k$ proves the displayed recurrence.

Now define
\[
c_k:=\max_{1\le t\le R_k}\bigl(Q_k g(t,k)\bigr)^{1/t}.
\]
For $m\le R_k$, the choice $t=m$ shows $g(m,k)\le c_k^m$. For $m>R_k$, choose $r\le R_k$ from the hypothesis and use the recurrence together with induction:
\[
g(m,k)\le Q_k g(r,k)g(m-r,k)\le c_k^r c_k^{m-r}=c_k^m.
\]
Thus $g(n,k)\le c_k^n$ for all $n$, and the displayed bound for $f(n,k)$ follows from $f(n,k)=g(n,k)+1$.
\end{proof}

\begin{definition}[Trace-class decomposition number]
Let $\mathcal X$ be a class of uniform families. For a uniform family $\mathcal A$, define
\[
\delta_{\mathcal X}(\mathcal A):=\min\Bigl\{q: \mathcal A=\bigsqcup_{i=1}^q \mathcal A_i,\ \mathcal A_i\in \mathcal X\Bigr\}.
\]
\end{definition}

\begin{theorem}[Bounded-rank solved-trace decomposition lemma]\label{thm:bounded-rank-solved-trace}
Fix $k$, and let $\mathcal X_k$ be a class of trace families. Write
\[
g_{\mathcal X_k}(r):=\sup\{ |\mathcal A| : \mathcal A\in \mathcal X_k,\ \mathcal A\text{ has rank }r\}.
\]
Assume there are constants $R_k,Q_k$ such that every $k$-sunflower-free $m$-uniform family $G$ contains a support piece $P$ of rank $r$, with $1\le r\le R_k$, and
\[
\delta_{\mathcal X_k}(\Tr_P(G))\le Q_k.
\]
Then, with $g(m,k):=f(m,k)-1$ and $g(0,k):=1$,
\[
g(m,k)\le \max_{1\le r\le R_k} Q_k\,g_{\mathcal X_k}(r)\,g(m-r,k).
\]
If in addition $g_{\mathcal X_k}(r)\le a_k^r$ for all $r$, then there exists $c_k<\infty$ such that
\[
g(n,k)\le c_k^n
\qquad\text{and hence}\qquad
f(n,k)<(c_k+1)^n
\]
for all $n$.
\end{theorem}

\begin{proof}
The proof is the same as above. Decompose $\Tr_P(G)$ into at most $Q_k$ pieces that belong to $\mathcal X_k$, bound each such piece by $g_{\mathcal X_k}(r)$, and bound each projected branch by $g(m-r,k)$. This gives the displayed recurrence. If $g_{\mathcal X_k}(r)\le a_k^r$, define
\[
c_k:=\max\Bigl\{a_k,\ \max_{1\le t\le R_k}\bigl(Q_k g_{\mathcal X_k}(t)\bigr)^{1/t}\Bigr\},
\]
and repeat the induction from Theorem~\ref{thm:corrected-trace-piece}.
\end{proof}

\begin{remark}
Taking $\mathcal X_k$ to be the class of $k$-sunflower-free trace families recovers Theorem~\ref{thm:corrected-trace-piece}. Taking $\mathcal X_k$ to be the recursive closure generated by the exact block-product, positive-density support-box, and fixed-memory finite-state classes gives the strongest version of this statement currently justified by the proved framework.
\end{remark}

\status{The naive $Q_k=1$ trace conjecture is stronger than necessary. The exact recursion only needs bounded rank together with bounded decomposition complexity of the trace.}

\subsection*{Fractional solved-trace pressure}
The bounded decomposition lemmas above still use an integral partition of the trace. Since exact trace recursion is additive, the sharp local invariant is instead a \emph{fractional} cover by solved subtraces.

\begin{definition}[Solved-trace pressure]\label{def:solved-trace-pressure}
Fix $k$, and let $\mathfrak S_k$ be a solved trace class in the sense of Section~\ref{sec:union-classes}, with one-step bound function $B_{\mathfrak S_k}(r,k)$. Equivalently, whenever $G$ is $k$-sunflower-free, $P$ is a support piece of rank $r$, and $U\subseteq \Tr_P(G)$ lies in $\mathfrak S_k$, one has
\[
|G[U]|\le B_{\mathfrak S_k}(r,k)\,g(m-r,k),
\]
where $m$ is the uniformity of $G$ and $g(t,k):=f(t,k)-1$. For a finite trace family $T$, define its \emph{solved-trace pressure} by
\[
\operatorname{sp}_k(T):=
\inf\left\{
\sum_{i=1}^s \lambda_i :
\lambda_i\ge 0,\ U_i\subseteq T,\ U_i\in \mathfrak S_k,\ 
\sum_{i:\,\tau\in U_i}\lambda_i\ge 1\ \text{for every }\tau\in T
\right\}.
\]
Thus $\operatorname{sp}_k(T)$ is the fractional cover number of $T$ by solved subtraces.
\end{definition}

\begin{theorem}[Bounded-rank solved-trace pressure lemma]\label{thm:solved-trace-pressure}
Fix $k$ and a solved trace class $\mathfrak S_k$. Suppose there exist constants $R_k,M_k$ such that every cone-reduced $k$-sunflower-free $m$-uniform family $G$ has a support piece $P$ of rank $r$, with $1\le r\le R_k$, for which
\[
\operatorname{sp}_k(\Tr_P(G))\le M_k.
\]
Set
\[
C_k:=M_k\max_{1\le r\le R_k} B_{\mathfrak S_k}(r,k).
\]
Then, with $g(m,k):=f(m,k)-1$,
\[
g(m,k)\le C_k\max_{1\le r\le R_k} g(m-r,k)
\qquad\text{for every }m\ge 1.
\]
Consequently there exists $c_k<\infty$ such that
\[
g(n,k)\le c_k^n
\qquad\text{and hence}\qquad
f(n,k)<(c_k+1)^n
\]
for all $n$.
\end{theorem}

\begin{proof}
Fix a cone-reduced $k$-sunflower-free $m$-uniform family $G$, and choose a support piece $P$ of rank $r\le R_k$ such that
\[
T:=\Tr_P(G)
\qquad\text{satisfies}\qquad
\operatorname{sp}_k(T)\le M_k.
\]
For a subtrace $U\subseteq T$, write
\[
G[U]:=\bigcup_{\tau\in U} G_\tau.
\]
Then $G[U]$ is still $k$-sunflower-free, $P$ is a support piece of rank $r$ for $G[U]$, and $\Tr_P(G[U])=U$. Therefore Theorem~\ref{thm:unionclass-closure} gives
\[
|G[U]|\le B_{\mathfrak S_k}(r,k)\,g(m-r,k)
\le
\Bigl(\max_{1\le s\le R_k} B_{\mathfrak S_k}(s,k)\Bigr) g(m-r,k).
\]

Fix $\eps>0$, and choose a fractional solved cover $(U_i,\lambda_i)_{i=1}^s$ witnessing
\[
\sum_{i=1}^s \lambda_i\le M_k+\eps.
\]
Since every $\tau\in T$ is covered with total weight at least $1$,
\[
|G|
=
\sum_{\tau\in T}|G_\tau|
\le
\sum_{i=1}^s \lambda_i\sum_{\tau\in U_i}|G_\tau|
=
\sum_{i=1}^s \lambda_i |G[U_i]|.
\]
Applying the solved one-step bound to each $U_i$ yields
\[
|G|
\le
(M_k+\eps)\Bigl(\max_{1\le s\le R_k} B_{\mathfrak S_k}(s,k)\Bigr) g(m-r,k).
\]
Letting $\eps\to 0$ gives
\[
g(m,k)\le C_k\,g(m-r,k)\le C_k\max_{1\le s\le R_k} g(m-s,k),
\]
which is the claimed recurrence.

Now define
\[
c_k:=\max\left\{C_k,\ \max_{1\le t\le R_k} g(t,k)^{1/t}\right\}.
\]
For $1\le m\le R_k$, the choice of $c_k$ gives $g(m,k)\le c_k^m$. For $m>R_k$, induction on $m$ and the recurrence yield
\[
g(m,k)\le C_k\max_{1\le s\le R_k} g(m-s,k)
\le
C_k\max_{1\le s\le R_k} c_k^{m-s}
\le
C_k c_k^{m-1}
\le
c_k^m,
\]
because $C_k\le c_k$. Hence $g(n,k)\le c_k^n$ for all $n$, and therefore
\[
f(n,k)=g(n,k)+1<(c_k+1)^n.
\]
\end{proof}

\begin{definition}[Solved-trace packing number]\label{def:solved-trace-packing}
With $\mathfrak S_k$ fixed as above, define
\[
\operatorname{pack}_k(T):=
\sup\left\{
\sum_{\tau\in T} x_\tau :
x_\tau\ge 0,\ 
\sum_{\tau\in U} x_\tau\le 1\ \text{for every solved subtrace }U\in \mathfrak S_k,\ U\subseteq T
\right\}.
\]
\end{definition}

\begin{proposition}[LP duality for solved-trace pressure]\label{prop:solved-trace-pressure-dual}
For every finite trace family $T$,
\[
\operatorname{sp}_k(T)=\operatorname{pack}_k(T).
\]
\end{proposition}

\begin{proof}
The definition of $\operatorname{sp}_k(T)$ is a finite linear program whose variables are the cover weights $\lambda_U$ indexed by solved subtraces $U\subseteq T$ and whose constraints require
\[
\sum_{U\ni \tau}\lambda_U\ge 1
\qquad(\tau\in T).
\]
Its dual has variables $x_\tau\ge 0$ and constraints
\[
\sum_{\tau\in U}x_\tau\le 1
\qquad(U\in \mathfrak S_k,\ U\subseteq T),
\]
with objective $\sum_{\tau\in T}x_\tau$. This is exactly the program defining $\operatorname{pack}_k(T)$. Standard finite-dimensional linear-programming duality therefore gives the claimed equality.
\end{proof}

\begin{corollary}[Uniform solved-trace capture implies bounded pressure]\label{cor:solved-trace-capture}
Fix $k$ and a solved trace class $\mathfrak S_k$. Suppose there exist constants $R_k$ and $\delta_k>0$ such that every cone-reduced $k$-sunflower-free $m$-uniform family $G$ has a support piece $P$ of rank $1\le r\le R_k$ with the following property: for every probability measure $\mu$ on
\[
T:=\Tr_P(G),
\]
there exists a solved subtrace $U\subseteq T$, $U\in \mathfrak S_k$, such that
\[
\mu(U)\ge \delta_k.
\]
Then
\[
\operatorname{sp}_k(T)\le \frac{1}{\delta_k}
\]
for that support piece, so Theorem~\ref{thm:solved-trace-pressure} applies. In particular there exists $c_k<\infty$ such that
\[
f(n,k)<c_k^n
\qquad\text{for all }n.
\]
\end{corollary}

\begin{proof}
Let $(x_\tau)_{\tau\in T}$ be any feasible dual packing from Proposition~\ref{prop:solved-trace-pressure-dual}, and set
\[
m:=\sum_{\tau\in T} x_\tau.
\]
If $m=0$, there is nothing to prove. Otherwise define a probability measure on $T$ by
\[
\mu(\tau):=\frac{x_\tau}{m}.
\]
By hypothesis there exists a solved subtrace $U\subseteq T$ with $\mu(U)\ge \delta_k$. Since $x$ is dual-feasible,
\[
1\ge \sum_{\tau\in U} x_\tau = m\,\mu(U)\ge m\delta_k.
\]
Hence $m\le 1/\delta_k$. Taking the supremum over all dual-feasible packings and using Proposition~\ref{prop:solved-trace-pressure-dual} gives
\[
\operatorname{sp}_k(T)=\operatorname{pack}_k(T)\le \frac{1}{\delta_k}.
\]
Theorem~\ref{thm:solved-trace-pressure} now yields the exponential bound.
\end{proof}

\leannote{The dual-packing step, namely that uniform solved-class capture bounds every feasible packing by \(\delta^{-1}\), is formalized and proved in \texttt{FormalConjectures/Problems/Erdos/E20/Compendium/Recurrences/CapturePressure.lean} as \texttt{uniformSolvedCapture\_gives\_pressureBound}. The local mechanism producing the capture hypothesis and the subsequent global exponential conclusion are not encoded in that file as completed Lean theorems.}

\begin{corollary}[Positive-overlap solved covers suffice]\label{cor:solved-trace-positive-overlap}
Fix $k$, and let $\mathfrak S_k$ be as above. Suppose every cone-reduced $k$-sunflower-free family $G$ has a support piece $P$ and solved subtraces
\[
U_1,\dots,U_q\subseteq \Tr_P(G)
\]
such that every trace atom lies in at least $\theta q$ of them:
\[
\sum_{i=1}^q \mathbf 1_{U_i}(\tau)\ge \theta q
\qquad\text{for every }\tau\in \Tr_P(G).
\]
Then $\operatorname{sp}_k(\Tr_P(G))\le 1/\theta$, and hence $f(n,k)<c_k^n$ for some $c_k<\infty$.
\end{corollary}

\begin{proof}
Assign weight $\lambda_i:=(\theta q)^{-1}$ to each $U_i$. The multiplicity hypothesis gives
\[
\sum_{i:\,\tau\in U_i}\lambda_i\ge 1
\qquad\text{for every }\tau\in \Tr_P(G),
\]
so the family $(U_i,\lambda_i)$ is a valid fractional solved cover of total weight
\[
\sum_{i=1}^q \lambda_i = \frac{1}{\theta}.
\]
Thus $\operatorname{sp}_k(\Tr_P(G))\le 1/\theta$, and Theorem~\ref{thm:solved-trace-pressure} applies.
\end{proof}



\begin{remark}[Weighted solved-trace pressure]\label{rem:weighted-solved-trace-pressure}
If solved trace classes come with class-dependent costs, so that each solved subtrace $U$ only guarantees
\[
|G[U]|\le w(U)\,g(m-r,k),
\]
then one replaces Definition~\ref{def:solved-trace-pressure} by the weighted infimum
\[
\inf\left\{
\sum_i \lambda_i w(U_i):
\lambda_i\ge 0,\ U_i\in \mathfrak S_k,\ 1_T\le \sum_i \lambda_i 1_{U_i}
\right\}.
\]
The proof of Theorem~\ref{thm:solved-trace-pressure} is unchanged, and the same linear-programming duality yields the corresponding weighted packing formulation.
\end{remark}


\begin{remark}[Why pressure is the sharp local invariant]
The bounded-rank integral decomposition lemma is the special case of solved-trace pressure in which the covering weights are $0$ or $1$ and the solved pieces are disjoint. Because exact trace recursion is linear, no weaker local invariant can rule out the dual packing obstruction of Definition~\ref{def:solved-trace-packing}. In that precise sense, bounded solved-trace pressure is the sharp endpoint of the present trace-recursion method.
\end{remark}

\status{The integral decomposition lemmas remain rigorous and useful. The pressure formulation shows exactly what the trace recursion really needs: bounded rank together with bounded \emph{fractional} solved-cover complexity. In particular, bounded solved-trace pressure would imply an exponential sunflower bound, and the stronger uniform-capture property of Corollary~\ref{cor:solved-trace-capture} would imply bounded pressure. What is still missing is a proof of either local statement in full generality.}


\subsection*{OA quotients, Bellman propagation, and neutral obstructions}
The LP side of the project suggests a second route. Its exact rigorous core is the following support-function propagation lemma.

\begin{proposition}[Abstract Bellman-face propagation]\label{prop:bellman-face}
Fix a finite alphabet $\Omega$ and an integer $L\ge 1$. Suppose a class of projected branches is equipped with:
\begin{enumerate}[label=(\roman*),leftmargin=2.2em]
\item a profile map $p(B)\in \R^D$;
\item a compact convex set $\mathcal D\subseteq \R^D$;
\item a quantity $\tau(B)=h_{\mathcal D}(p(B))$, where $h_{\mathcal D}(x):=\sup_{y\in \mathcal D}\langle x,y\rangle$;
\end{enumerate}
such that for every bounded prefix split of length $1\le \ell\le L$ one has
\[
p(B)=\sum_{u\in \Omega^\ell}\lambda(u)\,p(B_u)
\]
for weights $\lambda(u)\ge 0$ summing to $1$. Then:
\begin{enumerate}[label=(\alph*),leftmargin=2.2em]
\item for every bounded prefix split,
\[
\tau(B)\le \sum_{u\in \Omega^\ell}\lambda(u)\tau(B_u);
\]
\item if equality holds for every bounded prefix split throughout the descendant tree of $B$, then there exists a single point $y_*\in \mathcal D$ that simultaneously maximizes $h_{\mathcal D}(p(B_u))$ for every reachable descendant $B_u$;
\item if the reachable profile set $\{p(B_u):u\text{ reachable}\}$ is finite, then the descendant tree of $B$ is generated by a finite-state recursive system.
\end{enumerate}
\end{proposition}

\begin{proof}
Part (a) is just convexity of the support function:
\[
\tau(B)=h_{\mathcal D}\Bigl(\sum_u \lambda(u)p(B_u)\Bigr)
\le \sum_u \lambda(u)h_{\mathcal D}(p(B_u))
=\sum_u \lambda(u)\tau(B_u).
\]
For (b), apply equality at the root. Choose $y_*\in \mathcal D$ maximizing $h_{\mathcal D}(p(B))$. Since
\[
\langle p(B),y_*\rangle=\sum_u \lambda(u)\langle p(B_u),y_*\rangle
\]
and the right-hand side is bounded above by $\sum_u \lambda(u)h_{\mathcal D}(p(B_u))=h_{\mathcal D}(p(B))$, every active child $B_u$ must satisfy
\[
\langle p(B_u),y_*\rangle=h_{\mathcal D}(p(B_u)).
\]
Thus the same $y_*$ maximizes every child. Iterating down the descendant tree shows that $y_*$ maximizes every reachable profile. Part (c) is then immediate: if only finitely many profiles occur, use those profiles as states and let the transitions be induced by one-block extensions.
\end{proof}

\begin{conjecture}[OA-factor routing trichotomy]\label{conj:oa-quotient}
For each fixed $k$ there exist constants $t,L,D$ and a continuation-profile model on projected branches such that, after quotienting a projected branch by its maximal exact OA$_t$ factor, one of the following occurs:
\begin{enumerate}[label=(\roman*),leftmargin=2.2em]
\item some prefix split of length at most $L$ is profitable;
\item the reachable continuation profiles form a finite-state recursive class;
\item all reachable continuation profiles share a common maximizing dual witness as in Proposition~\ref{prop:bellman-face}(b).
\end{enumerate}
\end{conjecture}

\begin{remark}[Proof strategy for Conjecture~\ref{conj:oa-quotient}]
The intended proof is to combine Proposition~\ref{prop:oa-dual} with Proposition~\ref{prop:bellman-face}. If no profitable prefix exists, the Bellman inequality should be tight on every bounded split, forcing a common maximizing dual witness down the entire descendant tree. A finite orbit then gives a solved finite-state class; an infinite orbit with a common witness is the residual neutral obstruction. Thus the real missing step is a \emph{neutral-rigidity theorem} converting that last case into finite-state behavior.
\end{remark}

\status{This route is not yet self-contained in the present paper. The missing ingredients are a formal branchwise maximal OA quotient, a quantitative Bellman gap that turns strict convexity into a uniform profitable prefix, and a proof of neutral rigidity.}

\subsection*{Weighted-profile OA quotients}
The support-only OA discussion from Section~\ref{sec:lp} becomes more precise after passing to the full weighted profile monoid seen by the exact LP.

\begin{definition}[Weighted-profile visible-data quotient]\label{def:weighted-oa-quotient}
Fix an exposed dual face $\phi$ of the normalized block-box LP, and let $\Sigma=\Sigma_k$ be the finite alphabet of local block types. A projected branch with local weighted profile
\[
x=(x_\sigma)_{\sigma\in\Sigma}\in \N^\Sigma
\]
is sent by the \emph{homogenized visible-data map}
\[
A_\phi:\Z^\Sigma\to \Gamma_\phi
\]
to its total mass together with every linear statistic visible to the LP on the face $\phi$. Define the corresponding exact OA subgroup and quotient monoid by
\[
L_{\mathrm{OA}}(\phi):=\ker A_\phi\cap \Z^\Sigma,
\qquad
Q_\phi:=\N^\Sigma/L_{\mathrm{OA}}(\phi).
\]
\end{definition}

\begin{proposition}[Kernel directions are exactly weighted exact-OA factors]\label{prop:weighted-oa-kernel}
With notation as above, the subgroup $L_{\mathrm{OA}}(\phi)$ records precisely the invisible weighted directions on the face $\phi$.
More explicitly:
\begin{enumerate}[label=(\roman*),leftmargin=2.2em]
\item if, after clearing denominators, a profile splits as
\[
x=x^{(1)}+\cdots+x^{(r)}
\]
and the pieces $x^{(i)}$ all have the same normalized visible data on the face $\phi$, then all cleared pairwise differences lie in $L_{\mathrm{OA}}(\phi)$;
\item conversely, if $u\in L_{\mathrm{OA}}(\phi)\setminus\{0\}$ and $u=u^+-u^-$ is its positive/negative decomposition, then
\[
A_\phi(u^+)=A_\phi(u^-),
\]
so $u^+$ and $u^-$ define an exact OA factor invisible to every face-visible statistic.
\end{enumerate}
\end{proposition}

\begin{proof}
For (i), equal normalized visible data means that after clearing denominators there exist positive integers $m_i$ with
\[
A_\phi(m_i x^{(i)})=A_\phi(m_j x^{(j)})
\qquad\text{for all }i,j.
\]
Therefore
\[
m_jx^{(i)}-m_ix^{(j)}\in \ker A_\phi\cap \Z^\Sigma=L_{\mathrm{OA}}(\phi).
\]

For (ii), write $u=u^+-u^-$ with $u^+,u^-\in \N^\Sigma$ having disjoint supports. Since $u\in \ker A_\phi$,
\[
A_\phi(u^+)=A_\phi(u^-).
\]
Hence the two weighted profiles $u^+$ and $u^-$ have identical visible data on the face $\phi$. Interpreting a hidden choice between $u^+$ and $u^-$ as an auxiliary bit produces an exact OA factor invisible to all LP-visible statistics on that face.
\end{proof}

\begin{proposition}[Finite generation of the weighted OA quotient]\label{prop:weighted-oa-fg}
For every exposed face $\phi$, the monoid $Q_\phi$ is finitely generated. Equivalently,
\[
Q_\phi\cong A_\phi(\N^\Sigma)\subseteq \Gamma_\phi
\]
is generated by the finitely many images $A_\phi(e_\sigma)$, where $e_\sigma$ ranges over the standard basis vectors of $\N^\Sigma$.
\end{proposition}

\begin{proof}
Every element of $\N^\Sigma$ is a finite nonnegative integer combination of the basis vectors $e_\sigma$. Applying the additive map $A_\phi$ shows that every element of $A_\phi(\N^\Sigma)$ is a finite nonnegative integer combination of the finitely many generators $A_\phi(e_\sigma)$. The displayed isomorphism is immediate from the definition of the quotient.
\end{proof}

\begin{proposition}[Conditional semilinear collapse under a regular-interface hypothesis]\label{prop:weighted-oa-semilinear}
Fix an exposed face $\phi$, and suppose that for each ordered pair of interface states $(q,q')$ the OA-reduced controlled branches from $q$ to $q'$ are generated by a regular language over a finite alphabet of primitive profile increments from $\Sigma$. Then the set of quotient profiles realized between $q$ and $q'$ is semilinear in $Q_\phi$; equivalently, it is a finite union of sets of the form
\[
a+\N b_1+\cdots+\N b_r.
\]
\end{proposition}

\begin{proof}
Choose a finite alphabet of primitive moves whose cumulative profile increment is recorded in $\N^\Sigma$ by the usual Parikh-count map. By hypothesis, the admissible words between the interface states $q$ and $q'$ form a regular language $L$. Parikh's theorem states that the set of Parikh vectors of a regular language is semilinear in $\N^\Sigma$. Applying the additive quotient map
\[
\N^\Sigma \longrightarrow Q_\phi\cong A_\phi(\N^\Sigma)
\]
sends semilinear sets to semilinear sets, because affine translates and finitely generated submonoids remain affine translates and finitely generated submonoids under additive homomorphisms. Hence the realized quotient-profile set is semilinear in $Q_\phi$.
\end{proof}

\begin{remark}[Corrected status of the OA-quotient route]\label{rem:weighted-oa-status}
Propositions~\ref{prop:weighted-oa-kernel}--\ref{prop:weighted-oa-semilinear} rigorously identify the right quotient and show that any \emph{regular-interface} controlled model collapses to a semilinear family of quotient profiles. What they do \emph{not} yet prove is that every residual neutral obstruction satisfies the regular-interface hypothesis, nor that every semilinear component has already been incorporated into the solved finite-state recursive closure. Thus the weighted-profile quotient repairs the earlier support-level ambiguity, but the full neutral-rigidity theorem remains conditional.
\end{remark}

\status{The exact OA kernel is now identified on weighted profiles rather than raw supports. The remaining gap is no longer the definition of the quotient, but the passage from a residual neutral branch to a regular-interface model whose quotient image is semilinear.}




\subsection*{Semilinear quotient profiles do not imply finite-state collapse}
The weighted-profile OA quotient identifies the correct visible monoid, but semilinearity of that monoid is still weaker than finite-state recursive control.

\begin{definition}[Continuation congruence on an OA-reduced branch machine]\label{def:oa-continuation-congruence}
Fix an OA-reduced branch machine with reachable quotient profiles $Q$, a finite alphabet $\Sigma$ of bounded prefix-split types, partial continuation maps
\[
\tau_a:Q\rightharpoonup Q
\qquad (a\in \Sigma),
\]
and the associated terminal/output data carried by the recursion.
Define a congruence relation $\equiv$ on $Q$ by declaring
\[
q\equiv q'
\]
if and only if for every finite word $w\in \Sigma^\ast$:
\begin{enumerate}[label=(\roman*),leftmargin=2.2em]
\item the continuation $w$ is legal from $q$ exactly when it is legal from $q'$;
\item whenever it is legal, the resulting terminal/output data agree.
\end{enumerate}
\end{definition}

\begin{proposition}[Exact OA-neutral rigidity criterion]\label{prop:oa-continuation-criterion}
An OA-reduced branch machine is a solved finite-state recursive class if and only if the continuation congruence of Definition~\ref{def:oa-continuation-congruence} has finite index.
\end{proposition}

\begin{proof}
Assume first that the continuation congruence has finite index.
Quotient $Q$ by $\equiv$.
Because legality, one-step continuation, and terminal/output data are all invariant under $\equiv$, the recursive dynamics descend to the finite quotient $Q/{\equiv}$.
This gives a finite-state recursive presentation.

Conversely, suppose the branch machine already has a correct finite-state recursive presentation with finite state set $S$.
Map each reachable quotient profile $q\in Q$ to the state $s(q)\in S$ assigned to it.
If $s(q)=s(q')$, then by correctness of the finite-state presentation every finite continuation word has the same legality and the same terminal/output data from $q$ and $q'$.
Thus equality of state labels is a finite-index congruence contained in $\equiv$, so $\equiv$ also has finite index.
\end{proof}

\begin{counterexample}[A semilinear one-counter survivor]\label{ce:oa-one-counter}
There exists an OA-reduced neutral branch machine whose reachable quotient-profile set is a single finitely generated semilinear ray, but whose continuation congruence has infinite index.
In particular, semilinear quotient profiles do \emph{not} force finite-state collapse.
\end{counterexample}

\begin{proof}
Work in the cone
\[
C=\R_{\ge 0}^2
\]
and on the exposed face
\[
\phi=\{(x,0):x\ge 0\},
\]
cut out by the supporting functional $(x,y)\mapsto y$.
Let
\[
Q=\phi\cap \Z^2=\{q_m:=(m,0):m\in \N\}.
\]
Define the visible map on $\phi$ by
\[
A_\phi(q_m)=m.
\]
Then
\[
\ker A_\phi=\{0\},
\]
so the machine is already OA-reduced.
Its visible quotient monoid is just the one-generated semilinear set
\[
Q=q_0+\N(q_1-q_0).
\]

Take the finite control alphabet
\[
\Sigma=\{u,d\},
\]
and define partial continuation maps by
\[
u(q_m)=q_{m+1},
\qquad
d(q_m)=
\begin{cases}
q_{m-1}, & m\ge 1,\\
\text{undefined}, & m=0.
\end{cases}
\]
Let the terminal set be
\[
F=\{q_0\}.
\]
Give every state Bellman value $0$.
Then every bounded prefix split is tight, so the machine is neutral in the abstract Bellman-face sense.

Suppose $\sim$ is any congruence on $Q$ that is stable under $u$ and $d$ and preserves terminal data.
Take $m<n$ and assume toward contradiction that $q_m\sim q_n$.
Applying $d^m$ gives
\[
q_0=d^m(q_m)\sim d^m(q_n)=q_{n-m}.
\]
But $q_0\in F$, while $q_{n-m}\notin F$.
This contradicts preservation of terminal data.
Hence no two distinct states are congruent.
The continuation congruence has infinite index, even though the visible quotient-profile set is a single semilinear ray.
\end{proof}

\begin{remark}[What the OA one-counter survivor means]\label{rem:oa-one-counter}
Counterexample~\ref{ce:oa-one-counter} does \emph{not} refute the sunflower exponential conjecture.
It only refutes a specific inference:
\[
\text{semilinear quotient-profile set}
\ \Longrightarrow\
\text{finite-state recursive class}.
\]
The hidden issue is that static semilinearity does not control depth along a neutral ray.
An unbounded one-counter parameter can survive OA reduction and remain visible to the continuation congruence.
\end{remark}

\subsection*{Predictive-state neutral stability}
The neutral obstruction route also has a second formulation that works directly with future kernels rather than low-order marginals.

\begin{conjecture}[Predictive-state neutral stability]\label{conj:predictive-state}
Fix $k\ge 3$ and write $q:=k-1$. There exist constants $\eps_0,\delta>0$ and $M,R\in \N$ such that for every $0<\eps<\eps_0$ and every sufficiently large terminal kernel $K$ of rank $n$, the following holds. If $K$ has no $\delta$-profitable prefix block and no solved trace piece of rank at most $R$, then there is a partition
\[
[n]=G\sqcup B,
\qquad |B|\le \eps n,
\]
and a subkernel $K'\subseteq K$ with
\[
|K'|\ge e^{-\eps n}|K|,
\]
such that on the good coordinates $G$ every fiber of $K'$ is either:
\begin{enumerate}[label=(\roman*),leftmargin=2.2em]
\item an exact block product of width at most $q$, with width exactly $q$ on all but $O(\eps n)$ good coordinates; or
\item a $q$-ary finite-state branch generated by at most $M$ states.
\end{enumerate}
Moreover the description is hereditary under passing to heavy prefixes and to long good-coordinate suffixes.
\end{conjecture}

\begin{remark}[Proposed proof strategy for Conjecture~\ref{conj:predictive-state}]
The intended proof has four stages.
\begin{enumerate}[label=(\roman*),leftmargin=2.2em]
\item Use Theorem~\ref{thm:positive-density-box} together with the absence of profitable prefixes to force nearly $q$-regular, nearly balanced branching on most coordinates.
\item Define \emph{predictive states} using a hereditary box metric on full future kernels rather than bounded-order marginals. Counterexample~\ref{ce:parity-proper-marginals} shows that bounded-order marginals cannot suffice.
\item Show that if two heavy prefixes have genuinely different predictive states, then the first disagreement occurs at bounded depth. The corresponding difference tensor lies in a top Moebius layer and should produce a bounded-rank solved trace piece, contradicting the hypothesis.
\item Once only finitely many predictive states remain, the good part is a finite-state branch. If only one predictive state remains, the good part collapses to the width-$q$ block-product template.
\end{enumerate}
\end{remark}

\begin{corollary}[Conditional consequence of predictive-state stability]\label{cor:predictive-state-exp}
If Conjecture~\ref{conj:predictive-state} holds, then for every fixed $k$ there exists $c_k<\infty$ such that
\[
f(n,k)<c_k^n
\qquad\text{for all }n.
\]
\end{corollary}

\begin{proof}
Let $u_k(n):=M_{\ker}(n,k)$. Choose $c_k$ larger than all currently certified exponential bases coming from the profitable-prefix classes, the solved trace classes, the exact block-product class, and the finite-state classes with at most $M$ states. We prove by induction on $n$ that $u_k(n)\le c_k^n$.

Let $K$ be a terminal kernel of rank $n$. If $K$ already lies in one of the certified profitable-prefix or solved-trace classes, the choice of $c_k$ gives the bound immediately. Otherwise apply Conjecture~\ref{conj:predictive-state}. We obtain a partition $[n]=G\sqcup B$ with $|B|\le \eps n$ and a subkernel $K'\subseteq K$ with $|K'|\ge e^{-\eps n}|K|$ such that each good-coordinate fiber is either an exact width-$\le q$ block product or one of finitely many $M$-state branches. Let
\[
\Theta_k:=\max\{q,\Lambda_k\},
\]
where $\Lambda_k$ is the largest certified exponential base among the $M$-state branches. Then the number of possible good-coordinate traces is at most $\Theta_k^{|G|}$, while each corresponding bad-coordinate fiber has rank at most $|B|$ and therefore size at most $(k-1)^{|B|}u_k(|B|,k)$ by Theorem~\ref{thm:exact-reduction}. Hence
\[
|K|\le e^{\eps n}\Theta_k^{|G|}(k-1)^{|B|}u_k(|B|,k)
\le e^{\eps n}\Theta_k^{(1-\eps)n}\bigl((k-1)c_k\bigr)^{\eps n}.
\]
Since $c_k>\Theta_k$, choosing $\eps>0$ sufficiently small gives
\[
e^{\eps}\Theta_k^{1-\eps}\bigl((k-1)c_k\bigr)^\eps<c_k,
\]
so $|K|<c_k^n$. This proves $u_k(n)\le c_k^n$. Finally,
\[
M(n,k)\le (k-1)^n u_k(n)\le \bigl((k-1)c_k\bigr)^n
\]
by Theorem~\ref{thm:exact-reduction}, and therefore $f(n,k)<\bigl((k-1)c_k+1\bigr)^n$.
\end{proof}

\status{The predictive-state route is still open. The current gaps are the hereditary box metric, the bounded witness-scale lemma, and the extraction of a bounded-rank solved trace piece from the first-disagreement tensor.}

\begin{theorem}[Exact predictive-state core]\label{thm:predictive-state-core}
Fix $k\ge 3$, write $q:=k-1$, and let $K$ be a terminal kernel on $n$ coordinates. Suppose there exist a set of good coordinates $G\subseteq [n]$ and a subkernel $K'\subseteq K$ such that:
\begin{enumerate}[label=(\roman*),leftmargin=2.2em]
\item $|G|\ge (1-\eps)n$ and $|K'|\ge e^{-\eps n}|K|$;
\item every hereditary restriction of $K'|_G$ belongs to one of at most $M$ \emph{exact predictive states};
\item if two prefixes have the same predictive state, then their continuation kernels on the remaining good coordinates are exactly identical;
\item every retained good coordinate has width at most $q$.
\end{enumerate}
Then $K'|_G$ is generated by a finite-state branching system with at most $M$ states. If in fact $M=1$, then
\[
K'|_G=\bigotimes_{i\in G}\mu_i
\]
for suitable one-coordinate measures $\mu_i$; and if every retained coordinate has width exactly $q$, this is an exact width-$q$ block product. In particular there exists $C=C(k,\eps,M)$ such that the number of good-coordinate traces of $K'|_G$ is at most $C^n$.
\end{theorem}

\begin{proof}
Take the predictive states themselves as the internal memory states of a branching process. By assumption (iii), once the current state is known, the continuation kernel on the remaining good coordinates is known exactly. Therefore for each state $s$ and each next good coordinate $i\in G$, the continuation kernel determines:
\begin{enumerate}[label=(\alph*),leftmargin=2.2em]
\item which symbols are allowed at the coordinate $i$;
\item for each allowed symbol, which predictive state the corresponding child prefix enters.
\end{enumerate}
This produces a well-defined finite transition system on at most $M$ states whose recursive expansion reproduces the good-coordinate part $K'|_G$ exactly. Hence $K'|_G$ is a finite-state branching model with at most $M$ states.

Now assume $M=1$. Then every heavy prefix has the same continuation kernel on the remaining good coordinates. In particular, the conditional law of the next good coordinate does not depend on the earlier good-coordinate prefix. Writing the good coordinates in some order $i_1,\dots,i_{|G|}$ and iterating the chain rule, the measure on $K'|_G$ factors as
\[
\mu_{i_1}\otimes \mu_{i_2}\otimes \cdots \otimes \mu_{i_{|G|}}.
\]
This is the claimed exact product decomposition. If every retained coordinate has width exactly $q$, then each factor $\mu_i$ is supported on exactly $q$ symbols, so the product is a genuine width-$q$ block product.

Finally, in a finite-state branching system with at most $M$ states and alphabet size at most $q$, the number of traces of length $|G|$ is controlled by products of $M\times M$ transition matrices whose row sums are at most $q$. Thus the trace count is at most $(Mq)^{|G|}\le (Mq)^n$. Any larger constant $C=C(k,\eps,M)$ also works.
\end{proof}

\begin{remark}[Audit of the current predictive-state proof attempt]\label{rem:predictive-state-audit}
The present predictive-state route does \emph{not} yet prove the conjectured exponential sunflower bound from first principles. Counterexample~\ref{ce:balanced-multislice} and Proposition~\ref{prop:mesoscopic-slice-density} show that the naive first step is actually false: an arbitrary large terminal kernel need not contain any positive-density dense box on $\Omega(n)$ microscopic coordinates. Any repaired entry point must either allow thin neutral slice obstructions or work at the weaker $e^{-o(n)}$ density scale. Even after such a repair, two further ingredients would still be needed: a bounded witness-scale lemma for disagreements, and a theorem that the first bounded disagreement produces a bounded-rank solved trace piece. What Theorem~\ref{thm:predictive-state-core} isolates is the exact part of the route that \emph{is} already rigorous: once hereditary exact predictive states are present, finite-state branching follows formally, and the one-state case collapses to an exact product.
\end{remark}

\status{The predictive-state route therefore has a completely rigorous endgame but no valid entry point in its original positive-density form. The naive dense-box extraction is false; any repair must admit thin neutral slice obstructions or work at the weaker $e^{-o(n)}$ density scale. The bounded witness-scale and first-disagreement steps remain open even after that repair.}




\subsection*{Bounded-memory quotient states as the sharpened predictive target}
The predictive-state route can be sharpened further.
The right target is not a global inverse theorem on low-order marginals, and not even eventual periodicity at the outset.
It is a bounded-memory quotient-state theorem.

\begin{conjecture}[Bounded-memory quotient-state theorem]\label{conj:quotient-state}
Fix $k\ge 3$.
There exist constants $R=R(k)$ and $\eta=\eta(k)>0$ such that every OA-reduced neutral branch $B$ satisfies one of the following alternatives:
\begin{enumerate}[label=(\roman*),leftmargin=2.2em]
\item $B$ has a profitable split or prefix of depth at most $R$, with Bellman gain at least $\eta$;
\item for any two descendants $u,v$ of $B$,
\[
q(u)=q(v)
\quad\text{and}\quad
K_{\le R}(u)=K_{\le R}(v)
\qquad\Longrightarrow\qquad
K(u)=K(v),
\]
where $q(\cdot)$ is the OA-visible quotient profile and $K_{\le R}(\cdot)$ is the future kernel restricted to cylinders of depth at most $R$.
\end{enumerate}
In words, after quotienting out the exact OA factor, the full predictive state is determined by bounded-depth future data unless a bounded profitable split already appears.
\end{conjecture}

\begin{remark}[Why Conjecture~\ref{conj:quotient-state} would be enough]\label{rem:quotient-state-enough}
If the solved finite-state package is enlarged to include hereditary exact bounded-memory predictive classes, then Conjecture~\ref{conj:quotient-state} already supplies the required routing theorem:
either a bounded profitable prefix occurs, or the branch enters a bounded-memory solved class.
The stronger finite-state or eventual-periodic conclusions become downstream corollaries rather than primitive hypotheses.
\end{remark}

\begin{remark}[Programmatic proof strategy for Conjecture~\ref{conj:quotient-state}]\label{rem:quotient-state-strategy}
The intended proof has six steps.
\begin{enumerate}[label=(\arabic*),leftmargin=2.2em]
\item \emph{First disagreement.}
Assume there is no profitable bounded split, and choose two descendants $u,v$ of a neutral branch with
\[
q(u)=q(v),\qquad K_{\le R}(u)=K_{\le R}(v),\qquad K(u)\ne K(v),
\]
for which the first disagreement depth is minimal.

\item \emph{M\"obius compression.}
Apply M\"obius inversion on the future-cylinder tree.
Minimality of the disagreement depth forces all lower layers to cancel, so the discrepancy compresses to a top witness supported at the first level where the future kernels differ.

\item \emph{Pass to a neutral terminal.}
Face propagation keeps all descendants on the same Bellman face.
After quotienting by the exact OA kernel, the top witness becomes a finite collection of one-step child discrepancies
\[
\partial_{P,c}\in V_T
\]
attached to a neutral terminal $T$, together with positive pressure increments $\rho_{P,c}>0$.

\item \emph{Weighted coherence.}
The child discrepancies should satisfy the convex-balance relations from Corollary~\ref{cor:exact-convex-lift-discrepancy}.
If the corresponding coherence polytope were empty, a separating functional would create a bounded profitable split, contradicting neutrality.

\item \emph{Bounded witness scale.}
An extreme point of the coherence polytope has support bounded in terms of the fixed local alphabet.
After clearing denominators, this yields a bounded integer witness supported on only $O_k(1)$ child types.

\item \emph{Exact box extraction and uniform Bellman gap.}
The bounded integer witness glues to an exact local box product with strictly positive total pressure increment.
Because only finitely many bounded witness types occur, the minimum positive Bellman gain is uniform, giving the profitable split in alternative~(i).
\end{enumerate}
The difficult steps are the finite-scale realization of the top witness and the passage from the local discrepancy polytope to a genuine bounded profitable split.
\end{remark}

\subsection*{Positive circuits and exact future-kernel one-counter survivors}
The OA route admits two exact repairs.
The first replaces the incorrect ``extreme point plus denominator clearing'' step by a finite positive-circuit extraction on the top discrepancy alphabet.
The second shows that any genuine rank-$1$ future-kernel survivor is already exponentially tame.

\begin{definition}[Positive circuit of the top discrepancy alphabet]\label{def:top-positive-circuit}
Let $\Lambda_k$ be a finite top discrepancy alphabet and let
\[
\partial:\Lambda_k\to \Z^{d_k}
\]
be the corresponding discrepancy map after the exact OA quotient.
A vector
\[
n=(n_\lambda)_{\lambda\in \Lambda_k}\in \N^{\Lambda_k}\setminus\{0\}
\]
is called a \emph{positive circuit} if
\begin{enumerate}[label=(\roman*),leftmargin=2.2em]
\item
\[
\sum_{\lambda\in \Lambda_k} n_\lambda\,\partial_\lambda=0;
\]
\item the support of $n$ is inclusion-minimal among nonzero nonnegative relations;
\item the coefficients are primitive:
\[
\gcd\{n_\lambda:\lambda\in \supp(n)\}=1.
\]
\end{enumerate}
\end{definition}

\begin{proposition}[Positive-circuit extraction]\label{prop:top-positive-circuit}
Let $x=(x_\lambda)_{\lambda\in \Lambda_k}\in \R_{\ge 0}^{\Lambda_k}$ satisfy
\[
\sum_{\lambda\in \Lambda_k} x_\lambda\,\partial_\lambda=0,
\]
and assume that some $\lambda\in \supp(x)$ has $\partial_\lambda\neq 0$.
Then the support of $x$ contains the support of a nonzero positive circuit $n$.
Moreover:
\begin{enumerate}[label=(\roman*),leftmargin=2.2em]
\item
\[
|\supp(n)|\le d_k+1;
\]
\item there exists a constant $B_k<\infty$, depending only on the finite alphabet $\Lambda_k$, such that
\[
\sum_{\lambda\in \Lambda_k} n_\lambda\le B_k
\]
for every positive circuit.
\end{enumerate}
\end{proposition}

\begin{proof}
Normalize $x$ by its total mass:
\[
\alpha_\lambda:=\frac{x_\lambda}{\sum_\mu x_\mu}.
\]
Then $\alpha_\lambda\ge 0$, $\sum_\lambda \alpha_\lambda=1$, and
\[
\sum_{\lambda\in \Lambda_k}\alpha_\lambda\,\partial_\lambda=0.
\]
Hence the origin lies in the convex hull of the discrepancy vectors appearing in $\supp(x)$.

Choose an inclusion-minimal subset $S\subseteq \supp(x)$ such that
\[
0\in \operatorname{conv}\{\partial_\lambda:\lambda\in S\}.
\]
By minimality, the points $\{\partial_\lambda:\lambda\in S\}$ are affinely independent.
Therefore
\[
|S|\le d_k+1.
\]
The barycentric coordinates of $0$ with respect to this simplex are unique and strictly positive:
\[
\sum_{\lambda\in S} \beta_\lambda\,\partial_\lambda=0,
\qquad
\sum_{\lambda\in S}\beta_\lambda=1,
\qquad
\beta_\lambda>0.
\]
Because the discrepancy vectors are integral, the coefficients $\beta_\lambda$ are rational.
Clear denominators and then divide by the common gcd.
This produces a primitive nonzero nonnegative integer relation
\[
n=(n_\lambda)_{\lambda\in \Lambda_k}
\]
supported on $S$.
By minimality of $S$, the support of $n$ is inclusion-minimal among nonzero nonnegative relations.
So $n$ is a positive circuit.

This proves the support bound.
Since $\Lambda_k$ is finite, there are only finitely many possible supports $S$.
For each support the primitive circuit is unique up to scaling and has finite $\ell^1$-norm.
Taking the maximum over all possible supports gives the constant $B_k$.
\end{proof}

\begin{remark}[Audit correction for the bounded-scale step]\label{rem:positive-circuit-audit}
The correct bounded witness-scale invariant is the positive circuit of the finite top discrepancy alphabet, not an arbitrary extreme point of the full coherence polytope.
The latter can carry denominators depending on the bookkeeping right-hand side.
Positive circuits avoid that defect and give a genuinely finite witness type.
\end{remark}

\begin{theorem}[Finite-phase weighted one-counter future-kernel systems are exponentially tame]\label{thm:future-kernel-one-counter}
Let $Q$ be a finite phase set and let $D\subset \Z$ be a finite jump set.
For each $\delta\in D$, let
\[
A_\delta(z)\in \operatorname{Mat}_Q(\R_{\ge 0}[[z]])
\]
have positive radius of convergence and satisfy $A_\delta(0)=0$.
Assume that after absorbing a finite boundary strip, the total generating series of a branch class is produced by a half-line system on states $(q,h)\in Q\times \N$ whose one-step transitions change the counter by $\delta\in D$ and contribute the local block series $A_\delta(z)$.

Define the tilted phase kernel
\[
K(z,u):=\sum_{\delta\in D} u^\delta A_\delta(z)
\]
for $u>0$, and set
\[
\Psi(z):=\inf_{u>0}\rho(K(z,u)),
\]
where $\rho$ denotes spectral radius.
Then the total generating series of the branch class is analytic at every $z$ with $\Psi(z)<1$.
In particular it has positive radius of convergence, and its coefficients are exponentially bounded.
\end{theorem}

\begin{proof}
Let $\mathcal H$ be the space of sequences $x=(x_h)_{h\ge 0}$ with $x_h\in \R^Q$.
Define the transfer operator
\[
(T(z)x)_g
:=
\sum_{\substack{\delta\in D\\ g-\delta\ge 0}} A_\delta(z)\,x_{g-\delta}.
\]
After absorbing the finite boundary strip into analytic start and end data, the total generating series of the system has the form
\[
F(z)=b(z)+\alpha(z)^\top (I-T(z))^{-1}\beta(z),
\]
where $b,\alpha,\beta$ are finite-dimensional analytic series.
So it suffices to show that $I-T(z)$ is invertible whenever $\Psi(z)<1$.

Fix $u>0$ and a strictly positive row vector $\ell\in \R_{>0}^Q$.
For a sequence $x\in \mathcal H$, define the weighted norm
\[
\|x\|_{u,\ell}:=\sum_{h\ge 0} u^h\,\ell^\top |x_h|.
\]
Then
\[
\begin{aligned}
\|T(z)x\|_{u,\ell}
&=
\sum_{g\ge 0} u^g \ell^\top
\left|
\sum_{\substack{\delta\in D\\ g-\delta\ge 0}} A_\delta(z)\,x_{g-\delta}
\right|\\
&\le
\sum_{g\ge 0} u^g \ell^\top
\sum_{\substack{\delta\in D\\ g-\delta\ge 0}} A_\delta(z)\,|x_{g-\delta}|\\
&=
\sum_{h\ge 0} u^h\,\ell^\top
\left(\sum_{\delta\in D} u^\delta A_\delta(z)\right)|x_h|\\
&=
\sum_{h\ge 0} u^h\,\ell^\top K(z,u)\,|x_h|.
\end{aligned}
\]
If $K(z,u)$ is irreducible, choose $\ell$ to be a positive left Perron vector:
\[
\ell^\top K(z,u)=\rho(K(z,u))\,\ell^\top.
\]
Then
\[
\|T(z)x\|_{u,\ell}\le \rho(K(z,u))\,\|x\|_{u,\ell}.
\]
If $K(z,u)$ is reducible, decompose it into strongly connected components and work in a dominant block; the same estimate follows with the global spectral radius.
Hence whenever $\rho(K(z,u))<1$, the operator $T(z)$ is a strict contraction in the weighted norm.
Therefore the Neumann series
\[
(I-T(z))^{-1}=\sum_{m\ge 0} T(z)^m
\]
converges in operator norm, and $F(z)$ is analytic at that value of $z$.

Finally, because each $A_\delta(0)=0$, one has
\[
K(0,u)=0
\qquad\text{for every }u>0.
\]
So $\rho(K(z,u))\to 0$ as $z\to 0$.
In particular, for all sufficiently small positive $z$ one has $\Psi(z)<1$.
Hence $F(z)$ has positive radius of convergence.
By Cauchy's estimate, its coefficients are exponentially bounded.
\end{proof}

\begin{corollary}[Corrected OA endgame]\label{cor:future-kernel-one-counter}
Any OA route that reduces every residual neutral survivor either to a bounded profitable split or to an exact finite-phase weighted future-kernel one-counter system already yields an exponential bound.
\end{corollary}

\begin{proof}
Bounded profitable splits are handled by the exact recursive machinery already developed earlier in this paper.
If a residual survivor is instead represented by an exact finite-phase weighted future-kernel one-counter system, Theorem~\ref{thm:future-kernel-one-counter} shows that its generating series has positive radius of convergence and therefore contributes only exponential growth.
Taking the maximum of the finitely many resulting exponential bases over the solved cases closes the OA route.
\end{proof}

\begin{remark}[The first genuine remaining OA gap]\label{rem:future-kernel-gap}
Theorem~\ref{thm:future-kernel-one-counter} does \emph{not} follow from semilinearity of reachable quotient profiles alone.
Counterexample~\ref{ce:oa-one-counter} already shows that a semilinear rank-$1$ shadow can retain infinite continuation congruence.
So the structural gap is the lift from a semilinear rank-$1$ shadow to an exact finite-phase weighted future-kernel one-counter representation.
Once that lift is available, the rank-$1$ survivor is harmless.
\end{remark}





\subsection*{Bounded-window realization and dynamic OA-neutral rigidity}
The OA route also admits a sharpened local target. If every rank-$1$ survivor gains only through bounded windows on bounded variable supports, then a finite-type sunflower argument forces a uniform gain bound.

\begin{definition}[Bounded-window realization on a rank-$1$ survivor ray]\label{def:oa-window-realization}
Fix $k$. Let $x$ be a state on a real positive-slope rank-$1$ OA survivor ray, and let $\Gamma(x)$ denote its Bellman gain. A \emph{bounded-window realization} of $\Gamma(x)$ consists of a family $\mathcal W_x$ of actual OA-distinct continuation windows together with constants $L_k,S_k,B_k$ such that:
\begin{enumerate}[label=(\roman*),leftmargin=2.2em]
\item every $W\in \mathcal W_x$ has length at most $L_k$;
\item the \emph{variable local support}
\[
\sigma_x(W)
\]
of $W$ relative to the forced local core at $x$ has size at most $S_k$;
\item
\[
\Gamma(x)\le B_k\,|\mathcal W_x|.
\]
\end{enumerate}
\end{definition}

\begin{theorem}[Bounded-window realization plus typed sunflower lifting imply dynamic OA-neutral rigidity]\label{thm:oa-window-rigidity}
Fix $k$. Assume there exist constants $L_k,S_k,B_k$ such that every state on every real positive-slope rank-$1$ survivor ray admits a bounded-window realization in the sense of Definition~\ref{def:oa-window-realization}. Assume moreover the following typed sunflower lifting principle:

\begin{quote}
Whenever $k$ continuation windows at the same state have the same local OA type, agree on the same forced local core, and have pairwise disjoint variable petals, they lift to a genuine forbidden $k$-sunflower.
\end{quote}

Then there exists $C_k=O_k(1)$ such that
\[
\Gamma(x)\le C_k
\]
for every state $x$ on every real positive-slope rank-$1$ survivor ray.
\end{theorem}

\begin{proof}
Fix such a state $x$ and let $\mathcal W_x$ be a bounded-window realization of its Bellman gain. Let $K_x$ denote the forced local core already determined by the exact terminal-kernel and trace data at $x$. For each continuation window $W\in \mathcal W_x$, the variable support
\[
\sigma_x(W)
\]
has size at most $S_k$ by hypothesis.

\emph{Step 1: only finitely many local OA types occur.}
For each $W\in \mathcal W_x$, canonically relabel the atoms of $\sigma_x(W)$ by order of first appearance and record:
\begin{enumerate}[label=(\alph*),leftmargin=2.2em]
\item the start and end control states;
\item the control path of length at most $L_k$;
\item the bounded counter increments;
\item the incidence pattern of the relabelled variable atoms in the local Bellman transitions.
\end{enumerate}
Because the control is finite and both the length and variable support are bounded, only finitely many such labelled patterns occur. Let
\[
T_{k,L_k,S_k}
\]
be the finite set of local OA types, and write
\[
M_k:=|T_{k,L_k,S_k}|<\infty.
\]
Partition
\[
\mathcal W_x=\bigsqcup_{\tau\in T_{k,L_k,S_k}} \mathcal W_x(\tau)
\]
according to local OA type.

\emph{Step 2: a large type class forces a typed local sunflower.}
Fix one type $\tau$ and consider the support family
\[
\Sigma_\tau:=\{\sigma_x(W): W\in \mathcal W_x(\tau)\}.
\]
Every member of $\Sigma_\tau$ has size at most $S_k$. Let $\mathsf{SF}(S,r)$ be any sunflower threshold for set systems of size at most $S$; for instance one may take
\[
\mathsf{SF}(S,r)\le S!(r-1)^S.
\]
Let $A_{S_k}$ be a crude bound on the number of possible placements of a common core inside a canonical support template of size at most $S_k$; for example
\[
A_{S_k}\le 2^{S_k}S_k!
\]
suffices.

Assume that
\[
|\mathcal W_x(\tau)|>\mathsf{SF}(S_k,kA_{S_k}).
\]
Then the family $\Sigma_\tau$ contains a sunflower of size $kA_{S_k}$:
\[
\sigma_x(W_i)=C\sqcup P_i
\qquad (1\le i\le kA_{S_k}),
\]
with common core $C$ and pairwise disjoint petals $P_i$. Among these windows, the way the actual common core $C$ sits inside the canonical labelled template of type $\tau$ has at most $A_{S_k}$ possibilities. By the pigeonhole principle, there exist $k$ windows
\[
W_1,\dots,W_k
\]
with the same local OA type, the same placement of the common core, and pairwise disjoint petals.

These $k$ windows therefore agree on the forced base $K_x\cup C$ and are pairwise disjoint outside it. By the typed sunflower lifting hypothesis, they lift to a genuine forbidden $k$-sunflower. This contradicts that $x$ lies on a real survivor ray. Hence
\[
|\mathcal W_x(\tau)|\le \mathsf{SF}(S_k,kA_{S_k})
\qquad\text{for every }\tau.
\]

\emph{Step 3: the Bellman gain is uniformly bounded.}
Summing over all types gives
\[
|\mathcal W_x|
\le
M_k\,\mathsf{SF}(S_k,kA_{S_k}).
\]
By bounded-window realization,
\[
\Gamma(x)\le B_k\,|\mathcal W_x|
\le
B_k\,M_k\,\mathsf{SF}(S_k,kA_{S_k}).
\]
The right-hand side depends only on $k$, so setting
\[
C_k:=B_k\,M_k\,\mathsf{SF}(S_k,kA_{S_k})
\]
proves the theorem.
\end{proof}

\begin{corollary}[Conditional OA closure from bounded-window realization]\label{cor:oa-window-exp}
Assume the hypotheses of Theorem~\ref{thm:oa-window-rigidity}. If every residual OA survivor is either reduced to a bounded profitable split or lies on a real positive-slope rank-$1$ survivor ray with bounded-window realization, then the OA route yields an exponential sunflower bound.
\end{corollary}

\begin{proof}
The bounded profitable split case is already handled by the exact recursive machinery in this document. In the rank-$1$ survivor case, Theorem~\ref{thm:oa-window-rigidity} gives a uniform bound on the Bellman gain along the ray. Iterating that uniform gain bound along the exact OA recursion produces only exponential growth. Taking the maximum over the finitely many local solved classes therefore gives an exponential sunflower bound.
\end{proof}

\begin{remark}[What this bypasses]\label{rem:oa-window-bypass}
Theorem~\ref{thm:oa-window-rigidity} bypasses the false inference
\[
\text{semilinear quotient profile}
\Longrightarrow
\text{finite-state recursion}.
\]
The real issue is not semilinearity but the possibility that gain is witnessed only by longer and longer return words or by windows whose variable support drifts to infinity. Once real branch-realized witnesses are compressed to bounded windows on bounded variable supports, the rank-$1$ OA survivor becomes harmless.
\end{remark}

\status{The OA lane now has two clean local endgames. The one-counter theorem earlier in this section handles exact finite-phase survivors, while the new bounded-window theorem shows that a uniform bounded-window realization would already force dynamic OA-neutral rigidity. The remaining OA gap is therefore a realization problem: produce bounded windows on bounded variable supports for every real rank-$1$ survivor.}


\subsection*{Recursive low-excess support-product extraction}
The balanced multislice counterexamples show that the naive microscopic dense-box statement is false.
A more robust programmatic invariant is the support-product excess of a balanced block partition.

\begin{definition}[Support-product excess]\label{def:support-product-excess}
Let $K$ be a terminal kernel on a coordinate set $I$, and let
\[
\mathcal P=(I_1,\dots,I_t)
\]
be an ordered partition of $I$.
Define the \emph{support-product excess} of $K$ across $\mathcal P$ by
\[
\operatorname{Exc}_{\mathcal P}(K)
:=
\sum_{r=1}^t \log |\pi_{I_r}K|-\log |K|.
\]
Equivalently,
\[
\left|\prod_{r=1}^t \pi_{I_r}K\right|
=
e^{\operatorname{Exc}_{\mathcal P}(K)}|K|,
\]
so $\operatorname{Exc}_{\mathcal P}(K)$ is exactly the logarithmic loss incurred by replacing $K$ with the support box $\prod_r \pi_{I_r}K$.
\end{definition}

\begin{conjecture}[Recursive low-excess extraction theorem]\label{conj:low-excess-extraction}
Fix $k\ge 3$.
There exists $\eta=\eta(k)>0$ such that for every $\varepsilon>0$, every sufficiently large terminal kernel $K$ on $n$ coordinates admits, after conditioning on $O_k(1)$ memory, one of the following:
\begin{enumerate}[label=(\roman*),leftmargin=2.2em]
\item a profitable prefix;
\item bounded solved-trace pressure;
\item a positive-density support box on a coordinate set of size at least $\eta n$, in the sense of Theorem~\ref{thm:positive-density-box};
\item an $\eta$-balanced partition
\[
I=I_1\sqcup\cdots\sqcup I_t,
\qquad
2\le t\le O_k(1),
\qquad
\eta n\le |I_r|\le (1-\eta)n,
\]
such that
\[
\operatorname{Exc}_{\mathcal P}(K)\le \varepsilon n.
\]
\end{enumerate}
The fourth alternative is the recursive branch: it supplies a support-box decomposition with only $e^{\varepsilon n}$ overhead.
\end{conjecture}

\begin{remark}[Why support-product excess is the right quantity]\label{rem:low-excess-telescoping}
Fix an ordered balanced partition $\mathcal P=(I_1,\dots,I_t)$ and write
\[
P_r:=I_1\cup\cdots\cup I_r.
\]
Define
\[
d_r:=\log |\pi_{I_r}K|+\log |\pi_{P_{r-1}}K|-\log |\pi_{P_r}K|.
\]
Then $d_r\ge 0$ and
\[
\sum_{r=1}^t d_r
=
\sum_{r=1}^t \log |\pi_{I_r}K|-\log |K|
=
\operatorname{Exc}_{\mathcal P}(K).
\]
So a large support-product excess forces a linear local deficit on at least one block.
Programmatically, such a deficit should be converted into one of three already meaningful alternatives:
an exceptional profitable prefix, a bounded-memory solved-trace concentration, or a positive-density support box on a linear block.
If none of those outcomes occurs, the minimizing partition should have sublinear excess, which is exactly the recursive branch in Conjecture~\ref{conj:low-excess-extraction}.
\end{remark}

\begin{remark}[Sharpness and the residual leaf problem]\label{rem:low-excess-sharpness}
Balanced $q$-ary slices show that this conjecture cannot be replaced by a universal positive-density microscopic box theorem.
For a balanced slice one has
\[
\log |S_{m,q}|=m\log q-\frac{q-1}{2}\log m+O(1),
\]
so the natural balanced partition pays only $O(\log m)$ excess.
Products of mesoscopic slices therefore give examples with total excess $o(n)$ but no positive-density microscopic box on $\Omega(n)$ coordinates.

Thus even if Conjecture~\ref{conj:low-excess-extraction} is correct, one still needs a final theorem for the irreducible neutral leaves produced by repeated low-excess splitting.
At present those leaves may be thin slices, mesoscopic products, or some as-yet-unclassified obstruction class.
\end{remark}

\subsection*{Trace-atomic neutral leaves}
The recursive low-excess program has a sharply defined local endpoint.

\begin{definition}[Bad leaf]\label{def:bad-leaf}
Assume the repaired low-excess extraction theorem and its descendant re-entry property.
A leaf is called \emph{bad} if it is irreducible neutral and yields neither a solved witness nor an exposed-core witness under further admissible trace descent.
\end{definition}

\begin{lemma}[Minimal bad leaf lemma]\label{lem:minimal-bad-leaf}
Assume the repaired low-excess extraction theorem is stable under descendant re-entry and that witness heredity propagates from descendants to ancestors.
Let $W$ be a bad leaf of minimal complexity among all bad leaves.
Then for every nontrivial admissible trace $\psi$, if one starts from
\[
U:=\operatorname{Tr}_\psi(W)
\]
and re-runs the repaired extraction theorem on $U$, every resulting leaf is small.
\end{lemma}

\begin{proof}
Suppose some nontrivial admissible trace $\psi$ produces a descendant $U=\operatorname{Tr}_\psi(W)$ whose re-extracted tree contains a leaf $L$ that is not small.
There are then two possibilities.

If $L$ is won, then by witness heredity the ancestor $W$ yields a solved witness, contradicting that $W$ is bad.

If $L$ is irreducible neutral, then $L$ lies strictly below the proper trace $U$, so it has smaller complexity than $W$.
By minimality of $W$, the leaf $L$ is not bad.
Hence $L$ yields either a solved witness or an exposed-core witness.
Again witness heredity lifts that witness back to $W$, contradicting badness.

Therefore no non-small leaf can occur below a proper trace of $W$.
So every leaf appearing below every proper admissible trace of $W$ is small.
\end{proof}

\begin{definition}[Trace-atomic neutral leaf]\label{def:trace-atomic-leaf}
An irreducible neutral leaf $W$ is called \emph{trace-atomic} if for every nontrivial admissible trace $\psi$, the repaired extraction started from $\operatorname{Tr}_\psi(W)$ has only small leaves.
\end{definition}

\begin{proposition}[Recursive reduction to trace-atomic neutral leaves]\label{prop:trace-atomic-reduction}
Assume the repaired low-excess extraction theorem together with the two closure properties used in Lemma~\ref{lem:minimal-bad-leaf}.
Then every obstruction to the exposed-core / winning dichotomy can be represented by a trace-atomic neutral leaf.
\end{proposition}

\begin{proof}
If the dichotomy fails somewhere, there exists a bad leaf in the sense of Definition~\ref{def:bad-leaf}.
Choose one of minimal complexity.
By Lemma~\ref{lem:minimal-bad-leaf}, every proper admissible trace of that leaf re-extracts only to small leaves.
This is exactly the defining property of a trace-atomic neutral leaf.
\end{proof}

\begin{remark}[The first local gap after recursive low-excess extraction]\label{rem:trace-atomic-gap}
The recursive part of the low-excess route therefore closes.
What remains is purely local:
one must show that a trace-atomic neutral leaf is either won or else carries a nontrivial self-trace forcing an exposed neutral core.
So the honest residual class after the repaired extraction is not a generic irreducible neutral leaf, but the sharper class of trace-atomic neutral leaves.
\end{remark}



\subsection*{Lifted circuit defects, solved-cost explosive atoms, and primitive calibrated blockers}
The trace-atomic endpoint of the recursive low-excess program can be sharpened in three directions: the same-boundary realization problem collapses to one defect line, admissible trace covers are controlled by exact descendant costs, and the corresponding LP residue has a canonical blocker decomposition.

\begin{definition}[Strong fiber substitution hypothesis]\label{def:strong-fiber-substitution}
Let $T$ be a low-coherence support-irreducible trace-atomic neutral leaf with canonical exposed simplicial core
\[
E_T=\{e_1,\dots,e_m\}.
\]
We say that $T$ satisfies the \emph{strong fiber substitution hypothesis} if whenever an exposed support fiber over one vertex $e_i$ admits a profitable child, one may replace that fiber by the child while preserving the active quotient circuit data and strictly improving the lifted child-to-parent margin.
\end{definition}

\begin{proposition}[Singleton fibers under strong substitution]\label{prop:singleton-fibers}
Assume the strong fiber substitution hypothesis. Let
\[
\pi_T:\operatorname{Supp}(T)\longrightarrow E_T
\]
be the canonical factor map from the realized support to the abstract exposed core. Then every fiber
\[
F_i:=\pi_T^{-1}(e_i)
\]
is a singleton.
\end{proposition}

\begin{proof}
Because $E_T$ is exposed and simplicial, each vertex $e_i$ is exposed by some boundary functional $\varphi_i$ on the quotient core space. Pulling back $\varphi_i$ gives an exposed support piece
\[
P_i:=\{x\in \operatorname{Supp}(T):\varphi_i(\pi_T(x))=\max \varphi_i\},
\]
and this is exactly the fiber $F_i$.

Assume $F_i$ is not a singleton. Then $F_i$ is a proper support piece of the neutral terminal $T$, so by Theorem~\ref{thm:neutral-terminal-child} it has a profitable child. By the strong fiber substitution hypothesis, replacing the fiber by that profitable child preserves the active quotient circuit data and strictly improves the lifted witness margin. This yields either a profitable parent prefix, contradicting neutrality, or a strictly smaller failed support configuration, contradicting support-irreducibility. Therefore $F_i$ must be a singleton.
\end{proof}

\begin{proposition}[Defect-line reformulation of same-boundary realization]\label{prop:defect-line}
Assume the strong fiber substitution hypothesis. Let
\[
\lambda=(\lambda_1,\dots,\lambda_m),
\qquad
\sum_{i=1}^m \lambda_i=0,
\qquad
\sum_{i=1}^m \lambda_i e_i=0,
\]
be the signed affine circuit relation of the canonical exposed simplicial core $E_T$, and write
\[
\pi_T^{-1}(e_i)=\{x_i\}
\qquad (1\le i\le m)
\]
as in Proposition~\ref{prop:singleton-fibers}. Define the \emph{lifted circuit defect}
\[
\delta_T:=\sum_{i=1}^m \lambda_i x_i.
\]
Then the canonical exposed core is realized on an actual support subset in the same boundary metric if and only if
\[
\delta_T=0.
\]
Moreover, if $\delta_T\neq 0$, then the set
\[
R_T:=\{x_1,\dots,x_m\}
\]
is affinely independent.
\end{proposition}

\begin{proof}
Because $E_T$ is a circuit, its only affine dependence is the relation $\lambda$. By Proposition~\ref{prop:singleton-fibers}, the actual support above $E_T$ is the set $R_T$. Realizing the canonical core on an actual support subset in the same boundary metric means exactly that the same affine relation closes upstairs:
\[
\sum_{i=1}^m \lambda_i x_i=0.
\]
This is precisely the condition $\delta_T=0$.

Now suppose $\delta_T\neq 0$ and
\[
\sum_{i=1}^m \mu_i x_i=0,
\qquad
\sum_{i=1}^m \mu_i=0.
\]
Pushing this relation forward by $\pi_T$ gives an affine dependence among the $e_i$. Since $E_T$ is a circuit, there exists $c\in \R$ such that $\mu=c\lambda$. Hence
\[
0=\sum_{i=1}^m \mu_i x_i
=
c\sum_{i=1}^m \lambda_i x_i
=
c\,\delta_T.
\]
Because $\delta_T\neq 0$, one has $c=0$, so $\mu=0$. Therefore the $x_i$ are affinely independent.
\end{proof}

\begin{corollary}[Simplex-over-circuit failure shape]\label{cor:simplex-over-circuit}
Assume the hypotheses of Proposition~\ref{prop:defect-line}. If $\delta_T\neq 0$, then the actual support lift $R_T$ is an affine simplex, and the canonical exposed core $E_T$ is obtained from that simplex by quotienting out the defect line $\R\delta_T$.
\end{corollary}

\begin{proof}
The affine independence of $R_T$ is exactly the statement that it forms a simplex. By construction, the only affine dependence appearing after quotienting by the line $\R\delta_T$ is the original circuit relation $\lambda$, so the quotient is exactly the canonical exposed core.
\end{proof}

\begin{theorem}[Minimal defect terminal]\label{thm:minimal-defect-terminal}
Assume the strong fiber substitution hypothesis together with the repaired low-excess extraction, the finite base-window reduction package, and the reduction to trace-atomic neutral leaves. If some low-coherence neutral terminal has nonzero lifted circuit defect, then there exists a lexicographically minimal such terminal $T_*$ with the following properties:
\begin{enumerate}[label=(\roman*),leftmargin=2.2em]
\item $T_*$ is a trace-atomic support-irreducible low-coherence neutral leaf;
\item $T_*$ has a canonical rational exposed simplicial core
\[
E_*=\{e_1,\dots,e_m\};
\]
\item every fiber over $E_*$ is a singleton
\[
\pi_*^{-1}(e_i)=\{x_i\};
\]
\item its lifted circuit defect
\[
\delta_*:=\sum_{i=1}^m \lambda_i x_i
\]
is nonzero;
\item the actual lift
\[
R_*:=\{x_1,\dots,x_m\}
\]
is a simplex whose quotient by the line $\R\delta_*$ is $E_*$;
\item every proper support piece of $T_*$ has zero lifted circuit defect and therefore admits a profitable child;
\item the remaining obstruction carried by $\delta_*$ survives only through the residual rank-$1$ OA lane or the gauge-aware $(5,5)$ Hall lane.
\end{enumerate}
\end{theorem}

\begin{proof}
Choose a low-coherence neutral terminal with nonzero lifted circuit defect that is lexicographically minimal in trace depth, support size, ambient dimension, and common denominator of the circuit coefficients. The repaired low-excess extraction and the trace-atomic reduction allow us to take this minimal object to be trace-atomic and support-irreducible.

Items~(ii)--(v) follow from Proposition~\ref{prop:defect-line} and Corollary~\ref{cor:simplex-over-circuit}. For item~(vi), if some proper support piece had nonzero lifted circuit defect, then that piece would provide a strictly smaller counterexample to the same-boundary realization problem, contradicting minimality. Since every proper support piece of a neutral terminal has a profitable child, it follows that all proper pieces are locally harmless.

Once the quotient circuit relation has been accounted for, the only unresolved obstruction is the hidden defect line $\delta_*$. By the finite base-window reduction package, every remaining unresolved object enters one of the residual local lanes. The continuation-cone and solved-trace lanes have already been factored out, so the only surviving possibilities are the residual rank-$1$ OA lane or the gauge-aware $(5,5)$ Hall lane.
\end{proof}

\begin{definition}[Solved-cost admissible trace cover]\label{def:solved-cost-trace-cover}
Let $W$ be a trace-atomic neutral leaf of residual rank $r$ with no nontrivial self-trace. For every nontrivial admissible trace $\psi$, define its \emph{exact descendant cost} by
\[
c_W(\psi):=\operatorname{Sol}\bigl(\operatorname{Tr}_\psi(W)\bigr),
\]
where $\operatorname{Sol}$ is the exact solved-trace value furnished by the recursive package.
For a finite admissible trace family $\Psi$, write $a_\psi(w)\in\{0,1\}$ for the indicator that $\psi$ is admissible on $w\in W$, and define the \emph{trace-cover multiplicity}
\[
\tau_\Psi(W):=\min_{w\in W}\sum_{\psi\in\Psi} a_\psi(w).
\]
\end{definition}

\begin{theorem}[Solved-cost admissible trace-cover inequality]\label{thm:solved-cost-trace-cover}
Let $W$ be a minimal bad trace-atomic neutral leaf of residual rank $r$ with no nontrivial self-trace, and let $B_r$ denote the solved-trace threshold at rank $r$. Then every finite admissible trace family $\Psi$ satisfies
\[
\tau_\Psi(W)\,|W|
\le
\sum_{\psi\in\Psi} |\operatorname{Tr}_\psi(W)|
\le
\sum_{\psi\in\Psi} c_W(\psi)
\le
\sum_{\psi\in\Psi} B_{r-|\psi|}.
\]
In particular,
\[
\inf_{\Psi}
\frac{1}{\tau_\Psi(W)}
\sum_{\psi\in\Psi} c_W(\psi)
\ge B_r.
\]
\end{theorem}

\begin{proof}
For each $w\in W$, the number of traces in $\Psi$ admissible on $w$ is at least $\tau_\Psi(W)$. Summing over $w$ and reversing the order of summation gives
\[
\tau_\Psi(W)\,|W|
\le
\sum_{w\in W}\sum_{\psi\in\Psi} a_\psi(w)
=
\sum_{\psi\in\Psi}\sum_{w\in W} a_\psi(w)
=
\sum_{\psi\in\Psi} |\operatorname{Tr}_\psi(W)|.
\]

Because $W$ is trace-atomic, every proper admissible trace reruns only to solved descendants, so the exact solved-trace recursion applies to $\operatorname{Tr}_\psi(W)$. Hence
\[
|\operatorname{Tr}_\psi(W)|\le c_W(\psi)\le B_{r-|\psi|}.
\]
Summing over $\psi\in \Psi$ proves the full chain of inequalities.

Finally, if some admissible trace family $\Psi$ satisfied
\[
\frac{1}{\tau_\Psi(W)}
\sum_{\psi\in\Psi} c_W(\psi)<B_r,
\]
the displayed inequality would imply $|W|<B_r$, contradicting the minimal-bad calibration at rank $r$. Therefore the solved-cost ratio is always at least $B_r$.
\end{proof}

\begin{corollary}[Threshold-critical solved-cost explosions have no slack]\label{cor:solved-cost-threshold}
With the notation of Theorem~\ref{thm:solved-cost-trace-cover}, define
\[
\mathfrak C(W):=
\inf_{\Psi}
\frac{1}{\tau_\Psi(W)}
\sum_{\psi\in\Psi} c_W(\psi).
\]
If $\mathfrak C(W)=B_r$ and the infimum is attained by some admissible trace family $\Psi^\ast$, then every active trace in $\Psi^\ast$ satisfies
\[
c_W(\psi)=|\operatorname{Tr}_\psi(W)|.
\]
In particular, threshold-critical solved-cost explosions are built entirely from exact-threshold solved descendants.
\end{corollary}

\begin{proof}
At the threshold,
\[
B_r\le |W|
\le
\frac{1}{\tau_{\Psi^\ast}(W)}
\sum_{\psi\in\Psi^\ast} |\operatorname{Tr}_\psi(W)|
\le
\frac{1}{\tau_{\Psi^\ast}(W)}
\sum_{\psi\in\Psi^\ast} c_W(\psi)
=
B_r.
\]
Hence equality holds throughout. The final inequality is a sum of nonnegative terms
\[
c_W(\psi)-|\operatorname{Tr}_\psi(W)|\ge 0,
\]
so equality of the sums forces equality term by term on every active trace.
\end{proof}

\begin{definition}[Solved-cost cover LP and calibrated blocker]\label{def:calibrated-blocker}
Let $W$ be as in Definition~\ref{def:solved-cost-trace-cover}. Define its solved-cost cover value by
\[
\kappa(W):=
\inf\left\{
\sum_{\psi} c_W(\psi)x_\psi:
x_\psi\ge 0,\ 
\sum_{\psi} x_\psi a_\psi(w)\ge 1\ \text{for every }w\in W
\right\}.
\]
We call $W$ \emph{calibrated at threshold} if
\[
\kappa(W)=B_r.
\]
\end{definition}

\begin{proposition}[Two-sided scaling and weighted Hall inequalities]\label{prop:calibrated-blocker-scaling}
Assume $W$ is calibrated at threshold. Let $(x_\psi)$ be a primal optimum for $\kappa(W)$ and let $(\pi(w))_{w\in W}$ be a dual optimum normalized so that $\sum_{w\in W}\pi(w)=1$. Let
\[
\Psi:=\{\psi:x_\psi>0\},
\qquad
U:=\supp \pi,
\qquad
A_{w\psi}:=a_\psi(w),
\qquad
c_\psi:=c_W(\psi).
\]
Define
\[
M_{w\psi}:=\frac{B_r A_{w\psi}}{c_\psi},
\qquad
y_\psi:=\frac{x_\psi c_\psi}{B_r}.
\]
Then $y$ is a probability measure on $\Psi$ and
\[
My=\mathbf 1_U,
\qquad
M^\top \pi=\mathbf 1_\Psi.
\]
If
\[
N(S):=\{\psi\in\Psi:\exists\, w\in S \text{ with }A_{w\psi}>0\},
\]
then every subset $S\subseteq U$ satisfies the weighted Hall inequality
\[
\pi(S)\le y(N(S)).
\]
Equality holds if and only if every active trace touching $S$ is supported entirely inside $S$.
\end{proposition}

\begin{proof}
The primal LP defining $\kappa(W)$ has dual
\[
\sup\left\{
\sum_{w\in W}\alpha_w:
\alpha_w\ge 0,\ 
\sum_{w\in W}\alpha_w A_{w\psi}\le c_\psi \ \text{for every }\psi
\right\}.
\]
Because $W$ is calibrated at threshold, strong duality gives
\[
\sum_{\psi} c_\psi x_\psi = B_r = \sum_{w\in W}\alpha_w
\]
for some optimal dual vector $\alpha$. Normalize by writing
\[
\alpha_w=B_r\,\pi(w)
\]
with $\sum_w \pi(w)=1$.

Complementary slackness now yields
\[
\sum_{\psi\in\Psi} x_\psi A_{w\psi}=1
\qquad (w\in U)
\]
and
\[
c_\psi=B_r\sum_{w\in U}\pi(w)A_{w\psi}
\qquad (\psi\in \Psi).
\]
Multiplying the first equation by $1/B_r$ and substituting the definition of $y_\psi$ gives
\[
\sum_{\psi\in \Psi} M_{w\psi} y_\psi =1
\qquad (w\in U),
\]
that is,
\[
My=\mathbf 1_U.
\]
Similarly,
\[
\sum_{w\in U} M_{w\psi}\pi(w)
=
\frac{B_r}{c_\psi}\sum_{w\in U}\pi(w)A_{w\psi}
=1,
\]
so
\[
M^\top \pi=\mathbf 1_\Psi.
\]
Finally,
\[
\sum_{\psi\in\Psi} y_\psi
=
\frac{1}{B_r}\sum_{\psi\in\Psi} x_\psi c_\psi
=
1,
\]
so $y$ is a probability measure.

For the Hall inequality, define
\[
P_{w\psi}:=\pi(w)x_\psi A_{w\psi}.
\]
Then
\[
\sum_{\psi} P_{w\psi}=\pi(w)
\qquad\text{and}\qquad
\sum_w P_{w\psi}=y_\psi.
\]
Hence
\[
\pi(S)=\sum_{w\in S,\ \psi} P_{w\psi}
\le
\sum_{\psi\in N(S)} y_\psi
=
y(N(S)).
\]
Equality holds precisely when no mass of $P$ leaves $S$ through a trace touching $S$, which is equivalent to the statement that every active trace touching $S$ is supported entirely inside $S$.
\end{proof}

\begin{corollary}[Tight blocks uncross]\label{cor:tight-blocks-uncross}
In the setting of Proposition~\ref{prop:calibrated-blocker-scaling}, define a subset $S\subseteq U$ to be \emph{tight} if
\[
\pi(S)=y(N(S)).
\]
Then the tight subsets are closed under unions and intersections. Consequently the inclusion-minimal nonempty tight subsets are pairwise disjoint and partition $U$.
\end{corollary}

\begin{proof}
Let $S,T\subseteq U$ be tight. Add the equalities
\[
\pi(S)=y(N(S)),
\qquad
\pi(T)=y(N(T)).
\]
Since $\pi$ is a measure,
\[
\pi(S)+\pi(T)=\pi(S\cup T)+\pi(S\cap T).
\]
Since $N(S\cup T)=N(S)\cup N(T)$ and $N(S\cap T)\subseteq N(S)\cap N(T)$,
\[
y(N(S))+y(N(T))
\ge
y(N(S\cup T))+y(N(S\cap T)).
\]
Combining this with the Hall inequalities for $S\cup T$ and $S\cap T$ forces equality everywhere. Therefore both $S\cup T$ and $S\cap T$ are tight. The final statement is the standard lattice consequence for minimal nonempty tight sets.
\end{proof}

\begin{definition}[Primitive calibrated blocker]\label{def:primitive-calibrated-blocker}
A tuple
\[
\mathfrak B=(U,\Psi,M,\pi,y)
\]
is called a \emph{primitive calibrated blocker} if
\begin{enumerate}[label=(\roman*),leftmargin=2.2em]
\item $\pi>0$ on $U$ and $y>0$ on $\Psi$;
\item
\[
My=\mathbf 1_U,
\qquad
M^\top \pi=\mathbf 1_\Psi;
\]
\item every proper nonempty subset $S\subsetneq U$ satisfies the strict Hall inequality
\[
\pi(S)<y(N(S)).
\]
\end{enumerate}
\end{definition}

\begin{theorem}[Reduction to primitive calibrated blockers]\label{thm:primitive-calibrated-blocker}
Assume the repaired low-excess extraction package. Then every minimal bad family with a solved-cost explosive no-selftrace leaf reduces, up to the $e^{o(n)}$ overhead of the extraction, to a primitive calibrated blocker.
\end{theorem}

\begin{proof}
Start from a minimal bad family and descend to a trace-atomic no-selftrace leaf $W$ that is calibrated at threshold. Proposition~\ref{prop:calibrated-blocker-scaling} produces a two-sided scaled nonnegative kernel
\[
(U,\Psi,M,\pi,y)
\]
together with the tight-block lattice described in Corollary~\ref{cor:tight-blocks-uncross}.

Choose an inclusion-minimal nonempty tight block $U_0\subseteq U$. By the equality clause in Proposition~\ref{prop:calibrated-blocker-scaling}, every active trace touching $U_0$ is supported entirely inside $U_0$, so the LP residue splits along that block. Minimality of the original bad family allows us to replace $W$ by the induced obstruction on $U_0$. After renormalizing the corresponding restricted measures, the resulting tuple still satisfies
\[
My=\mathbf 1,
\qquad
M^\top \pi=\mathbf 1,
\]
and now has no proper nonempty tight subset. This is exactly the definition of a primitive calibrated blocker.
\end{proof}

\begin{remark}[What remains in the admissible-trace LP route]\label{rem:primitive-blocker-gap}
The LP and transport parts of the admissible-trace route are now fully explicit. What is still missing is the geometric theorem turning a transport-closed tight block into an actual admissible local core with the same descendant costs. Equivalently, one must show that every primitive calibrated blocker is either already an exposed neutral core or else enters one of the separator-locked primitive core models discussed earlier in this section.
\end{remark}

\status{The trace-atomic endpoint is now split into rigorous local reductions plus one clearly isolated structural gap. Under strong fiber substitution, the same-boundary realization problem is equivalent to vanishing of the lifted circuit defect $\delta_T$. Independently, the admissible-trace LP residue is a primitive calibrated blocker, and solved-cost explosions are the exact local obstruction below the recursive trace-cover theorem.}


\subsection*{Fixed-memory controlled mass inside multiplicative decompositions}
The fixed-memory route has a clean rigorous core, but only after one isolates the multiplicative part from the still-open decomposition theorem.

\begin{theorem}[Fixed-memory controlled mass gives a gap inside a multiplicative decomposition]\label{thm:controlled-mass-gap}
Fix $k\ge 3$, write $q:=k-1$, and fix $M\in \N$ and $\eps>0$. Suppose $\mathcal C_{k,M}$ is a class of factors with a certified base $a_{k,M}<q$, meaning that every factor $L\in \mathcal C_{k,M}$ of residual rank $r$ satisfies
\[
|L|\le a_{k,M}^r.
\]
Suppose a family $\PP$ admits an exact multiplicative decomposition
\[
|\PP|=\prod_{i=1}^t |L_i|,
\]
where each factor $L_i$ has residual rank $n_i$, the ranks sum to $n$, and every uncontrolled factor satisfies $|L_i|\le q^{n_i}$. If the total controlled rank obeys
\[
\sum_{i:\,L_i\in \mathcal C_{k,M}} n_i\ge \eps n,
\]
then
\[
|\PP|\le b_{k,M,\eps}^n,
\qquad
b_{k,M,\eps}:=a_{k,M}^{\eps}q^{1-\eps}<q.
\]
\end{theorem}

\begin{proof}
Let $N_C:=\sum_{i:\,L_i\in \mathcal C_{k,M}} n_i$ and $N_U:=n-N_C$. By hypothesis,
\[
|\PP|=\prod_{i=1}^t |L_i|\le a_{k,M}^{N_C}q^{N_U}.
\]
Since $N_C\ge \eps n$, this gives
\[
|\PP|\le a_{k,M}^{\eps n}q^{(1-\eps)n}=\bigl(a_{k,M}^{\eps}q^{1-\eps}\bigr)^n.
\]
Because $a_{k,M}<q$, the base $b_{k,M,\eps}:=a_{k,M}^{\eps}q^{1-\eps}$ is strictly smaller than $q$.
\end{proof}

\begin{corollary}[Near-neutral multiplicative examples avoid every fixed controlled class]\label{cor:controlled-mass-contrapositive}
Fix $k\ge 3$ and write $q:=k-1$. Let $\mathcal C_{k,M}$ be any class of factors with certified base $a_{k,M}<q$. If a sequence of multiplicatively decomposed examples $\PP_n$ has
\[
|\PP_n|=(q-o(1))^n,
\]
then for every fixed $M$ the proportion of total rank carried by factors from $\mathcal C_{k,M}$ tends to $0$.
\end{corollary}

\begin{proof}
If some $\eps>0$ occurred infinitely often, then Theorem~\ref{thm:controlled-mass-gap} would force those examples to have exponential base at most $a_{k,M}^{\eps}q^{1-\eps}<q$, contradicting the assumption $|\PP_n|=(q-o(1))^n$.
\end{proof}

\begin{remark}
Theorem~\ref{thm:controlled-mass-gap} is the rigorous core of the fixed-memory controlled-mass route. What is still missing is not merely a proof of positive-density extraction: Counterexample~\ref{ce:balanced-multislice} and Proposition~\ref{prop:mesoscopic-slice-density} show that such a statement is false in general. Any repair must either work after a weaker $e^{-o(n)}$-density extraction or insert a separate branch for thin neutral slice / mesoscopic block obstructions before attempting a multiplicative decomposition. The current $(4,4)$ and $(5,5)$ diagonalizers should \emph{not} be counted as certified controlled classes here, because their final global exponential constants remain open.
\end{remark}

\subsection*{Dense-box specialization and the residual obstruction class}
The abstract multiplicative theorem becomes more informative once one names the fixed-memory classes for which a strict exponential gap has already been proved.

\begin{definition}[Fixed-memory proved-gap classes]\label{def:proved-gap-classes}
Fix $k\ge 3$ and write $q:=k-1$. For $M\in \N$, let $\mathcal G_{k,M}^{\mathrm{gap}}$ denote the family of local block classes of memory at most $M$ for which there is a proved global bound
\[
F_\tau(m)\le \rho_\tau^m
\qquad\text{for all }m\ge 1
\]
with some constant $\rho_\tau<q$. Set
\[
\rho_{k,M}:=\max_{\tau\in \mathcal G_{k,M}^{\mathrm{gap}}}\rho_\tau<q.
\]
\end{definition}

\begin{theorem}[Dense-box fixed-memory controlled mass theorem]\label{thm:dense-box-controlled-mass}
Fix $k\ge 3$, write $q:=k-1$, and fix $M\in \N$ and $\eps>0$. Let $K$ be a terminal kernel of rank $n$. Suppose that, after passing to a positive-density support box $K^\square\subseteq K$ with
\[
|K^\square|\ge \delta |K|,
\]
the exact block-product theorem gives blocks
\[
B_1,\dots,B_t
\]
of coordinate lengths
\[
n_1,\dots,n_t,
\qquad
\sum_{j=1}^t n_j=n,
\]
and local classes $\tau_1,\dots,\tau_t$ such that
\[
|K^\square|\le \prod_{j=1}^t |B_j|,
\qquad
|B_j|\le F_{\tau_j}(n_j).
\]
If the total controlled coordinate mass
\[
m_{\mathrm{ctrl}}:=\sum_{j:\,\tau_j\in \mathcal G_{k,M}^{\mathrm{gap}}} n_j
\]
satisfies $m_{\mathrm{ctrl}}\ge \eps n$, then
\[
|K|\le \delta^{-1} c_{k,M,\eps}^n,
\qquad
c_{k,M,\eps}:=q^{1-\eps}\rho_{k,M}^{\eps}<q.
\]
\end{theorem}

\begin{proof}
By the hypotheses on the blocks,
\[
|K^\square|
\le
\prod_{j:\,\tau_j\in \mathcal G_{k,M}^{\mathrm{gap}}} F_{\tau_j}(n_j)
\prod_{j:\,\tau_j\notin \mathcal G_{k,M}^{\mathrm{gap}}} F_{\tau_j}(n_j).
\]
For a controlled block one has
\[
F_{\tau_j}(n_j)\le \rho_{k,M}^{n_j},
\]
while for an uncontrolled block we use only the ambient neutral bound
\[
F_{\tau_j}(n_j)\le q^{n_j}.
\]
Therefore
\[
|K^\square|
\le
\rho_{k,M}^{m_{\mathrm{ctrl}}} q^{\,n-m_{\mathrm{ctrl}}}
\le
\rho_{k,M}^{\eps n} q^{(1-\eps)n}
=
\bigl(q^{1-\eps}\rho_{k,M}^{\eps}\bigr)^n.
\]
Since $|K^\square|\ge \delta |K|$, it follows that
\[
|K|\le \delta^{-1}\bigl(q^{1-\eps}\rho_{k,M}^{\eps}\bigr)^n.
\]
Because $\rho_{k,M}<q$, the base $c_{k,M,\eps}:=q^{1-\eps}\rho_{k,M}^{\eps}$ is strictly smaller than $q$.
\end{proof}

\begin{corollary}[Near-neutral dense boxes avoid every fixed proved-gap class]\label{cor:dense-box-controlled-mass}
Fix $k\ge 3$, write $q:=k-1$, and let $(K_n)$ be a sequence of terminal kernels of ranks $n$ such that
\[
\limsup_{n\to\infty}|K_n|^{1/n}=q.
\]
For any fixed $M$, in every positive-density support-box / exact block-product decomposition of the form above, the proportion of coordinates carried by blocks from $\mathcal G_{k,M}^{\mathrm{gap}}$ tends to $0$.
\end{corollary}

\begin{proof}
Fix $M$. If some $\eps>0$ occurred infinitely often, then Theorem~\ref{thm:dense-box-controlled-mass} would give
\[
|K_n|\le \delta_n^{-1} c_{k,M,\eps}^{\,n}
\]
with $c_{k,M,\eps}<q$. Any positive-density support-box theorem has a density parameter bounded below independently of $n$, so the prefactor does not affect the exponential base. This contradicts the assumption that $|K_n|^{1/n}$ has limsup $q$. Hence the controlled proportion must tend to $0$.
\end{proof}

\begin{remark}[What currently remains in the dense-box regime]\label{rem:dense-box-residual-class}
At the present state of the project, the exact $(4,4)$ and $(5,5)$ diagonalizers do \emph{not} lie in $\mathcal G_{k,M}^{\mathrm{gap}}$, because their exact local reductions are rigorous but their final strict exponential bases are still open. Thus the residual dense-box obstruction class consists of near-neutral decompositions whose macroscopic coordinate mass escapes every fixed-memory class with a proved strict gap. This is the cleanest currently rigorous formulation of the ``minimal viable obstruction class'' for the dense-box route.
\end{remark}

\status{The fixed-memory controlled-mass theorem is therefore rigorous only after one conditions on an actual positive-density support box and counts \emph{only} those block classes for which a strict base below $k-1$ is already proved. Under that repair, near-neutral examples must avoid every fixed proved-gap class.}


\section{Audit of the initial Lean formalization program}\label{sec:lean-audit}
At the beginning of the project, a Lean formalization archive was audited. The audit yielded several useful structural conclusions.

\begin{enumerate}[label=(\alph*),leftmargin=2.2em]
\item The support-class and trace-class decompositions are real, and the basic ``trace classes are cone pieces'' bookkeeping is valid.
\item The existing seed-profile formalization was too weak to drive the exponential bound because the seed-averaging lemma allowed the empty seed $S=\varnothing$, making the main extraction vacuous.
\item Some later gap-analysis objects were placeholders: notably, a balanced-sparse predicate and some cover/LP shells existed formally but did not yet contain the missing combinatorial theorem.
\item The formalization was therefore best read as a map of the battlefield, not as a near-complete proof of an exponential sunflower bound.
\end{enumerate}

This early audit was valuable because it correctly identified the first serious logical gap in the then-current support-piece route: the missing extraction of a \emph{genuine nonempty} quantitatively useful seed.

\section{Concluding synthesis}\label{sec:conclusion}
The project now has a stable core.

\begin{itemize}[leftmargin=2em]
\item The exact reduction to terminal kernels is completely understood.
\item The exact heavy-set / prefix-product formalism is completely understood.
\item Broad support-piece casework, low-order LP inverse theorems, low-correlation rigidity, and bounded-local-coordinate extraction all fail in natural forms.
\item Exact block products, positive-density near-products, and fixed-memory finite-state classes are genuine theorem-producing model classes.
\item The right corrected bounded-rank trace statement is a bounded \emph{decomposition} theorem for traces, not the naive $Q_k=1$ sunflower-trace conjecture.
\item For the exact $(4,4)$ diagonalizer, the currently honest global rank output is only the ambient base $4$, and the local support is already balanced enough to kill the previous instability count.
\item The fixed-alphabet local-tensor lane remains the most promising route to new hard-model upper bounds, but its current diagonal theorems still require a repaired global rank argument.
\end{itemize}

In particular, the unrestricted sunflower problem now appears to have the following structure:
\begin{enumerate}[label=(\roman*),leftmargin=2.2em]
\item exact preprocessing reduces everything to terminal kernels,
\item graph-like bounded-degree parts are polynomially harmless,
\item the surviving obstruction universe splits into solved local tensor pieces, profitable-prefix pieces, and neutral obstruction pieces,
\item a universal routing theorem from arbitrary large terminal kernels into those classes would finish the exponential conjecture.
\end{enumerate}

The single sharp missing theorem, in the authors' view, is therefore a universal routing/profitable-prefix theorem for terminal kernels, or an equivalent bounded-rank sunflower-trace piece lemma.

\appendix

\section{Claim ledger}\label{app:ledger}
The following table records the main unique claims touched in the project. ``Rigorous'' means proved in this document. ``Conditional'' means exact once the stated hypothesis is granted. ``Refuted'' means disproved by an explicit counterexample. ``Exploratory'' means locally rigorous but globally incomplete.

\renewcommand{\arraystretch}{1.15}
\begin{longtable}{p{0.08\textwidth}p{0.47\textwidth}p{0.16\textwidth}p{0.18\textwidth}}
\toprule
ID & Claim & Status & Where in this paper\\
\midrule
\endfirsthead
\toprule
ID & Claim & Status & Where in this paper\\
\midrule
\endhead
R1 & Exact reduction $M_{\ker}(n,k)\le M(n,k)\le (k-1)^n M_{\ker}(n,k)$ & Rigorous & Theorem~\ref{thm:exact-reduction}\\
R2 & Exact one-round formula $m_K+s_K$ for sunflower size after stripping & Rigorous & Theorem~\ref{thm:one-round-sunflower}\\
R3 & Prefix-tree example sharp for $(k-1)^n$ stripping loss & Rigorous & Section~\ref{sec:leaf-stripping}\\
R4 & Exact trace recursion $|G|=\sum_{\tau}|G_\tau|$ & Rigorous & Theorem~\ref{thm:trace-recursion}\\
R5 & If $\Tr_P(G)$ is sunflower-free, then $|G|\le f(r,k)f(m-r,k)$ & Rigorous & Theorem~\ref{thm:trace-recursion}\\
R6 & One-block empty-core support pieces can encode arbitrary smaller coreless families & Rigorous & Proposition~\ref{prop:one-block-universality}\\
R7 & Exact heavy-set lemma $M_t(F)\ge |F|\prod_{i<t}\beta_i(F)$ & Rigorous & Theorem~\ref{thm:heavy-set}\\
R8 & Exact recursive upper bound via $B_{n,k}(t)$ & Rigorous & Corollary~\ref{cor:exact-recursion}\\
R9 & Adaptive kernel-prefix criterion implies exponential bound & Rigorous conditional theorem & Theorem~\ref{thm:adaptive-kernel-prefix}\\
R10 & Weighted choice-block invariant theorem & Rigorous & Theorem~\ref{thm:weighted-choice-block}\\
R11 & Raw uniformity / 2-core size is not the right correlation invariant & Rigorous counterexample & Section~\ref{sec:correlation}\\
R12 & Constant-scale correlation bound false even on the $2$-core & Rigorous counterexample & Counterexample~\ref{ce:simplex}\\
R13 & Sharp universal replacement $\sum d^2\ge m^2/(k-1)+(r-1)m$ & Rigorous & Section~\ref{sec:correlation}\\
R14 & Bounded-degree kernels are polynomially harmless & Rigorous & Theorem~\ref{thm:bounded-degree}\\
R15 & Degree-two kernels classified by simplices & Rigorous & Theorem~\ref{thm:degree-two}\\
R16 & Cheap bounded-block covers force heavy blocks & Rigorous & Section~\ref{sec:lp}\\
R17 & Exact dual classification of normalized block-box LP & Rigorous & Proposition~\ref{prop:oa-dual}\\
R18 & Product/correlation dichotomy on projected branches is false & Rigorous counterexample & Counterexample~\ref{ce:parity-projective}\\
R19 & Exact block-product theorem and stability & Rigorous & Theorem~\ref{thm:block-products}\\
R20 & Positive-density support-box theorem & Rigorous & Theorem~\ref{thm:positive-density-box}\\
R21 & Exact one-extra-symbol recurrence in $F_{q,k}(m)$ & Rigorous & Theorem~\ref{thm:extra-symbol}\\
R22 & Sharper randomized deletion lift & Rigorous & Theorem~\ref{thm:randomized-deletion}\\
R23 & Finite-state profitable-prefix theorem at rate $N_*(t)^{-1}$ / $\rho(T)^{-t}$ & Rigorous & Theorem~\ref{thm:finite-state-prefix}\\
R24 & Fixed-state-count entropy gap & Rigorous & Theorem~\ref{thm:finite-state-statecount}\\
R25 & Fixed-memory entropy gap with sharp extremizer & Rigorous & Theorem~\ref{thm:finite-state-fixed-order}\\
R26 & No $k$-only entropy gap across all memories & Rigorous counterexample & Theorem~\ref{thm:finite-state-fixed-order}\\
R27 & Broad support-piece / balanced-sparse route is fake & Rigorous counterexample & Proposition~\ref{prop:one-block-universality}\\
R28 & Low-$\sigma$ rigidity toward products is false & Rigorous counterexample & Section~\ref{sec:counterexamples}\\
R29 & Bounded-local-coordinate extraction is false & Rigorous counterexample & Section~\ref{sec:counterexamples}\\
R30 & A/B bridge false; trichotomy needed & Rigorous counterexample + synthesis & Section~\ref{sec:counterexamples}\\
R31 & $F_{4,4}(3)\ge 28$, disproving the naive power-of-two bound and forcing asymptotic base $\ge 28^{1/3}>3$ & Rigorous counterexample & Counterexample~\ref{ce:444-power-of-two-false} and Corollary~\ref{cor:444-asymp-lower}\\
R32 & Exact two-moment / quadratic-moment-curve reduction for $F_{5,5}(m)$ & Rigorous & Theorem~\ref{thm:5column-two-moment} and Corollary~\ref{cor:555-moment-curve}\\
R33 & Exact local scalar diagonalizer for the $(4,4)$ model & Rigorous & Proposition~\ref{prop:444-local-diagonalizer}\\
R34 & Honest global rank repair for the $(4,4)$ diagonalizer gives only $F_{4,4}(m)\le 4^m$ & Rigorous & Proposition~\ref{prop:444-honest-rank}\\
R35 & The $(4,4)$ local support is balanced; no centered one-body defect exists & Rigorous & Proposition~\ref{prop:444-balanced}\\
R36 & Parity slice with all proper marginals uniform & Rigorous counterexample & Counterexample~\ref{ce:parity-proper-marginals}\\
R37 & Forced residual trichotomy for solved / profitable / neutral terminals & Rigorous & Theorem~\ref{thm:forced-trichotomy}\\
R38 & Neutral terminals force a profitable child on every proper support piece & Rigorous & Theorem~\ref{thm:neutral-terminal-child}\\
R39 & Corrected bounded sunflower-trace decomposition lemma & Rigorous conditional theorem & Theorem~\ref{thm:corrected-trace-piece}\\
R40 & Bounded-rank solved-trace decomposition lemma & Rigorous conditional theorem & Theorem~\ref{thm:bounded-rank-solved-trace}\\
R41 & Abstract Bellman-face propagation for projected-branch LPs & Rigorous & Proposition~\ref{prop:bellman-face}\\
R42 & Fixed-memory controlled mass gives a multiplicative gap & Rigorous conditional theorem & Theorem~\ref{thm:controlled-mass-gap}\\
R43 & Coherent-lift reformulation of the missing terminal step; coherence threshold forces a heavy prefix, and genuine neutral terminals lie in a low-coherence / large-fractional-matching regime & Rigorous reformulation + structural reduction & Proposition~\ref{prop:coherent-lift}, Proposition~\ref{prop:coherence-heavy-vertex}, Corollary~\ref{cor:terminal-coherence-threshold}, and Corollary~\ref{cor:neutral-low-coherence}\\
R44 & Bounded-rank solved-trace pressure implies an exponential sunflower bound & Rigorous conditional theorem & Theorem~\ref{thm:solved-trace-pressure}\\
R45 & Solved-trace pressure equals the dual solved-trace packing number; uniform solved-trace capture implies bounded pressure & Rigorous & Proposition~\ref{prop:solved-trace-pressure-dual} and Corollary~\ref{cor:solved-trace-capture}\\
R46 & Weighted-profile OA kernel directions are the exact LP-invisible OA factors & Rigorous & Proposition~\ref{prop:weighted-oa-kernel} and Proposition~\ref{prop:weighted-oa-fg}\\
R47 & Regular-interface OA-reduced classes collapse to semilinear quotient-profile sets & Rigorous conditional theorem & Proposition~\ref{prop:weighted-oa-semilinear}\\
R48 & Exact predictive-state core: hereditary exact predictive states imply finite-state branching, and one state implies exact product & Rigorous & Theorem~\ref{thm:predictive-state-core}\\
R49 & Dense-box fixed-memory controlled mass forces a strict exponential gap below $k-1$ once such a box is given & Rigorous conditional theorem & Theorem~\ref{thm:dense-box-controlled-mass} and Corollary~\ref{cor:dense-box-controlled-mass}\\
R50 & Exact local correction, raw flattenings, and one-block monomial counts in the $(5,5)$ model all stall at base $5$ & Rigorous & Propositions~\ref{prop:555-unit-tensor}, \ref{prop:555-raw-tensor-obstruction}, and \ref{prop:555-local-monomial-obstruction}\\
R51 & Sharp two-coordinate theorem $F_{5,5}(2)=16$ and the resulting $F_{6,5}(2)<25$ bound & Rigorous & Theorem~\ref{thm:555-two-coordinate} and Corollary~\ref{cor:555-two-coordinate}\\
R52 & Reduction of the ordinary $(5,5)$ problem to the disjoint five-partite cross parameter $\Gamma_{5}^{\times}(m)$ & Rigorous & Proposition~\ref{prop:gamma555-reduction}\\
R53 & The disjoint five-partite cross parameter has asymptotic base $5$ & Rigorous obstruction theorem & Proposition~\ref{prop:gamma555-base5}\\
R54 & Perfect-matching data alone cannot tensorize the two-coordinate $(5,5)$ theorem & Rigorous obstruction theorem & Proposition~\ref{prop:555-perfect-matching-lineality}\\
R55 & Balanced multislices block positive-density microscopic support-box extraction; the sharp universal density scale is only $e^{-o(n)}$ & Rigorous counterexample & Counterexample~\ref{ce:balanced-multislice}, Corollary~\ref{cor:no-positive-density-extraction}, and Proposition~\ref{prop:mesoscopic-slice-density}\\

R56 & Exact convex lift theorem for weighted profitable children, together with its discrepancy form & Rigorous & Theorem~\ref{thm:exact-convex-lift} and Corollary~\ref{cor:exact-convex-lift-discrepancy}\\
R57 & Every unresolved terminal is already a support-irreducible low-coherence neutral terminal, and visible split forcing would close the terminal step & Rigorous structural reduction + conditional theorem & Proposition~\ref{prop:low-coherence-residual-class} and Theorem~\ref{thm:visible-split-reduction}\\
R58 & Genuine cone states satisfy an exact child-projection formula and strict rank-drop under transversal conditioning & Rigorous & Proposition~\ref{prop:cone-child-projection} and Corollary~\ref{cor:cone-rank-drop}\\
R59 & Raw exact-support conditioning can fail strict rank-drop & Rigorous counterexample & Counterexample~\ref{ce:raw-cone-rank}\\
R60 & Heavy exact links from bounded-transversal subcones are rigorous, and positive-mass bounded-transversal capture would imply an exponential bound & Rigorous local lemma + conditional theorem & Proposition~\ref{prop:heavy-exact-link} and Theorem~\ref{thm:bounded-transversal-capture}\\
R61 & The greedy heaviest support-pattern selector need not be hereditarily finite-state; the finite-profile criterion identifies the right replacement & Rigorous & Proposition~\ref{prop:selector-finite-profile} and Counterexample~\ref{ce:greedy-support-branch}\\
R62 & The two-block Hall quotient in the $(5,5)$ model has dimension $16$, but bare Hall profiles do not compose without gauge data & Rigorous obstruction theorem & Proposition~\ref{prop:555-hall-quotient} and Proposition~\ref{prop:555-hall-noncomposition}\\
R63 & Finite-state OA rigidity is equivalent to finite index of the continuation congruence, and semilinear quotient profiles do not by themselves imply finite-state collapse & Rigorous & Proposition~\ref{prop:oa-continuation-criterion} and Counterexample~\ref{ce:oa-one-counter}\\

C1 & Universal routing/profitable-prefix theorem for terminal kernels & Conditional / open & Section~\ref{sec:conditional}\\
C2 & Bounded sunflower-trace piece lemma & Conditional / open & Section~\ref{sec:conditional}\\
C3 & OA-factor routing trichotomy & Conditional / open & Conjecture~\ref{conj:oa-quotient}\\
C4 & Predictive-state neutral stability & Conditional / open & Conjecture~\ref{conj:predictive-state}\\
C5 & $\beta$-coherent lift hypothesis for neutral terminals & Conditional / open & Conjecture~\ref{conj:beta-coherent-lift}\\
C6 & One-step witness theorem for support-irreducible low-coherence neutral terminals & Conditional / open & Conjecture~\ref{conj:one-step-witness}\\
C7 & Visible split forcing for lifted child witnesses & Conditional / open & Conjecture~\ref{conj:visible-split-forcing}\\
C8 & Bounded-memory quotient-state theorem for OA-reduced neutral branches & Conditional / open & Conjecture~\ref{conj:quotient-state}\\
C9 & Recursive low-excess support-product extraction theorem & Conditional / open & Conjecture~\ref{conj:low-excess-extraction}\\
C10 & Positive-mass bounded-transversal capture theorem for continuation cones & Conditional / open & Theorem~\ref{thm:bounded-transversal-capture}\\
E1 & Global rank completion for the $(4,4)$ local diagonalizer, including the proposed base $16/(3\sqrt 3)$ & Exploratory; current slice count not certified & Remark~\ref{rem:444-bound-not-certified} and Section~\ref{subsec:exploratory-tensors}\\
E2 & Global rank/counting completion of the exact $(5,5)$ moment-curve reduction & Exploratory; exact reduction rigorous, final bound open & Section~\ref{subsec:exploratory-tensors}\\
E3 & $(6,6)$ local tensor with exact support and diagonalization & Exploratory; final rank step incomplete & Section~\ref{subsec:exploratory-tensors}\\
E4 & Matrix-valued power-of-two diagonalizer & Exploratory; global rank step missing & Section~\ref{subsec:exploratory-tensors}\\
E5 & Boolean-polynomial substitute for $(8,8)$ & Exploratory; lower-tail count not yet repaired & Section~\ref{subsec:exploratory-tensors}\\
\bottomrule
\end{longtable}

\section{Bibliographic background}\label{app:biblio}
The project repeatedly touches the following classical and modern landmarks.
\begin{thebibliography}{99}
\bibitem{ErdosRado1960}
P.~Erdos and R.~Rado,
\newblock Intersection theorems for systems of sets,
\newblock \emph{J. London Math. Soc.} \textbf{35} (1960), 85--90.

\bibitem{ALWZ2020}
R.~Alweiss, S.~Lovett, K.~Wu, and J.~Zhang,
\newblock Improved bounds for the sunflower lemma,
\newblock in \emph{Proceedings of STOC 2020}, 624--630; preprint arXiv:1908.08483.

\bibitem{Rao2020}
A.~Rao,
\newblock Coding for sunflowers,
\newblock \emph{Discrete Analysis} (2020), Article 2020:11; preprint arXiv:1909.04774.

\bibitem{BCW2021}
T.~Bell, S.~Chueluecha, and L.~Warnke,
\newblock Note on sunflowers,
\newblock \emph{Discrete Mathematics} \textbf{344} (2021), 112367; preprint arXiv:2009.09327.

\bibitem{Bloom2026}
T.~Bloom,
\newblock Erdos Problem \,\#20: uniform sets containing sunflowers,
\newblock \url{https://www.erdosproblems.com/20}, accessed April 20, 2026.

\bibitem{Fukuyama2025}
J.~Fukuyama,
\newblock Sunflower bound with a sub-logarithmic base,
\newblock preprint, arXiv:2510.19037, 2025.

\bibitem{Rao2025}
A.~Rao,
\newblock The story of sunflowers,
\newblock preprint, arXiv:2509.14790, 2025.

\bibitem{Tao2016}
T.~Tao,
\newblock Notes on the slice rank of tensors,
\newblock blog post, 2016, \url{https://terrytao.wordpress.com/2016/08/24/notes-on-the-slice-rank-of-tensors/}.

\bibitem{EllenbergGijswijt2017}
J.~S. Ellenberg and D.~Gijswijt,
\newblock On large subsets of $\mathbb F_q^n$ with no three-term arithmetic progression,
\newblock \emph{Annals of Mathematics} \textbf{185} (2017), 339--343.
\end{thebibliography}

\end{document}
